\documentclass[11pt]{article}
\usepackage[a4paper,margin=25mm]{geometry}
\usepackage[T1]{fontenc}
\usepackage{lmodern,microtype}
\usepackage{amsmath,amssymb,amsthm,mathtools,bm}
\usepackage{booktabs,array,longtable,enumitem}
\usepackage{xcolor}
\usepackage[colorlinks=true,linkcolor=blue!45!black,citecolor=blue!45!black,urlcolor=blue!45!black,filecolor=blue!45!black]{hyperref}
\usepackage{cleveref}
\hypersetup{
  pdftitle={Companion notes: degree correction and additional likelihood results},
  pdfauthor={spectral-lrt research project},
  pdfsubject={DC-SBM, historical decoders, and supplementary working material},
  pdfkeywords={spectral clustering, likelihood ratio, oracle risk, SBM, DC-SBM},
  bookmarksdepth=2
}
\setlist{itemsep=3pt,topsep=5pt,leftmargin=2em}
\setlength{\emergencystretch}{2em}
\allowdisplaybreaks[2]
\newtheorem{theorem}{Theorem}[section]
\newtheorem{lemma}[theorem]{Lemma}
\newtheorem{proposition}[theorem]{Proposition}
\newtheorem{corollary}[theorem]{Corollary}
\theoremstyle{definition}
\newtheorem{definition}[theorem]{Definition}
\newtheorem{assumption}[theorem]{Assumption}
\theoremstyle{remark}
\newtheorem{remark}[theorem]{Remark}
\newcommand{\R}{\mathbb R}
\newcommand{\E}{\mathbb E}
\newcommand{\PP}{\mathbb P}
\newcommand{\one}{\mathbf 1}
\newcommand{\diag}{\operatorname{diag}}
\newcommand{\rank}{\operatorname{rank}}
\newcommand{\tr}{\operatorname{tr}}
\newcommand{\op}{\mathrm{op}}
\newcommand{\Ber}{\operatorname{Bern}}
\newcommand{\KL}{\operatorname{KL}}
\newcommand{\TV}{\operatorname{TV}}
\newcommand{\col}{\operatorname{col}}
\newcommand{\supp}{\operatorname{supp}}
\newcommand{\argmax}{\operatorname*{arg\,max}}
\newcommand{\argmin}{\operatorname*{arg\,min}}
\newcommand{\norm}[1]{\left\lVert#1\right\rVert}
\newcommand{\abs}[1]{\left\lvert#1\right\rvert}
\newcommand{\ip}[2]{\left\langle#1,#2\right\rangle}
\newcommand{\ind}[1]{\mathbf1\{#1\}}
\title{Companion notes: degree correction\\and additional likelihood results\\
\large A working collection accompanying the SBM article}
\author{spectral-lrt research project}
\date{20 October 2026}
\begin{document}
\maketitle
\begin{abstract}
These notes collect material separated from the ordinary-SBM article.
The first part gives the bounded-parameter degree-corrected model,
the distinction between profile and known-target-degree oracles,
and the regularized fitted Bernoulli profile classifier. Its current
risk bound requires only $np\to\infty$ and $np^2\le C_1$, with no
logarithmic lower bound. The DC-specific nuisance estimation,
profile reduction, and tilted cancellation are retained here.
Shared degree-floor, spectral-row-tail, and local deleted-graph tools
are recalled from the SBM article without repeating their full proofs.
The later parts preserve the separate earlier unregularized decoders,
the wider-probability ordinary free-profile result, an obstruction for
the unrestricted known-target-degree oracle, and auxiliary arguments.
This is a working companion, not a second article with a unified scope.
Historical hypotheses are attached to their original decoders and do
not restrict the current regularized theorem.
\end{abstract}
\paragraph{Reading map.}
The first article is
\href{../sbm/sbm_spectral_likelihood.pdf}{\emph{Spectral likelihood classification in stochastic block models}}.
It is independent of this companion. The source package contains the
following overflow from the combined revision.
\begin{center}
\small
\begin{tabular}{>{\raggedright\arraybackslash}p{.46\textwidth}>{\raggedright\arraybackslash}p{.45\textwidth}}
\toprule
Original component & Present location\\
\midrule
Ordinary regularized theorem and shared full proofs
& The first article; shared statements recalled in Section~3 here.\\[3pt]
DC model, current regularized algorithm and DC fitting
& Sections~1--5; bounded degree parameters and no logarithmic lower bound.\\[3pt]
Earlier normalized geometry, ordinary free-profile decoder and probability closure
& Sections~6--8 and~10; their separate historical algorithms and full proofs.\\[3pt]
Earlier DC decoder and classical oracle bridge
& Section~9, with its probability closure in Section~10.\\[3pt]
Known-target obstruction and zero-degree fallback
& Section~11 and Section~12.1.\\[3pt]
Optional exponential-moment transfer and retained-spectrum identity
& Sections~3.9 and~12.2.\\[3pt]
Implementation scope, sparse history and audit records
& Sections~12.3--12.4; source directories \texttt{audits/} and \texttt{archive/}.\\
\bottomrule
\end{tabular}
\end{center}
\clearpage
\setcounter{tocdepth}{1}
\tableofcontents
\clearpage
\section{Degree-corrected models and oracle benchmarks}\label{sec:statements}
The target is a constant oracle-risk bound
\begin{equation}\label{eq:target}
 \E\ell(\widehat z,z)\le C R_*(n),
\end{equation}
with $C$ independent of $n$. The companion retains ordinary-SBM
notation where needed by the historical free-profile analysis.
The independent ordinary-SBM article contains the current ordinary
adjacency block-LRT theorem and its complete proof.

\subsection{Bernoulli models and fixed structural constants}

There are $n$ vertices and $K\ge2$ groups, with fixed $K$.
Let $z\in[K]^n$, let $Z$ be its indicator matrix, and let $n_a$ be the size
of group $a$. Throughout $n_a\ge\pi_-n$ for a fixed $\pi_->0$.
The edges above the diagonal are independent, the adjacency matrix $A$ is
symmetric, and $A_{ii}=0$. In DC-SBM,
\begin{equation}\label{eq:mean-model}
 A_{ij}\sim\Ber(\theta_i\theta_j p C_{z_i z_j}),\quad i<j,
 \qquad d=np,\qquad
 \Omega=\Theta Z(pC)Z^\top\Theta.
\end{equation}
Here $\Theta=\diag(\theta_i)$, $C=C^\top$, and the fixed structural bounds are
\begin{equation}\label{eq:structural}
 0<\theta_-\le\theta_i\le\theta_+<\infty,\quad
 \sum_{i:z_i=a}\theta_i=n_a,\quad
 0<c_-\le C_{ab}\le c_+<\infty,\quad
 \sigma_{\min}(C)\ge\kappa>0.
\end{equation}
The model class is nonempty only when $1\in[\theta_-,\theta_+]$.
All constants in asymptotic bounds may depend on these structural bounds,
but not on $n$. One may take a fixed $C$, or a sequence in this same compact
class. Ordinary SBM is the special case $\theta_i=1$, with block probability
matrix $P=pC$.

The completed matrix $\Omega$ includes a diagonal for algebraic convenience;
the actual mean is $\E A=\Omega-\diag(\Omega_{ii})$.
That deterministic difference has norm $O(p)$ and is included in all error
bounds. Put $D=\diag(A\one)$ and
$\lambda_i=(\E A\one)_i$. Thus $\lambda_i\asymp d$ and
$\Omega\one$ equals $d\theta_i S_{z_i}$ in coordinate $i$, where
\begin{equation}\label{eq:exposures}
 \pi_a=n_a/n,\qquad S_a=\sum_b\pi_b C_{ab}.
\end{equation}
The loss for the estimated partition is
\[
 \ell(\widehat z,z)=\frac1n\min_{\sigma\in\mathfrak S_K}
 \sum_i\ind{\sigma(\widehat z_i)\ne z_i}.
\]

\subsection{Three explicitly different local oracles}\label{sec:oracle}

All oracle classifiers below use the raw incident row $A_{i,-i}$ and the
true neighbor labels. They are local likelihood classifiers: an exact global
block-size constraint is not used to determine the one missing label.
For simplicity candidate weights are uniform. Any fixed weights bounded
above and below by positive constants give the same conclusions.

For ordinary SBM, the oracle chooses the maximizer of
\begin{equation}\label{eq:sbm-oracle}
 \ell_{i,*}(a)=\sum_{j\ne i}
 \{A_{ij}\log P_{a z_j}+(1-A_{ij})\log(1-P_{a z_j})\}.
\end{equation}
Its average error probability is $R_{\mathrm{SBM},*}$.
For DC-SBM, the \emph{known-target-degree oracle} replaces $P_{a z_j}$ in
\eqref{eq:sbm-oracle} by $\theta_i\theta_j p C_{a z_j}$ and is supplied
with every true $\theta_j$, including $\theta_i$. Its risk is
$R_{\theta,*}$.

The \emph{profile oracle} is supplied with true neighbor parameters and
$pC$, but profiles the target degree separately for every candidate:
\begin{equation}\label{eq:profile-oracle}
 \ell_{i,\mathrm{prof},*}(a)=
 \sup_{t\ge0}\sum_{j\ne i}
 \{A_{ij}\log(t\theta_jpC_{a z_j})
 +(1-A_{ij})\log(1-t\theta_jpC_{a z_j})\}.
\end{equation}
We restrict candidate edge probabilities to at most $q_{\max}=1/2$; the
constraint is inactive on the events used in the sparse theorem. Write
$R_{\mathrm{prof},*}$ for its average error probability. This is a
full-edge generalized likelihood-ratio reference. It is not described as a
Bayes test between two fixed known-$\theta_i$ laws. In all three definitions,
the risk averages the actual directional error at each fixed true vertex.
Ties are resolved by a fixed ordering; strict-tail arguments below cover
any such ordering.


\section{The current regularized DC-SBM result}\label{sec:sparse-restored}
\label{restore:algorithm}

Fix sufficiently small $a_0,\alpha>0$ and sufficiently large $\beta>0$,
depending only on the fixed structural parameter bounds. Choose them so
that the population degrees lie strictly inside the propensity-clamping
interval, and the population floor is inactive. Put
\[
 \mathbf d=A\one,\quad \bar D=n^{-1}\one^\top\mathbf d,\quad
 \tau=a_0\bar D,\quad G_i=\max(D_i,\tau).
\]
The all-zero graph has a fixed output. Otherwise $G$ is positive, even
when some vertices are isolated. Retain the $K$ eigenpairs with largest
absolute eigenvalues of $S_G=G^{-1/2}AG^{-1/2}$ and the raw degrees.
There is no requirement that a Perron degree vector be one of these
eigenvectors after flooring. Compute
\[
 R_0=G^{1/2}U\Lambda U^\top G^{1/2}=AH_G,
 \qquad H_G=G^{-1/2}UU^\top G^{1/2},
\]
and put, with $r=\diag(R_1)$ understood as a column vector,
\begin{equation}\label{restore:reconstruction}
 R_1=R_0+(\mathbf d-R_0\one)\one^\top/n,\qquad
 R^h=R_1+\frac{r\one^\top}{n-1}-\frac n{n-1}\diag(r).
\end{equation}
Thus $R^h\one=\mathbf d$ and $\diag(R^h)=0$ exactly. Only $(\mathbf d,U,\Lambda)$ are used
after encoding; $G$ is recovered from $\mathbf d$.

Initialize by normalizing the rows of $U$ and using the deterministic
radius-$1/8$, deletion-radius-$1/4$ count procedure in
\cref{dcsbm:initializer}. A zero row is assigned any fixed unit direction
for this initialization, without stopping other vertices.

Define
\[
 T_i=\min\{\max(D_i,\alpha\bar D),\beta\bar D\},\qquad
 \widehat\theta_i=T_i/\overline T_{\widehat z_i},\qquad
 \overline T_a=\widehat n_a^{-1}\sum_{j:\widehat z_j=a}T_j.
\]
Fit the symmetric block matrix by the exact signed Bernoulli working
criterion \eqref{dcsbm:block-criterion}, with $R^h$ and these degree
estimates, under the same fixed candidate-probability cap. Classify each
row by exact Bernoulli target profiling, with the degree parameter
profiled separately for every candidate community.

At least one fit-and-LRT update is performed. The procedure may use any
finite number of these hard refits, including finite data-dependent
stopping. The invariant basin below covers every such endpoint
simultaneously. An undefined empty-group or parameter-fit operation may
have a fixed fallback; those operations are well-defined on the proved
aggregate event. No individual zero degree triggers a graph-wide fallback.

\begin{theorem}[Regularized DC-SBM at every diverging expected-degree
scale]\label{restore:dc-theorem}\label{thm:regularized}
Under the fixed structural assumptions \eqref{eq:mean-model}--%
\eqref{eq:structural}, suppose
\[
 d=np\longrightarrow\infty,\qquad np^2\le C_1<\infty.
\]
The regularized algorithm of \cref{restore:algorithm}, using the
DC-SBM profile LRT, satisfies
\[
 \E\ell(\widehat z,z)\le C R_{\mathrm{prof},*}(n)
\]
for all sufficiently large $n$. The constant may depend on the fixed
structural bounds, clamping constants and $C_1$, but not on $n$.
Under the finite-$n$ equal-exposure condition
\eqref{eq:equal-exposure}, the right side may be replaced by
$C R_{\theta,*}(n)$. In particular, the classical balanced planted
partition with fixed separated $q/p$ is included. There is no lower
bound proportional to $\log n$ in this statement.
\end{theorem}


The reference $R_{\mathrm{prof},*}$ uses the full incident edge row.
The second comparison is available only under the stated equal-exposure
condition. The impossibility result in \cref{thm:imposs} shows why a
known-target-degree oracle is not a valid unrestricted benchmark.
Only retained degrees and one eigensystem enter the algorithm;
deleted graphs below are proof devices.

\section{Shared tools and their degree-corrected specialization}\label{sec:sparse-extension}
The ordinary-SBM article now uses a high-degree-capped adjacency
encoder. Its sharp Bernoulli transfer, same-threshold spectral-row tail,
and deterministic local scalar comparison also apply here. We recall
those tools with their degree-corrected premises. The small-floor
geometry and its exponentially accurate event are proved in full below,
since they differ from the large-threshold adjacency geometry in that
article. The DC classifier in these notes retains its normalized
encoder; no large-threshold DC extension is asserted.

\subsection{Sharp likelihood transfer}\label{sec:sharp}
\begin{lemma}[Bounded tilted error preserves the sharp prefactor]
\label{lem:sharp}
Let $X_j$ be independent Bernoulli variables with probabilities $p_j$,
and let
\[
 T=b_0+\sum_jw_jX_j,\qquad
 K(s)=\log\E e^{-sT}.
\]
Suppose $\abs{w_j}\le W$, the minimizer $s_*$ satisfies
$\tau\le s_*\le1-\tau$ for a fixed $\tau>0$, and under the probability
measure $Q$ proportional to $e^{-s_*T}$,
\begin{equation}\label{eq:sharp-premises}
 \E_Q T=0,\quad vd\le\operatorname{Var}_Q(T)\le V_0d,
 \quad\sum_jr_j\le V_0d,
\end{equation}
where $r_j=Q(X_j=1)$. Put $I=-K(s_*)$.
For every fixed $L<\infty$ and all sufficiently large $d$, uniformly over
$|x|\le L$,
\begin{equation}\label{eq:sharp-oracle}
 c_Ld^{-1/2}e^{-I}\le\PP(T<x)
 \le\PP(T\le x)\le C_Ld^{-1/2}e^{-I}.
\end{equation}

Let $\mathcal W$ be independent of $(X_j)_j$ and let
\[
 \widehat T=T+b+\sum_je_jX_j
\]
have coefficients measurable with respect to $\mathcal W$. On a measurable
training event suppose
\begin{equation}\label{eq:tilt-error}
 \max_j|e_j|\le M,\qquad
 \left|b+\sum_jr_je_j\right|\le B,\qquad
 \sum_jr_je_j^2\le V.
\end{equation}
For fixed $M,V$ and every real $x$,
\begin{equation}\label{eq:sharp-transfer}
 \PP(\widehat T\le x\mid\mathcal W)
 \le C\exp\{s_*x+B+e^M V/2\}\,d^{-1/2}e^{-I}.
\end{equation}
The constant is uniform over the training event. In particular $B$ can
be a nonnegative training-measurable random bound, provided its indicated
exponential moment is integrated when removing the conditioning.
\end{lemma}


\noindent\emph{Proof reference.}
Lemma~A.1 of the SBM article proves this by
Berry--Esseen under the exact Bernoulli tilt, followed by a second tilt
for the independent-training perturbation. That argument makes no
block-model assumption. It therefore applies directly to all scores
used below. See also Zhou and Li~\cite{zhouli} for sharp Chernoff
comparisons in community detection.

\subsection{Application to the actual Bernoulli SBM oracle}

For a true label $a$ and a competitor $b$, the oracle contrast is
\[
 T_{i,ab}=\ell_{i,*}(a)-\ell_{i,*}(b)
 =b_{i,ab}+\sum_{j\ne i}w^{ab}_{z_j}A_{ij},
\]
where
\begin{equation}\label{eq:sbm-weights}
 w_c^{ab}=\log\frac{P_{ac}(1-P_{bc})}{P_{bc}(1-P_{ac})},\qquad
 b_{i,ab}=\sum_{j\ne i}\log\frac{1-P_{az_j}}{1-P_{bz_j}}.
\end{equation}
Its exact Chernoff information is
\begin{equation}\label{eq:sbm-information}
 I_{i,ab}= -\min_{0\le s\le1}
 \sum_{j\ne i}\log\left[
 P_{az_j}^{1-s}P_{bz_j}^{s}
 +(1-P_{az_j})^{1-s}(1-P_{bz_j})^{s}\right].
\end{equation}
On the compact positive full-rank classes used here, probabilities are
comparable to $p$ and bounded away from one, the weights are bounded,
$s_*$ stays away from the endpoints, and the tilted variance is comparable
to $d$. To verify the last assertions, the derivatives of the transform
at its endpoints are respectively minus and plus the two directional
Bernoulli KL divergences, both of order $d$. Its second derivative is
bounded by $Cd$. The row separation implied by
$\sigma_{\min}(P)\ge\kappa p$ gives a lower bound $cd$ on that derivative
throughout $[0,1]$: tilted probabilities are comparable to $p$ and the
sum of squared log-odds differences is at least $cn$.
Thus \cref{lem:sharp} applies uniformly to every pair.

The same conclusion holds for the profile contrasts defined below, with
their own directional information $J_{i,a\to b}$.
It is not necessary for those two directional exponents to coincide.

\subsection{Fixed-\texorpdfstring{$K$}{K} comparison and label alignment}

Let $q_i$ be the error probability of the specified oracle at vertex $i$.
For any coherent family of its candidate scores, strict pairwise loss to a
competitor implies an oracle error; an oracle error implies non-strict
loss to some competitor. Hence
\[
 \max_{b\ne a}\PP(T_{i,ab}<0)
 \le q_i\le\sum_{b\ne a}\PP(T_{i,ab}\le0).
\]
By \eqref{eq:sharp-oracle},
\begin{equation}\label{eq:oracle-pairs}
 q_i\asymp\frac1{\sqrt d}\sum_{b\ne z_i}e^{-I_{i,ab}}
\end{equation}
for the SBM oracle; replace $I$ by $J$ for the profile oracle. The constants
may include a factor depending on fixed $K$.
Fitted partitions are aligned with the true groups by a single common
permutation on the good event. Uniform score comparisons under that
alignment imply
\[
 \E[\ell(\widehat z,z)\ind{\mathcal G}]
 \le\frac1n\sum_i\PP(\text{aligned fitted error at }i,\mathcal G).
\]
The ensuing pairwise bounds apply to fixed true coefficient vectors; a
data-dependent matching does not become an independence assumption.
Finally,
\begin{equation}\label{eq:bad-payment}
 \E\ell(\widehat z,z)
 \le C R_*+\PP(\mathcal G^c).
\end{equation}
\Cref{sec:probability} proves, for the historical decoders, $\PP(\mathcal G^c)=O(R_*)$ under the stated
model conditions. It is this comparison, rather than merely
$\PP(\mathcal G^c)=o(1)$, that completes each theorem.

\subsection{Empirical-floor geometry for bounded degree parameters}
Throughout this section $D_0=\diag(\Omega\one)$,
$G_i=\max(D_i,a_0\bar D)$, and $\norm v_G^2=v^\top Gv$.
The floor multiplier is fixed and sufficiently small.
\begin{proposition}[Sparse degree-floor geometry]\label{prop:sparse-floor}
Under the fixed structural conditions \eqref{eq:mean-model}--%
\eqref{eq:structural}, suppose $d=np\to\infty$ and $np^2\le C_1$.
For every fixed $r>0$, with probability at least $1-C_rn^{-r}$,
$cd\le G_i$, $cnd\le\sum_iD_i$, and
\begin{align}
 \norm{G^{-1/2}(A-\Omega)G^{-1/2}}_{\op}&\le C_r d^{-1/2},
 &\sum_iG_i&\le C_rnd,\label{eq:floor-noise}\\
 \sum_i\frac{(D_i-(D_0)_{ii})^2}{G_i}
 +\sum_i\frac{(G_i-(D_0)_{ii})^2}{G_i}
 &\le C_rn.\label{eq:floor-degree-energy}
\end{align}
Let $U$ contain the $K$ eigenvectors of $S_G$ with largest absolute
eigenvalues, let $\Pi=UU^\top$, and let
$H=G^{-1/2}\Pi G^{1/2}$. For all sufficiently large $n$, on the same
event the retained eigenvalues have a fixed separation from zero, the
remaining eigenvalues have magnitude at most $C_r/\sqrt d$, and
\begin{equation}\label{eq:floor-block-residual}
 \norm{(I-\Pi)G^{1/2}Z}_F\le C_r\sqrt n.
\end{equation}
For every block-constant $f$ with $\norm f_\infty\le C_f$,
\begin{equation}\label{eq:floor-weighted-errors}
 \norm{(I-H)f}_G\le C_rC_f\sqrt n,\qquad
 \norm{(I-H)f}_2^2\le C_rC_f^2n/d,
 \quad \norm v_G^2=v^\top Gv.
\end{equation}
After alignment, the orthonormal population basis
\[
 V=G^{-1/2}\Theta Z
       (Z^\top\Theta G^{-1}\Theta Z)^{-1/2}O
\]
satisfies $\norm{U-V}_F\le C_r/\sqrt d$ and
$\sqrt{G_i}\norm{V_{i\cdot}}_2\le C_r\sqrt{d/n}$.
The same assertions hold simultaneously for recomputed graphs after
any fixed number of deletions, after increasing the probability
parameter. No lower bound proportional to $\log n$ is imposed.
\end{proposition}


\begin{proof}
The Bernoulli transform of the total number of edges gives
$cnd\le\sum_iD_i\le Cnd$ except with probability $e^{-cnd}$.
For a sufficiently small fixed $a_0$, on this event
\[
 cd\le\tau=a_0\bar D\le Cd,
 \qquad \tau<\min_i(D_0)_{ii}.
\]
Let $d_* = n\max_{i\ne j}\E A_{ij}\asymp d$, and put
$W_{h,ii}=\sqrt{\min(1,2d_*/D_i)}$, with value one at zero degree.
The set $\mathcal S=\{i:D_i>2d_*\}$ has size at most $10n/d_*$
except with probability $e^{-cn}$. Indeed, if
$m=\lfloor10n/d_*\rfloor+1\le n$ vertices exceed the threshold,
some set $T$ of size $m$ satisfies $\sum_{i\in T}D_i>2md_*$.
This is an independent-edge sum with coefficients in $\{1,2\}$
and mean at most $md_*$. At $t=(\log2)/2$, its exponential
transform bounds this probability by
$\exp\{-(\log2-1/2)md_*\}$, since
$(e^{tc}-1)/c\le(e^{2t}-1)/2=1/2$ for $c=1,2$.
The union over $T$ costs at most
$\exp\{m\log(en/m)\}=\exp\{o(md_*)\}$.
If $m>n$, the assertion is automatic.

Every row of $W_hAW_h$ has squared Euclidean norm at most
$W_{h,ii}^2D_i\le2d_*$. The regularized-adjacency theorem of
Le, Levina and Vershynin
\cite[Theorem 2.1 and Remark 2.2]{llv2017}, in its row-norm form,
therefore gives
\[
 \norm{W_hAW_h-\E A}_{\op}\le C_r\sqrt d
\]
with failure probability at most $C_rn^{-r}$. It permits these
observed-graph-dependent weights. The difference
$W_h\Omega W_h-\Omega$ is supported on rows or columns in
$\mathcal S$, and every entry has magnitude at most $Cp$; hence
\[
 \norm{W_h\Omega W_h-\Omega}_{\op}
 \le Cp\sqrt{n|\mathcal S|}\le C\sqrt d.
\]
The completed-mean diagonal costs $O(p)$. Thus
$\norm{W_h(A-\Omega)W_h}_{\op}\le C_r\sqrt d$.
The exact factorization
\[
 G^{-1/2}=(2d_*)^{-1/2}FW_h,
 \qquad F_{ii}=\sqrt{\frac{\max(D_i,2d_*)}{G_i}}
\]
has uniformly bounded $F$, because $G_i\ge cd$ and $G_i\ge D_i$.
Conjugating proves the noise bound in \eqref{eq:floor-noise}.
Also $G_i\le D_i+\tau$ gives $\sum_iG_i\le Cnd$.

Applying the noise bound to $G^{1/2}\one$ gives
\[
 \sum_i\frac{(D_i-(D_0)_{ii})^2}{G_i}
 =\norm{G^{-1/2}(A-\Omega)\one}_2^2\le C_rn.
\]
Since $\tau<(D_0)_{ii}$, the maximum defining $G_i$ only reduces
the coordinatewise distance to $(D_0)_{ii}$. This proves the
second energy bound in \eqref{eq:floor-degree-energy}.

Let $J=Z^\top\Theta G^{-1}\Theta Z$. In each block,
Cauchy--Schwarz, the bounded degree parameters, and comparable
block masses give
\[
 \frac{cn}{d}
 \le\frac{(\sum_{i:z_i=a}\theta_i)^2}
          {\sum_{i:z_i=a}G_i}
 \le J_{aa}=\sum_{i:z_i=a}\frac{\theta_i^2}{G_i}
 \le\frac{Cn}{d}.
\]
The nonzero spectrum of $G^{-1/2}\Omega G^{-1/2}$ is that of
$pJ^{1/2}CJ^{1/2}$, which has a fixed separation from zero.
Subspace perturbation and \eqref{eq:floor-noise} give the claimed
eigenvalue and Frobenius bounds. The formula for $V$ and the lower
bound on $J$ give its weighted row bound.

The identity $(D_0)_{ii}=d\theta_iS_{z_i}$ shows that
$G^{-1/2}D_0Z$ belongs to
$\col(G^{-1/2}\Theta Z)$. Moreover
\[
 \norm{G^{1/2}Z-G^{-1/2}D_0Z}_F^2
 =\sum_i\frac{(G_i-(D_0)_{ii})^2}{G_i}\le C_rn,
 \qquad \norm{G^{1/2}Z}_F^2\le C_rnd.
\]
Combining this approximation with the $O(d^{-1/2})$ projector
error proves \eqref{eq:floor-block-residual}.
The matrix $H$ is the $G$-orthogonal projector corresponding to
$\Pi$. Writing $f=Zb$ then gives the first bound in
\eqref{eq:floor-weighted-errors}; the second follows from $G_i\ge cd$.

Deleting a fixed number of vertices changes each block mass and
each sum of degree parameters by $O(1)$. The preceding argument
uses their comparability, not exact mean-one normalization on a
minor, and applies with its recomputed degrees and threshold.
Increasing the fixed probability parameter before a finite-depth
union bound proves the simultaneous assertion.
\end{proof}

\subsection{Spectral-row tails and local deleted-graph comparison}
Let $\mathcal E_r$ denote the aggregate event just described, including
recomputed versions after at most three deletions. Each deleted graph
also has its own training-measurable event. The same-threshold row-tail
proof of Lemma~C.1 in the adjacency SBM article applies to any fixed
multiplier in $\tau/\bar D$, provided the aggregate premises hold.
Indeed, with $W=\tau^{1/2}G^{-1/2}$ and $\widetilde A=WAW$, the preceding
proposition gives $\norm{W(A-\Omega)W}_{\op}\le C\sqrt d$ and
an informative gap of order $d$. Also
\[
 W_i\le1,\quad W_i^2D_i\le\tau,\quad
 \norm{\widetilde A}_{\op}\le\tau,\quad
 \norm{\Omega_{i\cdot}}_2\le Cd/\sqrt n.
\]
These bounds and the exact threshold change
$\tau^{(-i)}-\tau=a_0(\bar D-2D_i)/(n-1)$ are all the
weight and population inputs to that proof. It uses neither
$\tau\ge\max_i(D_0)_{ii}$ nor equality of the degree parameters.
Thus the small multiplier $a_0$ used here gives
\begin{lemma}[Truncated energy for the degree-floor eigenvectors]
\label{restore:tail-energy-lemma}
There are fixed constants $R_0,C,c>0$ such that, for every fixed
$R\ge R_0$ and all sufficiently large $n$,
\begin{equation}
 \E\!\left[\ind{\mathcal E_r}
       \sum_{j=1}^n\norm{U_{j\cdot}}_2^2
                  \ind{\sqrt n\norm{U_{j\cdot}}_2>R}\right]
 \le CR^2e^{-cd\log R}.
 \label{restore:tail-energy}
\end{equation}
The constants $R_0,C,c$ depend on the fixed structural bounds,
$a_0$, and $r$, but not on $R$ or $n$.  The sufficiently-large-$n$
condition may depend on the chosen fixed $R$.  The same result holds
for each recomputed deleted graph with its training-measurable
geometry event.
\end{lemma}


The proof keeps the empirical-threshold perturbation of order $1/n$,
uses unrotated training coefficients in its conditional Bennett bound,
and absorbs the recurrence at the same truncation threshold. It therefore
also applies on the stronger aggregate event below. For an all-zero
training graph its proof-only convention is a deterministic orthonormal
fallback and identity weights; no conditioning on a full-graph event
is introduced.

\begin{corollary}[Exceptional coefficients are hit with exponentially small mass]
\label{restore:tail-hit-corollary}
On a training graph with the preceding floor, put
$H=G^{-1/2}UU^\top G^{1/2}$.
Let $f$ be any training-measurable vector satisfying
$\norm f_\infty\le F_0$, and let $e=(H-I)f$.
For every prescribed fixed $H_0>0$, there is a fixed $M<\infty$
such that, for any independent tested Bernoulli row with probabilities
$p_j\le Cd/n$,
\begin{equation}
 \E\!\left[\ind{\mathcal E_r}
                 \sum_{j:|e_j|>M}p_j\right]
 \le C_{H_0}e^{-H_0d}.
 \label{restore:tail-hit-mass}
\end{equation}
Here $\mathcal E_r$ is the training-measurable geometry event.
No exchangeability among degree parameters is required.
\end{corollary}

\begin{proof}
On that event, $G_j\ge a_-d$ and $\sum_jG_j\le C_Gnd$.
Consequently
\[
 |e_j|
 \le F_0+\frac{\norm{U_{j\cdot}}_2}{\sqrt{G_j}}
                  \norm{U^\top G^{1/2}f}_2
 \le F_0+F_0\sqrt{C_G/a_-}\sqrt n\norm{U_{j\cdot}}_2.
\]
For $M_R=F_0+F_0\sqrt{C_G/a_-}R$, the event
$|e_j|>M_R$ implies $\sqrt n\norm{U_{j\cdot}}_2>R$.
Write $H_R(U)=\sum_j\norm{U_{j\cdot}}_2^2
\ind{\sqrt n\norm{U_{j\cdot}}_2>R}$.
Then \cref{restore:tail-energy-lemma} gives
\begin{equation}
 \E\!\left[\ind{\mathcal E_r}
                    \sum_j\ind{|e_j|>M_R}\right]
 \le\frac n{R^2}\E[H_R(U)\ind{\mathcal E_r}]
 \le Cne^{-cd\log R}.
 \label{restore:tail-coefficient-count}
\end{equation}
Multiplication by $Cd/n$ bounds the expected hit mass by
$Cd e^{-cd\log R}$.
Choose a fixed $R$ with $c\log R>H_0+1$ and absorb the factor $d$.
Independence is required only to interpret the conditional probability
that the tested row hits the training exceptional set; no independence
between the training geometry event and its eigenvectors is asserted.
\end{proof}

\begin{corollary}[Average exceptional target-row mass]
\label{restore:tail-target-corollary}
For every fixed degree cutoff $L$ and every fixed $H_0>0$, there is
a fixed $L_0$ such that
\begin{equation}
 \frac1n\sum_i\PP\!\left\{
    \sqrt{G_i}\norm{U_{i\cdot}}_2>L_0\sqrt{d/n},\quad
    D_i\le Ld,\quad \mathcal E_r\right\}
 \le C_{H_0,L}e^{-H_0d}.
 \label{restore:tail-target-mass}
\end{equation}
\end{corollary}

\begin{proof}
On the displayed event, $G_i\le C_Ld$.
Choose $L_0=R\sqrt{C_L}$.
The target-row exception then implies
$\sqrt n\norm{U_{i\cdot}}_2>R$.
Apply \eqref{restore:tail-energy} and divide its corresponding
expected count by $n$.  Choosing fixed $R$ sufficiently large gives
\eqref{restore:tail-target-mass}.
\end{proof}

\subsection{The local scalar comparison}\label{restore:loo}
\begin{lemma}[A local floored scalar comparison]\label{restore:loo-scalar}
Write $A=\left(\begin{smallmatrix}0&x^\top\\x&B\end{smallmatrix}\right)$,
$k=\sum_jx_j\le Ld$, and let $\Delta$ be the diagonal degree floor
recomputed from $B$. Suppose
\[
 cd\le\Delta_j,\qquad \sum_j\Delta_j\le Cnd,
 \qquad S_B=\Delta^{-1/2}B\Delta^{-1/2}
   =U_B\Lambda_BU_B^\top+N,
\]
where $NU_B=0$, $\norm N_{\op}\le C/\sqrt d$, and the $K$ retained
eigenvalues have a fixed gap from zero. Let the full floor be
$G=\diag(g,\Delta+E)$, with
\[
 cd\le g\le Cd,\qquad E=E_\mathrm{loc}+E_\mathrm{glob},\qquad
 0\le E_\mathrm{loc}\le\diag(x),\quad
 \norm{E_\mathrm{glob}}_{\op}\le Cd/n.
\]
Here $d\to\infty$ and $d\le Cn$. Suppose only on hit coordinates that
\begin{equation}\label{restore:loo-hit}
 \norm{\Delta_j^{-1/2}(U_B)_{j\cdot}}_2
       \le R/\sqrt{nd}\quad\text{whenever }x_j=1.
\end{equation}
Let $U,\Lambda$ be the full informative eigensystem of
$G^{-1/2}AG^{-1/2}$, and put
\[
 H_G=G^{-1/2}UU^\top G^{1/2},\qquad
 H_B=\Delta^{-1/2}U_BU_B^\top\Delta^{1/2}.
\]
For every $\norm w_\infty\le W$, uniformly over the stated inputs,
\begin{equation}\label{restore:loo-main}
 (AH_Gw)_i=x^\top H_Bw_{-i}+O_{L,R,W}(1).
\end{equation}
No bound is assumed on unhit rows of $U_B$ or on $\max_j\Delta_j$.
\end{lemma}


\noindent\emph{Proof reference.}
Lemma~D.1 of the SBM article is deterministic
under precisely these premises. Its generalized eigenvector equations
contain the floored top degree $g$, not the raw incident degree $k$.
Consequently its proof applies unchanged here.

\begin{corollary}[Exact-degree and hollow comparison]\label{restore:loo-corrected}
Under the same local premises, the reconstruction
\eqref{restore:reconstruction} satisfies, for
$\bar w_{-i}=(n-1)^{-1}\sum_{j\ne i}w_j$,
\begin{equation}\label{restore:loo-corrected-action}
 (R^h w)_i=x^\top\{H_B(w_{-i}-\bar w_{-i}\one)
                        +\bar w_{-i}\one\}+O_{L,R,W}(1).
\end{equation}
\end{corollary}
The same deterministic proof, including its local diagonal bound,
is given in Corollary~D.2 of the SBM article. Uniformity in $w$ is algebraic; fitted
coefficients are first replaced by true block contrasts in
\cref{restore:fit-score} before conditioning on the deleted graph.

\subsection{Aggregate probabilities over all degree scales}
\begin{lemma}[Exponentially accurate aggregate geometry above a logarithmic scale]
\label{restore:high-aggregate}
Fix any $A_0>0$ and $H>0$. Under the fixed positive compact structural
conditions, if $d\ge A_0\log n$ and $np^2\le C_1$, there is an aggregate
event $\mathcal E_H$ with
\[
 \PP(\mathcal E_H^c)\le C_H e^{-Hd}
\]
on which the empirical-floor geometry, weighted degree-energy bounds,
two-sided total-degree bounds, and true block-total bounds used in
\cref{restore:fit-geometry-one,restore:fit-geometry-two,restore:fit-block-event}
hold for the full graph and every graph obtained by deleting at most three
vertices. The constants may depend on $A_0,H,C_1$ and the structural
bounds. No condition is imposed on the smallest raw degree or the
largest spectral row.
\end{lemma}


\begin{proof}
Bounded degree parameters give independent upper-triangular entries
with probabilities at most $Cd/n$. The centered bounded-entry norm
inequality of Bandeira--van Handel
\cite[Corollary 3.12 and Remark 3.13]{bvh} gives
\[
 \PP\{\norm{A-\Omega}_{\op}>L_H\sqrt d\}
 \le e^{-(H+2)d}
\]
after choosing a sufficiently large fixed $L_H$. Indeed, the row
variance sum is $O(d)$, the entry bound is one, and the prefactor $n$
in the norm-tail inequality is absorbed using $\log n\le d/A_0$.
The completed-mean diagonal has norm $O(p)$ and is absorbed in the
same threshold. No minimum raw degree is imposed.

The Bernoulli transform of the total number of edges gives, with
failure probability at most $Ce^{-cnd}$,
\[
 cnd\le\sum_iD_i\le Cnd.
\]
For each of the fixed $K^2$ true block totals, Bernstein's inequality
for the independent upper-triangular edges, whose coefficients have
absolute value at most two, gives
\[
 \PP\left\{\left|[Z^\top(A-\Omega)Z]_{ab}\right|>Ln\right\}
 \le 2\exp\left\{-\frac{cL^2n}{d+L}\right\}.
\]
The deterministic completed diagonal costs $O(d)$ and is included
by increasing $L$. Since $d\to\infty$ and $d^2/n\le C_1$,
the exponent is at least $c'L^2d/C_1$ for all sufficiently large $n$.
Increasing the fixed $L=L_H$ therefore makes the failure of all block
totals at most $C_He^{-(H+2)d}$. The total-degree failure is also
bounded by $Ce^{-(H+2)d}$ for large $n$.

These three full-graph events imply the required events on every
fixed-depth minor without a degree-wise probability union. To see this,
let $J=[n]\setminus T$ with $|T|\le3$, set $A_J=A_{JJ}$ and
$\Omega_J=\Omega_{JJ}$, and recompute all degrees and the empirical
floor on $J$. The centered matrix $A_J-\Omega_J$ is a principal
submatrix of $A-\Omega$, so its norm is at most $L_H\sqrt d$.
Removing $T$ changes the total raw degree by at most $2|T|n$;
as $d\to\infty$, the two-sided bound by fixed multiples of $nd$
therefore persists uniformly over $T$. Removing $T$ changes a true
block total of $A-\Omega$ by $O(|T|n)$ deterministically, since
the entries have bounded absolute value. Its $O_H(n)$ bound persists
as well. The surviving group masses and sums of degree parameters differ from
their original values by $O(1)$ and remain comparable to $n$.

The recomputed average degree is consequently comparable to $d$.
For the fixed sufficiently small floor multiplier $a_0$, its threshold
$\tau_J=a_0\overline D_J$ lies between positive fixed multiples of
$d$ and below every coordinate of $D_{0,J}=\Omega_J\one$.
Thus $G_{J,j}=\max(D_{J,j},\tau_J)\ge cd$ and
$\sum_jG_{J,j}\le Cnd$. The raw norm bound immediately gives
\begin{align*}
 \norm{G_J^{-1/2}(A_J-\Omega_J)G_J^{-1/2}}_{\op}
   &\le C_Hd^{-1/2},\\
 \sum_{j\in J}\frac{(D_{J,j}-D_{0,J,j})^2}{G_{J,j}}
   &=\norm{G_J^{-1/2}(A_J-\Omega_J)\one}_2^2
     \le C_Hn.
\end{align*}
Since the population floor is inactive,
$|G_{J,j}-D_{0,J,j}|\le|D_{J,j}-D_{0,J,j}|$; this proves the
other weighted energy bound. The fixed signal gap and the residual
$\norm{(I-U_JU_J^\top)G_J^{1/2}Z_J}_F\le C_H\sqrt n$
now follow by the deterministic Gram-matrix and projection arguments
in \cref{prop:sparse-floor}. Those arguments use only the floor,
total degree, normalized noise, and degree energy, all just proved.
Intersecting the stated aggregate consequences defines
$\mathcal E_H$. The three full-graph events imply it simultaneously
for every indicated minor, proving its claimed failure bound.

For a deleted graph used as independent training data, form the
analogous event using that graph and its further minors. The same
argument gives failure probability $C_He^{-Hd}$, and this event is
measurable from the training graph alone. The full-graph common event
is used only for deterministic implications, never as a condition
that would preserve independence of an incident row.
\end{proof}


\subsection{The degree-corrected tilted cancellation}

Work on a deleted graph so that all coordinates below index actual
neighbors of a separate tested vertex. For the bounded block-constant
linear candidate contrast used in the DC profile reduction,
the exact tilted Bernoulli probabilities have the decomposition
\begin{equation}\label{eq:floor-tilt-decomposition}
 r=D_0h+\epsilon,\qquad
 h\text{ block constant},\quad
 \norm h_\infty\le C/n,\quad \norm\epsilon_\infty\le Cp^2.
\end{equation}
For example the leading term of the exact tilt is
$\theta_i\theta_jpC_{a,z_j}e^{-s w_{z_j}}$, while
$(D_0)_{jj}=\theta_jp\sum_k\theta_k C_{z_j,z_k}$.
Dividing the former by the latter gives the required block-constant
$h_j$. The denominator in the exact Bernoulli tilt contributes
$O(p^2)$ uniformly. This uses the original Bernoulli law throughout.
The target coordinate is absent; artificially appending $r_i=0$ would
instead create an $O(p)$ error at that one coordinate.

For bounded block-constant $f$, put $e=(H-I)f$. The
$G$-orthogonality of $H$ gives the exact cancellation
\begin{align}
 r^\top e
 ={}&-\langle(I-H)h,(I-H)f\rangle_G
       +((D_0-G)h)^\top e+\epsilon^\top e.
       \label{eq:floor-tilted-cancellation}
\end{align}
The first term is $O(1)$ by \eqref{eq:floor-block-residual}, since
the block coefficients of $h$ are $O(1/n)$. For the second,
Cauchy--Schwarz and \eqref{eq:floor-degree-energy} give
\[
 |((D_0-G)h)^\top e|
 \le \norm h_\infty
 \left(\sum_j\frac{((D_0)_{jj}-G_j)^2}{G_j}\right)^{1/2}
 \norm e_G
 \le C.
\]
The last term is at most
$Cp^2\sqrt n\,\sqrt{n/d}\le Cnp^2$. Also $r_j\le Cp$, so
\begin{equation}\label{eq:floor-tilt-moments}
 |r^\top e|\le C(1+np^2),\qquad
 \sum_jr_je_j^2\le Cp\norm e_2^2\le C.
\end{equation}
Thus the first-order bias cancellation and bounded quadratic error do
not require logarithmic degree. The replacement for a uniform
coefficient bound is supplied by
\cref{restore:tail-energy-lemma,restore:tail-hit-corollary}.

\subsection{Two sharp transfers that do not assume every coefficient is bounded}

The next two lemmas give concrete alternatives to a global
row-incoherence event. They concern a row-independent, training-measurable
score, just as in \cref{lem:sharp}.

\begin{lemma}[Exponential-moment coefficient control]\label{lem:sparse-exp-transfer}
Assume the oracle hypotheses of \cref{lem:sharp}. Let
$\widehat T=T+b+\sum_je_jX_j$, with $(b,e)$ measurable from independent
training data. On a measurable training event suppose
\[
 \left|b+\sum_jr_je_j\right|\le B,\qquad
 \sum_jr_je_j^2e^{|e_j|}\le V,
\]
where $B,V$ are fixed. Then, for all real $x$ and all sufficiently
large $d$,
\[
 \PP(\widehat T\le x\mid\text{training})
 \le C e^{s_*x+B+V/2}d^{-1/2}e^{-I}.
\]
The constant is independent of the largest coefficient.
\end{lemma}

\begin{proof}
Taylor's inequality gives, under the original oracle tilt $Q$,
\[
 \log\E e^{-s_*\widehat T}
 \le-I-s_*\left(b+\sum_jr_je_j\right)
       +\frac12\sum_jr_je_j^2e^{|e_j|}
 \le-I+B+V/2.
\]
Tilt again by $e^{-s_*\widehat T}$. The new probabilities satisfy
\[
 \widehat r_j
 =\frac{r_je^{-s_*e_j}}{1-r_j+r_je^{-s_*e_j}}
 \le r_je^{s_*|e_j|}.
\]
For $|e_j|\le1$, this bounds the new centered third absolute moments
by $C_W\sum_jr_j=O(d)$. For $|e_j|>1$, use $s_*\le1-\tau$ and
\[
 e^{s_*|u|}(W+|u|)^3
 \le C_{W,\tau}u^2e^{|u|}\qquad (|u|>1).
\]
The remaining third moments are therefore $O(V)$. The upper variance
bound follows from the same estimates. There is no restriction that
every $e_j$ be bounded.

For the lower variance, restrict to $|e_j|\le1$. The original variance
discarded by this restriction is at most $W^2V$, and the weighted
$\ell_2$ norm of the remaining perturbation is at most $\sqrt V$.
Thus the square root of the original-weight variance of the remaining
$(w_j+e_j)$ is at least $\sqrt{vd-W^2V}-\sqrt V$.
On this restricted set, the new and old Bernoulli variances differ by
only fixed factors. The new total variance is therefore at least $cd$
for large $d$. Berry--Esseen gives a $C/\sqrt d$ bound for every unit
interval under the new tilt. Summing those intervals, using the first
displayed exponential-moment bound, proves the assertion as in
\cref{lem:sharp}.
\end{proof}

\begin{lemma}[Paying exceptional coefficients only when the row hits them]
\label{lem:sparse-truncated-transfer}
Assume the oracle hypotheses of \cref{lem:sharp}, and write $p_j$
for the un-tilted row probabilities. Suppose $p_j\le C_pr_j$, as
holds here for bounded oracle likelihood coefficients. On an independent
training event $\mathcal E$, suppose
\[
 \left|b+\sum_jr_je_j\right|\le B,\qquad \sum_jr_je_j^2\le V.
\]
Here $B,V$ are fixed constants. For any fixed $M>0$ define
\[
 h_M=\sum_{j:|e_j|>M}p_j,\qquad
 \eta_M=\sum_jr_je_j^2\ind{|e_j|>M}.
\]
Then, for every fixed bounded $x$,
\begin{equation}\label{eq:sparse-truncated-transfer}
 \PP(\widehat T\le x)
 \le C_Md^{-1/2}e^{-I}
       +\E[\ind{\mathcal E}h_M]
       +\PP(\mathcal E^c).
\end{equation}
Since $h_M\le C_p\eta_M/M^2$, an expected truncated quadratic error of order
$d^{-1/2}e^{-I}$ suffices. A bound $Ce^{-Hd}$ is a convenient
sufficient condition whenever $I\le C_Id$ and $H>C_I$.
The expected hit-mass condition on $h_M$ is weaker and is itself enough.
\end{lemma}

\begin{proof}
For analysis only, replace $e_j$ by
$e'_j=e_j\ind{|e_j|\le M}$. The truncated coefficients satisfy
\[
 \left|b+\sum_jr_je'_j\right|
 \le B+\eta_M/M\le B+V/M,\qquad
 \sum_jr_j(e'_j)^2\le V.
\]
Hence \cref{lem:sharp} bounds their score error by
$C_Md^{-1/2}e^{-I}$ on $\mathcal E$.
The original and truncated random scores are identical whenever
$X_j=0$ for every $j$ with $|e_j|>M$. Conditional on training, the
probability that this fails is at most
\[
 \sum_{j:|e_j|>M}p_j
 =h_M\le \frac{C_p}{M^2}\eta_M.
\]
Integrate this bound on $\mathcal E$ and pay at most
$\PP(\mathcal E^c)$ outside it. No coefficient truncation is an
additional input or computation required of the implemented decoder.
\end{proof}

For the degree-floor projection there is a concrete reduction to a
weighted row count. Define
\[
 \zeta_j=\sqrt{\frac{nd}{G_j}}\norm{U_{j\cdot}}_2.
\]
If $e=(H-I)f$, $\norm f_\infty\le C_f$ and
$\sum_jG_j\le C_Gnd$, then deterministically
\[
 |e_j|
 \le \frac{\norm{U_{j\cdot}}_2}{\sqrt{G_j}}
        \norm{U^\top G^{1/2}f}_2+C_f
 \le C_f\sqrt{C_G}\,\zeta_j+C_f.
\]
Choosing $M=C_f\sqrt{C_G}L+C_f$ thus gives
$\{|e_j|>M\}\subseteq\{\zeta_j>L\}$. Since $p_j\le Cd/n$,
the sufficient expected-count estimate
\begin{equation}\label{eq:sparse-row-count-target}
 \E\!\left[\ind{\mathcal E}
       \sum_j\ind{\zeta_j>L}\right]\le Cne^{-Hd},
 \qquad H>C_I,
\end{equation}
would imply
$\E[\ind{\mathcal E}h_M]\le Cd e^{-Hd}
=o(d^{-1/2}e^{-I})$ whenever $I\le C_Id$.
Equation \eqref{eq:sparse-row-count-target} is not a consequence of
aggregate geometry alone. It is proved by the same-threshold
truncated-energy argument in \cref{restore:tail-energy-lemma}
and the floor lower bound.


\section{DC fitting from aggregate geometry and local row bounds}
\label{restore:fit}

This section supplies the deterministic fitting module for the
degree-floor decoder. It uses the exact Bernoulli working criterion and
the exact-degree reconstruction $R^h$ in \eqref{restore:reconstruction}.
In particular, its rank-one degree correction and simple row hollowing
are different from the earlier correction $\diag(s)(H-I)$.
One profile refinement suffices. The optional finite repeated hard-refit
version is covered explicitly by \cref{restore:fit-invariant-basin}.

Put $d=np$, $B=pC$, $\Omega=\Theta ZBZ^\top\Theta$ and
$D_{0,i}=(\Omega\one)_i$. As before, $\mathbf d=A\one$ is the raw
degree vector, $D_i=\mathbf d_i$, and
\[
 G_i=\max(D_i,a_0\overline D),\qquad
 S=G^{-1/2}AG^{-1/2},\qquad
 R_0=G^{1/2}U\Lambda U^\top G^{1/2}.
\]
The constants $a_0>0$ and all subsequent clipping constants are fixed
using only the structural parameter bounds. They do not depend on $n$,
the unknown labels, or individual unobserved parameters.
Work on the two-sided total-degree event on which
$\overline D\asymp d$ and the population floor is inactive.
The aggregate geometric and block-total premises used below are
\begin{align}
 \norm{S-G^{-1/2}\Omega G^{-1/2}}_{\op}
   &\le C/\sqrt d,&
 \sum_iG_i&\le Cnd, \label{restore:fit-geometry-one}\\
 \sum_i\frac{(G_i-D_{0,i})^2}{G_i}
 +\sum_i\frac{(D_i-D_{0,i})^2}{G_i}
   &\le Cn,&
 \norm{(I-UU^\top)G^{1/2}Z}_F&\le C\sqrt n,
 \label{restore:fit-geometry-two}\\
 \max_{a,b}\left|[Z^\top(A-\Omega)Z]_{ab}\right|
   &\le Cn.&& \label{restore:fit-block-event}
\end{align}
Here $G_i\ge cd$, $D_{0,i}\asymp d$, and the retained signal has a fixed
gap from zero. The degree-floor proposition already supplies
\eqref{restore:fit-geometry-one}--\eqref{restore:fit-geometry-two}.
The fixed-$K$ Bernoulli block-total bound
\eqref{restore:fit-block-event} has failure probability
$C\exp\{-cL^2n/(d+L)\}$ when its fixed threshold constant is $L$.
Under $np^2=d^2/n\le C_1$, increasing $L$ pays this failure with any
prescribed fixed exponential rate in $d$. No degree-wise union bound
is needed for this block event.

Write $\Pi=UU^\top$. An aligned orthonormal basis of the random
population space has the exact form
\[
 V=G^{-1/2}\Theta Z
       (Z^\top\Theta G^{-1}\Theta Z)^{-1/2}O,
\]
where $O$ is orthogonal. The diagonal Gram matrix has entries comparable
to $n/d$. The signal gap and \eqref{restore:fit-geometry-one} imply
\begin{equation}\label{restore:fit-basis}
 \norm{U-V}_F\le C/\sqrt d,\qquad
 \sqrt{G_i}\norm{V_{i\cdot}}_2\le C\sqrt{d/n}.
\end{equation}
An orthogonal rotation used in this display does not change any row
norm or projection. Also $\norm S_{\op}\le1$: the inequality
$|x^\top Ax|\le\sum_iD_ix_i^2\le\sum_iG_ix_i^2$ proves this directly.
Hence $\norm\Lambda_{\op}\le1$.

\begin{lemma}[Uniform weighted mass of wrong-label sets]
\label{restore:fit-wrong-mass}
Fix $M<\infty$. For every possibly data-dependent set $E$ with
$|E|\le Mn/d$, on the aggregate event,
\begin{equation}\label{restore:fit-wrong-mass-bounds}
 \sum_{i\in E}G_i\le C_Mn,\qquad
 \sum_{i\in E}\sqrt{G_i}\norm{U_{i\cdot}}_2
       \le C_M\sqrt{n/d}.
\end{equation}
Moreover
$\sum_i\sqrt{G_i}\norm{U_{i\cdot}}_2\le C\sqrt{nd}$.
\end{lemma}

\begin{proof}
The elementary inequality
$G_i\le2D_{0,i}+(G_i-D_{0,i})^2/G_i$ gives
$\sum_{i\in E}G_i\le C(d|E|+n)\le C_Mn$.
Put $a_i=\sqrt{G_i}\norm{U_{i\cdot}}_2$.
By \eqref{restore:fit-basis} and Cauchy--Schwarz,
\[
 \sum_{i\in E}a_i
 \le C|E|\sqrt{d/n}
      +\left(\sum_{i\in E}G_i\right)^{1/2}\norm{U-V}_F
 \le C_M\sqrt{n/d}.
\]
The last assertion follows from
$\sum_i a_i\le(\sum_iG_i)^{1/2}\norm U_F$ and fixed $K$.
All estimates hold simultaneously for all such sets; they impose no
independence requirement on the selected labels.
\end{proof}

\begin{lemma}[Radius-count initialization under aggregate geometry]
\label{restore:fit-initializer}
The deterministic radius-count initializer specified earlier has at most
$C n/d$ errors on the aggregate event, after one common permutation.
This assertion requires neither a uniform degree upper bound nor
uniform spectral-row incoherence.
\end{lemma}

\begin{proof}
The normalized rows of $V$ are the $K$ orthogonal population centers.
Choose fixed $b_0>2\max_iD_{0,i}/d$, uniformly over the structural class.
For every $G_i>b_0d$,
$(G_i-D_{0,i})^2/G_i\ge G_i/4\ge b_0d/4$.
Thus at most $Cn/d$ such rows occur by
\eqref{restore:fit-geometry-two}. They contribute at most $4Cn/d$
to the squared directional error. For the remaining rows,
$\norm{V_{i\cdot}}_2\ge c/\sqrt n$. The row-normalization inequality
and \eqref{restore:fit-basis} bound their total squared directional
error by $Cn\norm{U-V}_F^2\le Cn/d$.
The existing radius-count lemma now applies with total squared error
$Cn/d$ and fixed positive community proportions. Its proof is
deterministic and unchanged.
\end{proof}

\begin{lemma}[Degree-corrected block totals and absolute mass]
\label{restore:fit-block-totals}
For $R^h$ as in \eqref{restore:reconstruction},
\begin{equation}\label{restore:fit-global-absolute}
 \sum_{i,j}|R^h_{ij}|\le Cnd.
\end{equation}
Uniformly over every fitted partition $\widehat Z$ having at most
$Mn/d$ errors under the common label matching,
\begin{equation}\label{restore:fit-total-accuracy}
 [\widehat Z^\top R^h\widehat Z]_{ab}
       =n_an_bB_{ab}+O_M(n).
\end{equation}
The omitted self-loops are included in the remainder.
\end{lemma}

\begin{proof}
Let $S_{\rm tail}=(I-\Pi)S(I-\Pi)$ and
$F_1=(I-\Pi)G^{1/2}\one$. Spectral invariance and the signal gap give
$\norm{S_{\rm tail}}_{\op}\le C/\sqrt d$ and
$\norm{F_1}_2\le C\sqrt n$.
Writing $v=\mathbf d-R_0\one$, one has exactly
$v=G^{1/2}S_{\rm tail}F_1$. Hence
\[
 \sum_i|v_i|
 \le\left(\sum_iG_i\right)^{1/2}
          \norm{S_{\rm tail}F_1}_2\le Cn.
\]
Since $\Pi_{ii}\le1$ and $\sum_i\Pi_{ii}=K$,
\[
 \sum_i|(R_0)_{ii}|
 \le\sum_iG_i\Pi_{ii}
 \le2(\max_iD_{0,i})K
       +\sum_i\frac{(G_i-D_{0,i})^2}{G_i}
 \le C(d+n)\le Cn.
\]
Consequently $\sum_i|(R_1)_{ii}|\le Cn$ and the simple hollowing
has total absolute cost at most $2\sum_i|(R_1)_{ii}|\le Cn$.
With $a_i$ as in the preceding lemma,
$|(R_0)_{ij}|\le a_ia_j$, so
$\sum_{i,j}|(R_0)_{ij}|\le(\sum_i a_i)^2\le Cnd$.
The rank-one correction has total absolute mass $\sum_i|v_i|\le Cn$.
This proves \eqref{restore:fit-global-absolute}.

Put $\widehat F=(I-\Pi)G^{1/2}\widehat Z$. By
\eqref{restore:fit-wrong-mass-bounds},
\[
 \norm{\widehat F}_F
 \le\norm{(I-\Pi)G^{1/2}Z}_F
       +\norm{G^{1/2}(\widehat Z-Z)}_F
 \le C_M\sqrt n .
\]
The exact quadratic identity
\[
 \widehat Z^\top(R_0-A)\widehat Z
       =-\widehat F^\top S_{\rm tail}\widehat F
\]
therefore has norm at most $C_Mn/\sqrt d$. Moving the mislabeled
vertices changes a raw block total by at most a fixed multiple of
$\sum_{i\in E}D_i\le\sum_{i\in E}G_i\le C_Mn$.
Equation \eqref{restore:fit-block-event} and the exact within-block
mean normalization of $\theta$ prove \eqref{restore:fit-total-accuracy}
with $R_0$ in place of $R^h$.
The rank-one correction changes a fitted block total by
$[\widehat Z^\top v]_a\widehat n_b/n$, where
$\widehat Z^\top v=\widehat F^\top S_{\rm tail}F_1=O_M(n/\sqrt d)$.
The already bounded hollowing cost completes the proof.
\end{proof}

\subsection{A bounded, normalized degree-parameter estimator}

Choose fixed positive $\alpha$ sufficiently small and $\beta$ sufficiently
large, using the structural bounds, so that on the two-sided
total-degree event
\[
 \alpha\overline D\le\tfrac12\min_iD_{0,i},\qquad
 \beta\overline D\ge2\max_iD_{0,i}.
\]
Define
\begin{equation}\label{restore:fit-clamped-theta}
 T_i=\min\{\max(D_i,\alpha\overline D),\beta\overline D\},
 \qquad
 \widehat\theta_i=
 \frac{T_i}{\widehat n_{\widehat z_i}^{-1}
                 \sum_{j:\widehat z_j=\widehat z_i}T_j}.
\end{equation}
This changes the nuisance estimator only. The retained degrees and the
degree identity $R^h\one=\mathbf d$ still use the original raw degrees.
The fitted parameters lie in $[\alpha/\beta,\beta/\alpha]$, and their
mean in every fitted block is exactly one.

\begin{lemma}[Clamped nuisance estimates]
\label{restore:fit-theta}
Uniformly over the $Mn/d$ count basin,
\begin{equation}\label{restore:fit-theta-rates}
 \frac{\norm{\widehat\theta-\theta}_2}{\sqrt n}
        \le C_Md^{-1/2},\qquad
 \widehat n_a^{-1}\sum_{\widehat z_i=a}T_i
        =dS_a+O_M(1),
\end{equation}
where $D_{0,i}=dS_{z_i}\theta_i$.
\end{lemma}

\begin{proof}
The chosen clamp interval contains $D_{0,i}$ with a fixed margin.
A direct split into the unclamped, lower-clamped, and upper-clamped
cases gives
\[
 |T_i-D_i|\le C\frac{(D_i-D_{0,i})^2}{G_i},\qquad
 |T_i-D_{0,i}|^2\le Cd\frac{(D_i-D_{0,i})^2}{G_i}.
\]
For the first inequality, an active lower clamp implies a degree
deviation of order $d$, and an active upper clamp implies a deviation
of order $D_i$. For the second, the unclamped range has $G_i\le Cd$;
projection onto an interval containing $D_{0,i}$ cannot increase
the deviation. Above the upper clamp,
$D_i((T_i-D_{0,i})/(D_i-D_{0,i}))^2\le T_i\le Cd$.
Summing yields
\begin{equation}\label{restore:fit-clamp-cost}
 \sum_i|T_i-D_i|\le Cn,\qquad
 \sum_i|T_i-D_{0,i}|^2\le Cnd.
\end{equation}
Since $\one=Z\one_K$, the raw block-total event implies
$\sum_{z_i=a}(D_i-D_{0,i})=O(n)$.
Clamping changes that sum by $O(n)$, and moving $Mn/d$ labels changes
a clamped degree total by at most $Cd(Mn/d)=O_M(n)$.
The fitted size differs from $n_a$ by $O_M(n/d)$.
Together with $\sum_{z_i=a}D_{0,i}=n_adS_a$, these facts prove the
group-mean assertion in \eqref{restore:fit-theta-rates}.
For correctly labeled vertices, division of
$T_i-D_{0,i}$ by this fitted mean, followed by
\eqref{restore:fit-clamp-cost}, gives squared error $O_M(n/d)$.
The incorrectly labeled vertices number $O_M(n/d)$ and have bounded
theta error. This proves the first assertion.
\end{proof}

\subsection{The exact Bernoulli block likelihood}

Fit a positive symmetric $B$ by maximizing the ordered Bernoulli
working criterion with $R^h$ and \eqref{restore:fit-clamped-theta},
divided by two, under the same fixed probability cap
$\widehat\theta_i\widehat\theta_jB_{\widehat z_i\widehat z_j}
\le q_{\max}<1$.
The next argument uses aggregate absolute mass instead of the stronger
and unnecessary assertion that every reconstructed entry is less than one.

\begin{lemma}[Unique Bernoulli block fit and its rate]
\label{restore:fit-block-root}
The fitted block roots are unique and interior for all large $n$.
Uniformly over the $Mn/d$ basin,
\begin{equation}\label{restore:fit-B-rate}
 \max_{a,b}\left|\frac{\widehat B_{ab}}{B_{ab}}-1\right|
       \le C_M(d^{-1}+p)\le C_Md^{-1}.
\end{equation}
\end{lemma}

\begin{proof}
For the ordered pairs $\mathcal J_{ab}$ associated with an unordered
block pair, put
\[
 E_{ab}=\sum_{\mathcal J_{ab}}R^h_{ij},\qquad
 t_{ij}=\widehat\theta_i\widehat\theta_j,\qquad
 T_{ab}=\sum_{\mathcal J_{ab}}t_{ij},\qquad
 \nu_{ab}=2-\mathbf1\{a=b\}.
\]
The preceding block-total lemma and exact fitted normalization imply
\[
 E_{ab}=\nu_{ab}n_an_bB_{ab}+O_M(n)\asymp n^2p>0,\qquad
 T_{ab}=\nu_{ab}\widehat n_a\widehat n_b+O(n).
\]
The exact score equation is $E_{ab}=F_{ab}(b)$, where
\[
 F_{ab}(b)=b\sum_{\mathcal J_{ab}}
              \frac{(1-R^h_{ij})t_{ij}}{1-bt_{ij}}.
\]
For every admissible $b$, bounded positive theta, the probability cap,
and \eqref{restore:fit-global-absolute} give
\[
 F_{ab}'(b)
 =\sum_{\mathcal J_{ab}}
          \frac{(1-R^h_{ij})t_{ij}}{(1-bt_{ij})^2}
 \ge cn^2-C\sum_{\mathcal J_{ab}}|R^h_{ij}|
 \ge cn^2-Cnd>0.
\]
Here $p=d/n\to0$. At zero, $F_{ab}(0)=0$; its value at the
positive cap endpoint is of order $n^2$, while $E_{ab}\asymp n^2p$.
There is a unique interior root. The likelihood derivative is
$(E_{ab}-F_{ab}(b))/b$, so the root is its unique global maximum.

For $b\asymp p$,
\[
 \sum_{\mathcal J_{ab}}
       \frac{(1-R^h_{ij})t_{ij}}{1-bt_{ij}}
       =T_{ab}+O(n^2p).
\]
The error follows by expanding the denominator and using
$\sum|R^h_{ij}|\le Cnd=Cn^2p$.
Bracketing the root between fixed multiples of $p$, then using this
expansion, gives
$\widehat B_{ab}=(E_{ab}/T_{ab})\{1+O(p)\}$.
The block-total and group-size errors prove the first inequality in
\eqref{restore:fit-B-rate}. The second uses $dp=np^2\le C_1$.
\end{proof}

\subsection{A local row condition suffices for every fitted stage}

For a tested vertex $i$, impose only the local conditions
\begin{equation}\label{restore:fit-local-row}
 D_i\le bd,\qquad
 a_i:=\sqrt{G_i}\norm{U_{i\cdot}}_2\le L\sqrt{d/n},
\end{equation}
where $b,L$ are fixed. These are not imposed on every neighbor
simultaneously.

\begin{lemma}[Local absolute mass and wrong-label stability]
\label{restore:fit-local-mass}
On the aggregate event and \eqref{restore:fit-local-row},
\begin{equation}\label{restore:fit-local-mass-bounds}
 \sum_j|R^h_{ij}|\le C_{b,L}d,\qquad
 \sup_{|E|\le Mn/d}\sum_{j\in E}|R^h_{ij}|\le C_{M,b,L}.
\end{equation}
\end{lemma}

\begin{proof}
Since $|(R_0)_{ij}|\le a_ia_j$,
\[
 \sum_j|(R_0)_{ij}|\le a_i\sum_j a_j\le Cd,\qquad
 \sum_{j\in E}|(R_0)_{ij}|
       \le a_i\sum_{j\in E}a_j\le C_M
\]
by \cref{restore:fit-wrong-mass}. Also
$|v_i|\le D_i+\sum_j|(R_0)_{ij}|\le Cd$.
The rank-one correction therefore adds at most $Cd$ to the first
bound and $|v_i||E|/n\le C_M$ to the second.
Finally $|(R_1)_{ii}|\le a_i^2+|v_i|/n\le Cp$.
The simple row hollowing changes either row bound by at most
$2|(R_1)_{ii}|\le Cp$. This proves the claim.
\end{proof}

\begin{lemma}[Bernoulli profile reduction for a signed row]
\label{restore:fit-profile}
Let $r$ be a hollow signed row satisfying
$\sum_jr_j=D_i\ge0$ and $\sum_j|r_j|\le Ld$.
Let the positive neighbor intensities $b_j$ lie between fixed
multiples of $p$. Maximizing the exact Bernoulli row criterion
\[
 \sum_j\{r_j\log(tb_j)+(1-r_j)\log(1-tb_j)\}
\]
over $t\ge0$ under the fixed cap $tb_j\le q_{\max}<1$ gives the
usual normalized block-profile contrast, with remainder $O(np^2)$.
No entrywise bound on $r_j$ is required.
\end{lemma}

\begin{proof}
After combining the coefficient of $\log t$, the derivative is
\[
 D_i/t-\sum_j\frac{(1-r_j)b_j}{1-tb_j}.
\]
Throughout the admissible interval,
\[
 \sum_j\frac{(1-r_j)b_j}{1-tb_j}
 \ge cnp-Cp\sum_j|r_j|
 \ge cd-CLdp\ge c'd.
\]
Thus every maximizing $t$ lies in a fixed bounded interval,
because $D_i\le Ld$. This reasoning does not require individual
nonedge weights $1-r_j$ to be nonnegative.
On that fixed interval, the exact Bernoulli expansion is uniform:
\[
 \ell_r(t)
   =D_i\log t+\sum_jr_j\log b_j
         -t\sum_jb_j+O(np^2).
\]
Restoring the edge terms in the linear exposure costs
$Cp\sum_j|r_j|=O(np^2)$, and the quadratic remainder costs
$Cp^2(n+\sum_j|r_j|)=O(np^2)$.
The leading optimizer $D_i/\sum_jb_j$ is bounded and admissible
for all large $n$, so taking maxima preserves the remainder.
The resulting label-dependent maximum is the standard normalized
block-profile score. If $D_i=0$, the derivative bound makes $t=0$
optimal after combining the $\log t$ terms; use the continuous
convention $0\log0=0$.
\end{proof}

\begin{proposition}[Uniform replacement by the true DC contrast]
\label{restore:fit-uniform-score}
\label{restore:fit-score}
Fix $M,b,L$. Use the clamped theta estimates
\eqref{restore:fit-clamped-theta}, the exact Bernoulli block fit, and
the exact Bernoulli target profile on $R^h$.
On the aggregate event and the local target condition
\eqref{restore:fit-local-row}, uniformly over every fitted partition
in the $Mn/d$ basin,
\begin{equation}\label{restore:fit-true-score}
 \widehat L_i^{ab}
       =[R^h w_{i,ab}]_i+O_{M,b,L}(1),\qquad
 (w_{i,ab})_j=
     \log\frac{q_{a,z_j}^{(i)}}{q_{b,z_j}^{(i)}}\quad(j\ne i).
\end{equation}
The true profiles $q^{(i)}$ are the self-excluded neighbor profiles
defined in \eqref{dcsbm:q-true}. The value of $w_{i,ab}$ at $i$ is immaterial since
$R^h_{ii}=0$.
\end{proposition}

\begin{proof}
By \cref{restore:fit-profile}, the profile reduction costs
$O(np^2)=O(1)$. Exact fitted-block normalization gives the fitted
neighbor mass
$\widehat M_c^{(i)}
 =\widehat n_c-\widehat\theta_i\mathbf1\{\widehat z_i=c\}$.
The fitted sizes differ from the true sizes by $O_M(n/d)$,
the fitted theta values are bounded, and
\eqref{restore:fit-B-rate} gives relative block error $O_M(1/d)$.
Smoothness of normalization and logarithms on the fixed positive
profile class therefore yields fitted log-profile error $O_M(1/d)$.
Multiplication by the row absolute mass in
\eqref{restore:fit-local-mass-bounds} costs $O_M(1)$.
Replacing the fitted labels by the true labels changes a bounded
coefficient vector on at most $Mn/d$ coordinates; the second bound
in \eqref{restore:fit-local-mass-bounds} gives another $O_M(1)$.
All inequalities hold simultaneously for every partition in the basin,
so data-dependent fitting does not introduce an independence assumption.
\end{proof}

\begin{lemma}[A uniform basin for finite hard refits]
\label{restore:fit-invariant-basin}
\label{restore:fit-invariant}
There is a fixed $M$ such that, on the aggregate event and for all
sufficiently large $n$, each hard nuisance-fit and profile-assignment
update maps the $Mn/d$ error basin into itself under the same label
matching. Outside a fixed target-exception set determined by the graph,
every error at every updated iterate is contained in
\[
 [R^h w_{i,ab}]_i\le C_M
 \quad\text{for some }b\ne z_i.
\]
The target-exception set can be taken to be the failure of the local
conditions \eqref{restore:fit-local-row}, with fixed sufficiently large
thresholds. Thus a finite number of these specified hard refits, including
adaptive finite stopping after at least one update, requires no union
over iterations in the final risk bound.
\end{lemma}

\begin{proof}
Fix sufficiently large constants $b,R$. The set
$\{i:D_i>bd\}$ has at most $Cn/d$ elements by the degree energy in
\eqref{restore:fit-geometry-two}, since $b$ exceeds twice the largest
population degree multiplier. Since $G_i\ge cd$,
\eqref{restore:fit-basis} gives $\norm{V_{i\cdot}}_2\le C/\sqrt n$.
For $R>2C$, the set $\{i:\sqrt n\norm{U_{i\cdot}}_2>R\}$
also has at most $Cn/d$ elements, by
$\norm{U-V}_F^2\le C/d$.
Every other target satisfies \eqref{restore:fit-local-row} for a
fixed $L$, because $G_i\le C_bd$ there.
This target-exception set depends only on the graph and its
cardinality bound is independent of $M$.

Write $M_G=G^{-1/2}\Omega G^{-1/2}$ and
$S_K=U\Lambda U^\top$. Best rank-$K$ approximation and
\eqref{restore:fit-geometry-one} imply
$\norm{S_K-M_G}_{\op}\le C/\sqrt d$. Hence
\[
 \norm{G^{-1/2}(R_0-\Omega)Z}_F^2
 \le \frac Cd\norm{G^{1/2}Z}_F^2\le Cn.
\]
Moreover
$\norm{G^{-1/2}v}_2^2
 =\norm{S_{\rm tail}F_1}_2^2\le Cn/d$.
The rank-one correction therefore preserves the bound
\begin{equation}\label{restore:fit-basin-action}
 \norm{G^{-1/2}(R_1-\Omega)Z}_F^2\le Cn.
\end{equation}
The simple hollowing changes any bounded block-constant contrast at
row $i$ by at most $C|(R_1)_{ii}|$.
Since $\sum_i|(R_1)_{ii}|\le Cn$, at most $Cn/d$ rows can incur
a hollowing contrast error larger than any prescribed fixed small
multiple of $d$.
This auxiliary counting set is needed only for the deterministic
basin proof; it is not added as a probability-cost event in the sharp
risk argument.

The true profile contrast has a uniform positive population margin
$[\Omega w_{i,ab}]_i\ge c_{\rm gap}d$ for $a=z_i$, $b\ne a$,
including its $O(p)$ self-exclusion correction.
Extend each $w_{i,ab}$ to a block-constant vector at coordinate $i$;
this does not affect $R^h w$ because $R^h_{ii}=0$.
Its $K$ block coefficients are uniformly bounded. For a regular
target, \cref{restore:fit-uniform-score} bounds the fitted discrepancy
by $C_M$, which is less than $c_{\rm gap}d/4$ for all sufficiently
large $n$ (the threshold may depend on $M$).
Outside the auxiliary hollowing set, a new erroneous label therefore
forces
\[
 \norm{[(R_1-\Omega)Z]_{i\cdot}}_2\ge c d.
\]
For the regular targets, $G_i\le C_bd$. Summing the squared
requirement divided by $G_i$ and applying
\eqref{restore:fit-basin-action} gives at most $C_0n/d$ such targets.
Including the two cardinality exceptions still gives $C_0n/d$
with $C_0$ independent of $M$.

Choose $M$ at least both this $C_0$ and the initializer constant.
The basin is invariant for all sufficiently large $n$.
For every regular target, the canonical error-event inclusion
follows directly from \eqref{restore:fit-true-score}, independently
of this counting argument. It is the same event for all partitions
in the fixed basin. Induction therefore covers each finite hard
refit and any data-dependent finite stopping rule without a new
probability union or a stopping-time independence claim.
\end{proof}


\section{Completion of the current DC risk bound}
\begin{proof}[Proof of the regularized DC risk bound below a fixed logarithmic scale]
Fix any constant $A_0<\infty$ and first work on
$d\le A_0\log n$, still with $d\to\infty$.
Let $\mathcal E$ collect the aggregate geometric conditions of
\cref{prop:sparse-floor}, their recomputed versions after deletion of
at most three vertices, two-sided total degree, and the true block-total event used
in the fitting lemmas. None imposes a simultaneous lower bound on all
raw degrees or a simultaneous incoherence bound. For every prescribed
fixed $r$, its failure probability is $O(n^{-r})$ after increasing the
fixed constants. For the block totals this also follows directly from
the independent-edge Bernstein bound at a sufficiently large constant
times $n$: its exponent is at least $cn/d$, and $np^2\le C_1$ gives
$n/d\ge d/C_1$; in the present range it is much larger than $\log n$.

The initializer and invariant-basin result in
\cref{restore:fit-invariant} apply deterministically on $\mathcal E$.
For a tested target $i$, impose only its upper-degree bound
$D_i\le Ld$, its local normalized row bound, and the condition that its
incident row hits no exceptional deleted eigenvector row. By
\cref{restore:tail-energy-lemma,restore:tail-hit-corollary,restore:tail-target-corollary}, the average probability cost
of the two row conditions is $C_He^{-Hd}$ for any prescribed fixed $H$,
after choosing their fixed thresholds sufficiently large. For the
LOO hit condition, the exceptional set consists of deleted rows with
$\sqrt n\norm{(U_B)_{j\cdot}}_2>R$; the floor implies
\eqref{restore:loo-hit} on its complement. Its expected hit mass
is bounded by the row-count argument in
\eqref{restore:tail-coefficient-count}.
The upper
degree tail has the same property by the single-row Bernoulli Chernoff
bound. Polynomial factors in $d$ are absorbed by choosing a larger
exponential rate at the preceding step. No union over the $n$ tested
targets is taken: their error probabilities are averaged.

On these local events, the fitted-score lemma
\cref{restore:fit-score} replaces every estimated, possibly
data-dependent contrast by a deterministic true block contrast at
$O(1)$ cost. The local LOO comparison
\cref{restore:loo-corrected} replaces its actual full-graph spectral
action by an independent training action at another $O(1)$ cost.
For the DC profile contrast, the new score is the true raw-row
linear profile contrast plus
\[
 \sum_{j\ne i}A_{ij}e_j+O(1),\qquad
 e=(H_B-I)(w-\bar w\one).
\]
The coefficients are measurable from the deleted graph and independent
of the incident Bernoulli row. For each target use the analogous
training event $\mathcal E_i$ formed only from that deleted graph and
its further deletions. In applying the independent-row bound,
drop the full-graph event after the deterministic error inclusion and
condition only on the training graph on $\mathcal E_i$.
The geometry cost is
\[
 \varepsilon_{\rm geom}
 =\PP(\mathcal E^c)+\frac1n\sum_i\PP(\mathcal E_i^c)
 \le C_rn^{-r}.
\]
This does not condition the Bernoulli row on a full-graph event.
The tilted mean and quadratic cost are
bounded by \eqref{eq:floor-tilt-moments}; here $np^2$ is bounded.
The exceptional coefficient hit-mass hypothesis of
\cref{lem:sparse-truncated-transfer} is now a conclusion of
\cref{restore:tail-hit-corollary}, rather than an unproved premise. That sharp
transfer gives, after averaging the local target exceptions and taking
a fixed-$K$ union over competitors,
\[
 \E\ell(\widehat z,z)
 \le \frac Cn\sum_i\sum_{b\ne z_i}
                  d^{-1/2}e^{-J_{i,z_i\to b}}
       +C_He^{-Hd}+\varepsilon_{\rm geom}.
\]
Here $J_{i,a\to b}$ is the exact directional profile exponent in
\eqref{dcsbm:profile-exponent}. By \cref{dcsbm:profile-tail}, the same
pairwise directional sharp scale is a lower bound, up to a constant,
for the actual profile oracle. Since $K$ is fixed, the
displayed sum is bounded by a constant times its average risk.
Only the averaged target-row tail is used; no pointwise uniform tail
for heterogeneous degree parameters is asserted.

Uniformly on the compact DC class, the reference average risk is at
least $c d^{-1/2}e^{-C_Jd}$. Choose $H>C_J+1$ and then choose fixed
$r>A_0(C_J+1)+2$. For $d\le A_0\log n$,
\[
 \frac{e^{-Hd}+n^{-r}}{d^{-1/2}e^{-C_Jd}}\longrightarrow0.
\]
This closes the unconditional expected-risk estimate throughout the
entire range below any fixed multiple of $\log n$. In particular, no
positive lower coefficient of $\log n$ is required.
\end{proof}

\begin{proof}[Completion over the remaining degree scales]
Fix $A_0>0$ and use the preceding sparse-range argument when
$d\le A_0\log n$. For $d>A_0\log n$, replace its polynomially
accurate aggregate event by $\mathcal E_H$ from
\cref{restore:high-aggregate}. The proof of
\cref{restore:tail-energy-lemma} uses only the deterministic aggregate
geometry on the event; the innovation concentration is applied for
every training graph before restricting to the event. It therefore
holds unchanged with $\mathcal E_H$ and the corresponding
training-measurable events. Its coefficient-hit and target-row
corollaries, the scalar LOO comparison, all fitted-score bounds,
and the invariant basin likewise apply.

The resulting risk upper bound is the appropriate oracle term plus
$C_{H'}e^{-H'd}+C_He^{-Hd}$, where $H'$ can be prescribed by choosing
the fixed local thresholds. On the fixed positive compact DC model
class the oracle risk is at least $cd^{-1/2}e^{-C_Jd}$.
Choosing $H,H'>C_J+1$ pays both remainders. This proves the stated
all-scale profile-oracle bound. The known-target-degree comparison follows
from the separate equal-exposure calculation in
\cref{dcsbm:equal-exposure}. No logarithmic degree lower bound is imposed.
\end{proof}

\clearpage
\part{Earlier decoders and additional results}
\section{Historical unregularized encoders and their separate results}\label{sec:historical}

This part preserves the separate analyses of the earlier unregularized
algorithms. Its degree-range assumptions describe those historical
encoders only; the current regularized theorem has no logarithmic
lower bound. Their logarithmic-range hypotheses
are not hypotheses of \cref{thm:regularized}. Their positive-degree
fallback convention is replaced by the explicit floor in the primary
algorithm.

When all degrees are positive, form
\[
 S=D^{-1/2}AD^{-1/2}=I-L.
\]
Retain the degree eigenvector $D^{1/2}\one/\norm{D^{1/2}\one}_2$ and the
$K-1$ other eigenvectors with largest absolute eigenvalues in its orthogonal
complement. Let $U^\top U=I_K$ and $SU=U\Lambda$, where
$\Lambda=\diag(\mu_1,\ldots,\mu_K)$. On the proved event this is exactly
the informative $K$-dimensional space. The encoder supplies only
$(D,U,\Lambda)$ to the decoder. Put
\begin{equation}\label{eq:encoder-action}
 X=D^{-1/2}U,\qquad AX=D^{1/2}U\Lambda.
\end{equation}
Thus row actions on the retained space require no later edge lookup.
The normalized Laplacian eigenvalues contain the same information via
$\mu=1-\lambda_L$. Exact arithmetic and algebraic retained spaces define
the theorem; a fixed numerical rank threshold is not its definition.

The ordinary-SBM decoder uses the \emph{signed} reconstruction
\begin{equation}\label{eq:euclidean-decoder}
 Q=\operatorname{orth}(X),\quad \Pi_X=QQ^\top,\quad R_E=A\Pi_X.
\end{equation}
It fits free Bernoulli row profiles with the finite growing/replacement
schedule and mandatory final E/M step specified in
\cref{sec:sbm-proof}. Only probabilities are clipped; $R_E$ is not clipped.
Neither true labels nor a correct initial partition are inputs.

For DC-SBM we make the explicit degree-metric correction
\begin{equation}\label{eq:metric-decoder}
 H_D=D^{-1/2}UU^\top D^{1/2},\qquad
 R_D=AH_D=D^{1/2}U\Lambda U^\top D^{1/2}.
\end{equation}
With $h_i=\norm{U_{i\cdot}}_2^2$ and
$s_i=(R_D)_{ii}/(1-h_i)$, its hollow version is
\begin{equation}\label{eq:hollow-decoder}
 R^\circ=R_D+\diag(s)(H_D-I).
\end{equation}
It satisfies
$\diag(R^\circ)=0$, $R^\circ\one=D\one$ and $R^\circ X=AX$.
It need not be symmetric after hollowing; a symmetric block parameter is
fitted with the ordered-row working criterion divided by two.

For the initial partition, normalize every row of $U$ to unit length.
Repeatedly select a remaining row having the largest number of remaining
rows within distance $1/8$, and remove rows within distance $1/4$ of the
selected row. After $K$ selections assign every row to the closest selected
row. The separation argument in \cref{sec:dcsbm-proof} proves its
$O(1/d)$ error; this is a deterministic procedure with $O(Kn^2)$ distance
comparisons, not an assumed global k-means solution.
Normalize degree estimates to mean one within each fitted group, fit
the symmetric block matrix by the exact Bernoulli working criterion, and
perform at least one exact Bernoulli profile-LRT assignment using
$R^\circ$. One pass suffices. The same theorem permits a finite number of
the stated hard refits, including adaptive stopping, by the invariant-basin
argument proved below.

Both algorithms return a fixed partition, or their initial partition,
if a prescribed operation is undefined. This includes zero degrees,
insufficient retained rank, a hollow denominator zero, an empty initial
group, or an inadmissible parameter fit. The required operations are all
well-defined on the good event, so this convention costs only its actual
failure probability.

\begin{theorem}[Earlier ordinary free-profile decoder]\label{thm:sbm}
Fix $K\ge2$, $\pi_->0$ and $\kappa>0$. Suppose
\[
 n_a\ge\pi_-n,\qquad
 p/2\le P_{ab}\le p\le1/4,\qquad
 \sigma_{\min}(P)\ge\kappa p,\qquad d=np\ge5\log n.
\]
Use \eqref{eq:euclidean-decoder} and the finite growing/replacement
likelihood decoder, with probability floor
$\varepsilon=c_0\max(1,n^{-1}\sum_iD_i)/n$ for a fixed sufficiently small
$c_0>0$ (for example any $c_0<1/16$ on this class).
There are constants $C,n_0$, depending only on the displayed fixed
structural constants and $c_0$, such that for $n\ge n_0$,
\begin{equation}\label{eq:sbm-risk}
 \E\ell(\widehat z_E,z)\le C R_{\mathrm{SBM},*}(n).
\end{equation}
The algorithm uses neither $P$, $p$, nor the true partition.
This includes fixed arbitrary $K$, unbalanced group sizes of comparable
order, and general symmetric full-rank block probability matrices in the
displayed range. It does not require all within-group probabilities, or all
between-group probabilities, to be identical.
\end{theorem}

\begin{theorem}[Earlier unregularized DC profile decoder]
\label{thm:dc}
Assume \eqref{eq:mean-model}--\eqref{eq:structural}. Fix $C_1<\infty$ and
suppose
\begin{equation}\label{eq:dc-range}
 C_0\log n\le d=np,\qquad np^2\le C_1,
\end{equation}
where $C_0$ is sufficiently large depending only on the fixed structural
bounds and $C_1$. The degree-metric decoder
\eqref{eq:metric-decoder}--\eqref{eq:hollow-decoder}, with the deterministic
initialization and estimated-parameter profile procedure just specified,
satisfies
\begin{equation}\label{eq:dc-risk}
 \E\ell(\widehat z_D,z)\le C R_{\mathrm{prof},*}(n)
\end{equation}
for all sufficiently large $n$, with $C$ independent of $n$.
The statement is uniform over the displayed compact parameter class.
No correct grouping, true degree parameters, true block matrix, sample
splitting, or leave-one-out eigensystem is an input to the algorithm.
\end{theorem}

\begin{corollary}[The known-target-degree oracle on equal exposures]
\label{cor:known}
Under \cref{thm:dc}, suppose additionally that the actual finite-$n$ block
exposures satisfy
\begin{equation}\label{eq:equal-exposure}
 \sum_b(n_b/n)C_{ab}=S\quad\text{for all }a.
\end{equation}
Then
\[
 R_{\mathrm{prof},*}\asymp R_{\theta,*},\qquad
 \E\ell(\widehat z_D,z)\le C R_{\theta,*}.
\]
Thus there are nontrivial fixed-$K$ SBM and DC-SBM classes on which the
spectral-only fitted likelihood procedure attains a constant multiple of
the original known-parameter local LRT risk. The ordinary free-profile
decoder, with a sufficiently small floor coefficient for its structural
parameter class, also satisfies \eqref{eq:sbm-risk} for general positive
full-rank $C$ under \eqref{eq:dc-range} and \eqref{eq:equal-exposure},
without the entry-ratio-two restriction of \cref{thm:sbm}.
\end{corollary}

Exact equality in \eqref{eq:equal-exposure} can be relaxed to exposure
differences whose contribution to all likelihood contrasts is uniformly
$O(1)$. For example fixed $C\pi=S\one$ with
$n_a=n\pi_a+O(1)$ suffices. Equality of limiting exposures with no rate
is not by itself sufficient. The general known-$\theta_i$ claim without
an additional restriction is false: \cref{thm:imposs} gives an explicit
diverging risk ratio for every full-graph algorithm.

\begin{remark}[What is and is not constant]
The factor may depend on signal conditioning, group proportions, the
allowed degree-parameter range and $C_1$. It is independent of $n$ on
the stated class. We do not assert a small numerical constant, a limit of
the risk ratio, or uniformity when $\sigma_{\min}(C)\to0$. The latter
weak-signal extension needs separate quantitative estimates.
\end{remark}


\section{Historical unregularized spectral geometry and leave-one-out control}
\label{sec:geometry}

This section proves the spectral premises used by both fitted decoders.
Leave-one-out (LOO) matrices occur only in the proof. The implemented
encoder retains one full-graph eigensystem.

\subsection{Raw concentration and deterministic spectral consequences}

For any prescribed fixed $H>0$, when $d\ge c\log n$ there is a fixed
$L_H$ such that
\begin{equation}\label{eq:raw-norm-tail}
 \PP\{\norm{A-\Omega}_{\op}>L_H\sqrt d\}\le e^{-Hd}.
\end{equation}
Indeed, apply Corollary 3.12 and its centered-entry extension in
Remark 3.13 of Bandeira--van Handel~\cite{bvh} to $A-\E A$.
For bounded centered entries that result bounds the tail at
$C\sqrt{\max_i\sum_j\operatorname{Var}(A_{ij})}+t$ by
$n\exp(-t^2/C')$. Here the row variance sum is $O(d)$, and the entries
have absolute value at most one. Choosing $t$ a sufficiently large fixed
multiple of $\sqrt d$ absorbs both $H$ and $\log n$.
The completed-mean diagonal costs $O(p)$. The centered-entry extension is
needed: centered Bernoulli noise does not have a symmetric distribution.

Fix constants $a,b,L>0$ and define the base event
\begin{equation}\label{eq:base-event}
 \mathcal E_0=
 \{\norm{A-\Omega}_{\op}\le L\sqrt d,
       \quad ad\le D_i\le bd\text{ for every }i\}.
\end{equation}
Its minimum-degree probability is deliberately not included in
\eqref{eq:raw-norm-tail}; it is treated separately in
\cref{sec:probability}.
On $\mathcal E_0$,
\begin{equation}\label{eq:degree-aggregate}
 \norm{(D-\bar D)\one}_2\le L\sqrt{nd},\qquad
 \sum_j\frac{(D_j-\bar D_j)^2}{D_j}\le Cn,
 \quad \bar D=\diag(\Omega\one).
\end{equation}
These follow by applying $A-\Omega$ to $\one$; no uniform vanishing
relative degree error is required.

The observed-degree-normalized completed mean
\[
 M=D^{-1/2}\Omega D^{-1/2}
\]
has rank $K$, column space $\col(D^{-1/2}\Theta Z)$, and its nonzero
eigenvalues have absolute values at least a fixed positive constant.
For the last statement,
$Z^\top\Theta D^{-1}\Theta Z$ is diagonal with entries comparable to
$n/d$, and the nonzero eigenvalues are those of
\[
 p(Z^\top\Theta D^{-1}\Theta Z)^{1/2}
 C(Z^\top\Theta D^{-1}\Theta Z)^{1/2}.
\]
The lower singular-value assumption supplies the fixed gap.
Since $\norm{S-M}_{\op}\le C/\sqrt d$, the retained eigenvalues are
bounded away from zero, the discarded spectrum is $O(d^{-1/2})$, and
\begin{equation}\label{eq:normalized-angle}
 \norm{UU^\top-\Pi_M}_{\op}\le C/\sqrt d,
\end{equation}
where $\Pi_M$ projects onto $\col(D^{-1/2}\Theta Z)$.
These gap estimates follow directly by expanding a retained vector in the
signal and complementary eigenbases, or by the usual subspace perturbation
inequality. Positive and negative signal eigenvalues are both retained.

The same conclusions hold for every vertex-deleted graph with its degrees
\emph{recomputed}: deleting a vertex preserves the raw principal-submatrix
norm bound, reduces each surviving degree by at most one, and changes a
block mass by only a bounded amount. All constants can be chosen common
for the full graph and these $n$ minors.

\subsection{Exact generalized equations for one incident row}

Put the tested vertex first and write
\begin{equation}\label{eq:arrowhead-A}
 A=\begin{pmatrix}0&x^\top\\x&B\end{pmatrix},\qquad
 k=\norm{x}_1,\quad \Delta=\diag(B\one),\quad
 D=\diag(k,\Delta+E_x),\quad E_x=\diag(x).
\end{equation}
Here $x\in\{0,1\}^{n-1}$. Suppose the deleted degrees lie between
fixed multiples of $d$, $0<k\le bd$, and
\[
 S_B=\Delta^{-1/2}B\Delta^{-1/2}
    =U_B\Lambda_BU_B^\top+N,
 \quad U_B^\top U_B=I_K,\quad NU_B=0,
\]
with signal eigenvalues separated from zero and
$\norm{N}_{\op}\le C/\sqrt d$. Write
\[
 V=\Delta^{-1/2}U_B,\quad Q_\perp=I-U_BU_B^\top,\quad
 F=\Delta^{-1/2}E_x\Delta^{-1/2},\quad h=\Delta^{-1/2}x.
\]
Then $\norm F\le C/d$ and $\norm h^2\le Ck/d$.
The full normalized matrix differs from $\diag(0,S_B)$ by
$O(d^{-1/2})$: its off-diagonal normalized row has squared norm at most
$C/d$, and the lower degree-renormalization perturbation is $O(d^{-1})$.
The full retained signal therefore has the same fixed gap.

Take full generalized eigenvectors $g_r=D^{-1/2}u_r$, so
$Ag_r=\mu_rDg_r$ and jointly $G^\top DG=I_K$.
Decompose them exactly as
\[
 g_r=\binom{t_r}{\Delta^{-1/2}(U_By_r+z_r)},
 \qquad U_B^\top z_r=0.
\]
Let $h_V=U_B^\top h$ and $h_Q=Q_\perp h$.
The top row and complementary equations give
\begin{align}
 h^\top(U_By_r+z_r)&=\mu_r k t_r,\label{eq:generalized-top}\\
 (\mu_rI-N+\mu_rQ_\perp FQ_\perp)z_r
   &=t_rh_Q-\mu_rQ_\perp FU_By_r.\label{eq:generalized-complement}
\end{align}
The inverse $T_r$ of the complementary operator has bounded norm.
Substituting into \eqref{eq:generalized-top} gives
\begin{equation}\label{eq:top-resolvent}
 (\mu_r k-h_Q^\top T_rh_Q)t_r
 =h_V^\top y_r-\mu_rh_Q^\top T_rQ_\perp FU_By_r.
\end{equation}
Its denominator has magnitude at least $ck$: the second term is at most
$Ck/d$. This remains true when $k$ is small.

Put $M_B=\norm{U_B}_{2\to\infty}$.
The coordinates $y_r$ have bounded norm because
$\Delta^{1/2}G_{-i}=(I+F)^{-1/2}U_{-i}$ has norm at most one.
Using
\[
 \norm{h_V}\le CkM_B/\sqrt d,\qquad
 \norm{FU_B}_F\le CM_B\sqrt k/d
\]
in the exact equations yields
\begin{equation}\label{eq:generalized-bounds}
 |t_r|\le CM_B/\sqrt d,\qquad
 \norm{z_r}_2\le CM_B\sqrt k/d.
\end{equation}
The matrix $Y=(y_1,\ldots,y_K)$ has bounded inverse: the full and padded
deleted signal spaces have angle $O(d^{-1/2})$, and the extra multiplier
$(I+F)^{-1/2}$ costs only $O(d^{-1})$.

\begin{lemma}[Normalized row bounds from model events]\label{lem:normalized-rows}
Under the structural model assumptions, for any prescribed $H>0$ there
are row-concentration events of total failure probability at most
$2Kn e^{-Hd}$ such that, on their intersection with $\mathcal E_0$,
the full normalized retained vectors and every deleted normalized retained
basis satisfy
\begin{equation}\label{eq:row-incoherence}
 \norm{U}_{2\to\infty}\le C_H/\sqrt n,
 \qquad \max_i\norm{U^{(-i)}}_{2\to\infty}\le C_H/\sqrt n.
\end{equation}
They also satisfy the gap and angle bounds above. In particular, with
$\Pi_0=Z(Z^\top Z)^{-1}Z^\top$,
\begin{equation}\label{eq:euclidean-angle}
 \norm{\Pi_X-\Pi_0}_{\op}\le C_H/\sqrt d,
 \quad \norm{Q}_{2\to\infty}\le C_H/\sqrt n,
\end{equation}
and the same holds for all recomputed deleted projectors.
\end{lemma}

\begin{proof}
Let $M_\mathrm{full}=\norm U_{2\to\infty}$ and
$M_i=\norm{U^{(-i)}}_{2\to\infty}$.
The lower rows of the full eigenbasis are
$(I+F)^{1/2}(U_BY+Z_\perp)$, with
$Z_\perp=(z_1,\ldots,z_K)$.
Using $Y^{-1}$ and \eqref{eq:generalized-bounds} gives
\[
 M_i\le C M_\mathrm{full}+CM_i/\sqrt d,
 \qquad M_i\le C' M_\mathrm{full}\quad\text{for large }d.
\]
This argument has not assumed row incoherence.
For the tested full row, \eqref{eq:top-resolvent}, $k\asymp d$ and
$u_{i,r}=\sqrt{k}\,t_r$ imply the sharper estimate
\begin{equation}\label{eq:tested-row-bootstrap}
 \norm{U_{i\cdot}}_2
 \le\frac C{\sqrt d}\norm{x^\top\Delta^{-1/2}U_B}_2+CM_i/d.
\end{equation}

Conditional on $B$, its eigenbasis and recomputed degrees are fixed and
$x$ is an independent Bernoulli row. Its mean action has norm at most
$C\sqrt{d/n}$, by $\norm{\E x}_2\le Cp\sqrt n$ and
$\norm{\Delta^{-1/2}U_B}\le C/\sqrt d$.
We use the following weighted Bernoulli bound from
Abbe--Fan--Wang--Zhong~\cite[Lemma 7, Eq.\ (C.2)]{afwz}:
for fixed $v$ and independent Bernoulli coordinates with success
probabilities at most $q$,
\begin{equation}\label{eq:afwz-row}
 \PP\left\{\left|\sum_jv_j(X_j-\E X_j)\right|
 >\frac{(2+\alpha)nq\norm v_\infty}
 {1\vee\log(\sqrt n\norm v_\infty/\norm v_2)}\right\}
 \le2e^{-\alpha nq}.
\end{equation}
Apply it to each of the fixed $K$ columns of $\Delta^{-1/2}U_B$, taking
$q=\min(Cp,1)$, which is comparable to $p$ on the structural class,
and selecting $\alpha$ to give rate $Hd$.
The column norms are at most $C/\sqrt d$ and $CM_i/\sqrt d$ in
$\ell_2$ and $\ell_\infty$, respectively. The function
$b/(1\vee\log(\sqrt n b/s))$ is nondecreasing in both admissible norm
arguments. Thus outside conditional probability $2Ke^{-Hd}$,
\[
 \norm{x^\top\Delta^{-1/2}U_B}_2
 \le C\sqrt{d/n}+
 \frac{C_H\sqrt d\,M_i}{1\vee\log(c\sqrt nM_i)}.
\]
A union over $i$, followed by \eqref{eq:tested-row-bootstrap}, yields
\[
 M_\mathrm{full}\le C/\sqrt n+
 \frac{C_H'M_\mathrm{full}}
 {1\vee\log(c'\sqrt nM_\mathrm{full})}
 +C'M_\mathrm{full}/d.
\]
For large $d$, absorb the last term. If $\sqrt nM_\mathrm{full}$ were
larger than a sufficiently large fixed constant, the logarithmic term
could also be absorbed, giving a contradiction. This proves
\eqref{eq:row-incoherence} and its deleted versions.

The conditioning is legitimate even though $\mathcal E_0$ depends on
$x$: define the conditional row event for every deleted graph satisfying
its measurable degree and gap premises, apply \eqref{eq:afwz-row} there,
and use $\mathcal E_0$ only for deterministic implications afterwards.

For \eqref{eq:euclidean-angle}, bounded-condition-number multiplication by
$D^{-1/2}$ in \eqref{eq:normalized-angle} compares $\col(X)$ with
$\col(D^{-1}\Theta Z)$ at angle $O(d^{-1/2})$.
Multiplying block $a$ of $D^{-1}\Theta Z$ by $dS_a/\sqrt{n_a}$ gives
entries $(\bar D_i/D_i)/\sqrt{n_a}$ on that block. By
\eqref{eq:degree-aggregate}, this basis differs from
$Z\diag(n_a^{-1/2})$ in Frobenius norm by $O(d^{-1/2})$.
This proves the angle to the block space. Finally,
$Q=X(X^\top X)^{-1/2}$ has the stated row bound since the eigenvalues of
$X^\top X=U^\top D^{-1}U$ are comparable to $1/d$.
The argument applies to DC-SBM as well as SBM.
\end{proof}

\subsection{A deterministic full-to-LOO score lemma}

\begin{lemma}[Both actual full-graph reconstructions have bounded LOO error]
\label{lem:loo}
Assume the deleted degree, spectral gap and discarded-spectrum conditions
following \eqref{eq:arrowhead-A}, as well as
$\norm{U_B}_{2\to\infty}\le C/\sqrt n$ and $k\le bd$.
For every $w\in\R^n$ with $\norm w_\infty\le W$, let
\[
 \Pi_B=V(V^\top V)^{-1}V^\top,
 \qquad H_B=VV^\top\Delta,\qquad V=\Delta^{-1/2}U_B.
\]
Uniformly over such incident rows and coefficient vectors,
\begin{align}
 (R_Ew)_i&=x^\top\Pi_Bw_{-i}+O(1),\label{eq:loo-euclidean}\\
 (R_Dw)_i&=x^\top H_Bw_{-i}+O(1).\label{eq:loo-weighted}
\end{align}
The constants depend only on the displayed fixed bounds and $W$.
For the hollow degree-metric reconstruction on the model good event,
\eqref{eq:loo-weighted} also holds with $R^\circ$ in place of $R_D$.
The right-side coefficients use degrees recomputed from $B$. When $w$
is fixed, as in the oracle contrasts, or is measurable from $B$, those
coefficients are independent of the tested row. Uniform scalar stability
alone does not make a full-row-dependent $w$ independent.
\end{lemma}

\begin{proof}
The case $k=0$ has zero row action and can be given that convention.
For $k>0$, \eqref{eq:generalized-bounds} now yields
\begin{equation}\label{eq:loo-refined}
 \norm t\le\frac C{\sqrt{nd}},\qquad
 \norm{\Delta^{-1/2}Z_\perp}_{\op}
 \le\frac{C\sqrt k}{d^{3/2}\sqrt n}.
\end{equation}
For the Euclidean action right-multiply $G$ by
$Y^{-1}(V^\top V)^{-1/2}$. This does not change its column space and
gives a basis
\[
 \mathcal F=\binom{\alpha}{Q_B+E},\quad
 Q_B=V(V^\top V)^{-1/2},\quad
 \norm\alpha\le C/\sqrt n,\quad
 \norm E\le C\sqrt k/(d\sqrt n).
\]
Its Gram matrix is $I+J$, where
$\norm J\le C\{\sqrt k/(d\sqrt n)+1/n\}$.
The deleted Euclidean basis has row norm at most $C/\sqrt n$, so
\begin{align*}
 \norm{x^\top Q_B}&\le Ck/\sqrt n,&
 \norm{x^\top E}&\le Ck/(d\sqrt n),\\
 \norm{Q_B^\top w_{-i}}&\le W\sqrt n,&
 \norm{\mathcal F^\top w-Q_B^\top w_{-i}}
 &\le C_W(1/\sqrt n+\sqrt k/d).
\end{align*}
Expand
\[
 (R_Ew)_i=(x^\top Q_B+x^\top E)(I+J)^{-1}\mathcal F^\top w
\]
about $x^\top Q_BQ_B^\top w_{-i}$.
The errors are bounded by a constant times
\begin{equation}\label{eq:loo-errors}
 k/d+k/n+k\sqrt k/(d\sqrt n)=O(1),
\end{equation}
because $k\le bd$ and $d\le Cn$. This proves
\eqref{eq:loo-euclidean}. A coarse $O(n^{-1/2})$ bound for the lower
rows alone would not suffice; the stronger $E$ bound is essential.

For the degree-metric action, put $v_r=\Delta^{-1/2}z_r$,
$f_r=Vy_r+v_r$ and $w_V=V^\top\Delta w_{-i}$.
The exact top equation gives $k\mu_rt_r=x^\top f_r$, whence
\begin{align}\label{eq:weighted-row-expansion}
 (R_Dw)_i
 =\sum_r(x^\top Vy_r+x^\top v_r)
 \{kt_rw_i+y_r^\top w_V+v_r^\top\Delta w_{-i}
                              +f_r^\top E_xw_{-i}\}.
\end{align}
The normalization $G^\top DG=I$ gives
\[
 Y^\top Y+Z_\perp^\top Z_\perp+ktt^\top
 +(VY+\Delta^{-1/2}Z_\perp)^\top E_x
            (VY+\Delta^{-1/2}Z_\perp)=I.
\]
All terms except $Y^\top Y$ have norm at most $Ck/(nd)$, by
\eqref{eq:loo-refined} and the row bound on $V$.
Since $Y$ is square, $\norm{YY^\top-I}\le Ck/(nd)$ as well.
The leading term in \eqref{eq:weighted-row-expansion} is
$x^\top VYY^\top w_V$; its difference from
$x^\top VV^\top\Delta w_{-i}$ is at most $C_Wk^2/(nd)$.
The other factors satisfy
\begin{align*}
 \norm{x^\top V}&\le Ck/\sqrt{nd},&
 \norm{w_V}&\le C_W\sqrt{nd},\\
 \norm{(x^\top v_r)_r}&\le Ck/(d^{3/2}\sqrt n),&
 \norm{(v_r^\top\Delta w_{-i})_r}&\le C_W\sqrt{k/d},\\
 \norm{(f_r^\top E_xw_{-i})_r}&\le C_Wk/\sqrt{nd},&
 \norm{kt}&\le Ck/\sqrt{nd}.
\end{align*}
Their products are bounded by constants times
$k/d$, $k^2/(nd)$ and $k\sqrt k/(d\sqrt n)$, all bounded.
This proves \eqref{eq:loo-weighted} with negative as well as positive
informative eigenvalues.

On the model good event, full row incoherence implies
$|(H_D)_{ij}|\le C/n$, $|(R_D)_{ii}|\le Cd/n=Cp$ and
$|s_i|\le Cp$. The hollow adjustment to a bounded row contrast is
$s_i\{(H_Dw)_i-w_i\}=O(p)$, proving the last assertion.
\end{proof}

The result compares the \emph{actual} full-graph normalized algorithm to
a row-independent score, with its changing degree normalizers included.
It does not treat a full-graph eigenvector as independent of the row that
created it. This is the dependence step needed before any conditional
sharp-tail argument is applied.

\section{The earlier ordinary free-profile decoder}
\label{sec:sbm-proof}
\label{sbm:section}

We work under the ordinary-SBM assumptions of Theorem~\ref{thm:sbm}.
In particular, $K\ge2$ is fixed, $n_a\ge c_*n$, the block matrix is
symmetric and uniformly full rank at scale $p$, and
$p/2\le P_{ab}\le p\le1/4$. Write $d=np$, $C_a=\{i:z_i=a\}$,
$M_0=\Omega=ZPZ^\top$, and
\[
 \Pi_0=Z(Z^\top Z)^{-1}Z^\top,
 \qquad p_{a*}=(P_{a,z_j})_{j=1}^n.
\]
The diagonal of $M_0$ is retained in these comparison profiles;
the observed graph has no self-loops. The resulting bounded diagonal
correction is included below.

\subsection{A finite decoder using only the retained information}
\label{sbm:algorithm}

Let $D=\diag(A\one)$, and retain $K$ normalized adjacency eigenpairs
$D^{-1/2}AD^{-1/2}U=U\Lambda$, including the Perron direction and
selecting the remaining signal eigenvalues by absolute value. On graphs
where the encoder is defined, put
\[
 X=D^{-1/2}U,\qquad
 Q=X(X^\top X)^{-1/2},\qquad
 Y=DX\Lambda(X^\top X)^{-1/2},\qquad
 R=YQ^\top.
\]
These are exact algebraic operations. The generalized eigenrelation gives
$Y=AQ$ and hence $R=A\Pi$, where $\Pi=QQ^\top$. The Perron direction
implies $\Pi\one=\one$ and $R\one=A\one$. In particular, the entire
decoder below is evaluable from $(D,U,\Lambda)$ without accessing an
individual adjacency entry after compression. The exceptional-graph
convention in the theorem makes the encoder and decoder total.

Fix $0<c_0<1/16$, let $\bar d=n^{-1}\sum_iD_{ii}$, and set
\[
 \varepsilon=c_0\max(1,\bar d)/n,\qquad
 \mathcal I=[\varepsilon,1-\varepsilon].
\]
All clipping in this section is coordinatewise to $\mathcal I$.
For $v\in\R^n$ and $\mu\in\mathcal I^n$, define
\[
 \ell(v;\mu)
 =\sum_j\{v_j\log\mu_j+(1-v_j)\log(1-\mu_j)\},
 \qquad \ell_i(\mu)=\ell(R_{i*};\mu).
\]
The entries of $R$ may be signed. The score is nevertheless finite and
has the retained-factor expression
\begin{equation}
 \ell_i(\mu)=Y_iQ^\top\operatorname{logit}(\mu)
             +\sum_j\log(1-\mu_j).
 \label{sbm:score}
\end{equation}
For $s$ fitted profiles $\mathcal C=(\mu_1,\ldots,\mu_s)$ and simplex
weights $\pi$, use
\[
 \mathcal L(\mathcal C,\pi)
 =\sum_i\log\sum_{a=1}^s\pi_a e^{\ell_i(\mu_a)}.
\]
An E-step and its subsequent M-step are
\begin{align}
 \gamma_{ia}
 &=\frac{\pi_a e^{\ell_i(\mu_a)}}
         {\sum_b\pi_b e^{\ell_i(\mu_b)}},
 &N_a&=\sum_i\gamma_{ia},\notag\\
 \mu_{aj}^{\rm new}
 &=\operatorname{clip}_{\mathcal I}
       \left(\frac{\sum_i\gamma_{ia}R_{ij}}{N_a}\right),
 &\pi_a^{\rm new}&=N_a/n.
 \label{sbm:em}
\end{align}
When $N_a=0$, the old profile is retained and the new weight is zero.
The inner fitting subroutine performs at most $100$ exact E/M updates.
It may stop earlier when its labels are unchanged and its objective
increment is at most $10^{-12}(1+|\mathcal L_{\rm new}|)$. Ties in
all maximizations are settled deterministically. The proof uses only
finiteness and monotonicity of these inner updates.

The candidate attached to vertex $j$ is
$\nu_j=\operatorname{clip}_{\mathcal I}(R_{j*})$. Start from the clipped
global mean and fit one component. For $s=2,\ldots,K$, append a candidate
maximizing
\[
 \sum_i\left(\ell_i(\nu_j)-\max_{a<s}\ell_i(\mu_a)\right)_+,
\]
excluding the vertex identifiers selected in previous growing stages,
and fit all $s$ components from uniform weights. This gives $K$ profiles.
Next perform at most $T=\lceil\log n\rceil$ replacement rounds. For each
deletion $a$ from the round's incumbent, put
\[
 h_i^{(-a)}=\max_{b\ne a}\ell_i(\mu_b)
\]
and choose a candidate maximizing
$\sum_i(\ell_i(\nu_j)-h_i^{(-a)})_+$ over all vertices. Replace profile
$a$ by this candidate, initialize all $K$ weights uniformly, and run
the inner fit. After trying every deletion, retain the fit with largest
$\mathcal L$, including the unchanged incumbent among the choices.
A round without strict improvement terminates replacement.

Finally, evaluate the selected fit's responsibilities, perform one
mandatory M-step in \eqref{sbm:em}, and return
\begin{equation}
 \widehat z_i\in\argmax_a
 \{\ell_i(\mu_a^{\rm new})+\log\pi_a^{\rm new}\}.
 \label{sbm:decode}
\end{equation}
The reference to an E/M update here is ordered: the E-step is evaluated
at the selected fit, its M-step follows, and the decoder uses the new
profiles and weights.

\subsection{A common box and deterministic reconstruction bounds}
\label{sbm:box-section}

Throughout the deterministic arguments, intersect the base event
$\mathcal E_0$ in \eqref{eq:base-event} with the row-concentration
events of Lemma~\ref{lem:normalized-rows}. Its constants are fixed,
independent of $n$. In addition to the degree bounds and normalized
signal gap, its consequences used here are
\begin{equation}
 \norm{A-M_0}_{\op}\le C\sqrt d,\qquad
 \norm{U}_{2\to\infty}\le C/\sqrt n,\qquad
 \norm{\Pi-\Pi_0}_{\op}\le C/\sqrt d.
 \label{sbm:common-event}
\end{equation}
The same conclusions hold for the vertex-deleted graphs on their
corresponding good events. Their probability costs are handled in the
probability closure, after the sharp row calculation.

\begin{lemma}[Reconstruction and generated probability box]
\label{sbm:box}
On the full-graph good event, for all sufficiently large $n$,
\begin{align}
 \norm{R-M_0}_F^2&\le Cd,
 &\max_{i,j}|R_{ij}|&\le Cp,
 &\norm{R_i(I-\Pi_0)}_2&\le C\sqrt p,
 \label{sbm:reconstruction-bounds}\\
 \norm{\Pi_0(A-R)\Pi_0}_{\op}&\le C.
 \label{sbm:block-residual}
\end{align}
There are fixed $a_-,a_+,\eta_->0$ such that every profile generated by
the decoder, every clipped true-class mean, and every true comparison
profile belongs to
\begin{equation}
 \mathcal P_n
 =\{\mu:\ a_-p\le\mu_j\le a_+p,\quad
                    \mu_j\le1-\eta_-\text{ for every }j\}.
 \label{sbm:box-def}
\end{equation}
Every true probability is at distance at least $cp$ from both endpoints
of the actual clipping interval $\mathcal I$.
\end{lemma}

\begin{proof}
Write $E=A-M_0$. Since $R=A\Pi$,
\[
 R-M_0=E\Pi+M_0(\Pi-I).
\]
Each term has rank at most $K$, operator norm at most $C\sqrt d$, and
hence Frobenius norm at most $C\sqrt d$, by
\eqref{sbm:common-event} and $\norm{M_0}_{\op}\le Cd$.
For the other row bounds, write $G=X^\top X$. The degree bounds imply
that the eigenvalues of $G$ are between constant multiples of $d^{-1}$.
Thus $Q=XG^{-1/2}$ has row norms at most $C/\sqrt n$, and
\[
 AQ=DQT_0,\qquad T_0=G^{1/2}\Lambda G^{-1/2},\qquad
 \norm{T_0}_{\op}\le C.
\]
It follows that $\norm{(AQ)_i}_2\le Cd/\sqrt n$, which proves the
entrywise bound for $R=(AQ)Q^\top$. Moreover,
\[
 \norm{R_i(I-\Pi_0)}_2
 \le\norm{(AQ)_i}_2\norm{Q^\top(I-\Pi_0)}_{\op}
 \le C\sqrt{d/n}.
\]

Both $\Pi$ and $\Pi_0$ are orthogonal projectors. Consequently,
\[
 \norm{\Pi_0(I-\Pi)\Pi_0}_{\op}
 =\norm{(I-\Pi)\Pi_0}_{\op}^2\le C/d.
\]
Using $M_0=\Pi_0M_0\Pi_0$ gives
\[
 \Pi_0(A-R)\Pi_0
 =M_0\Pi_0(I-\Pi)\Pi_0+\Pi_0E(I-\Pi)\Pi_0.
\]
The two terms have norms at most $Cd/d$ and
$C\sqrt d/\sqrt d$, respectively. This proves
\eqref{sbm:block-residual} without a commutation assumption on $A$
and $\Pi$.

The average degree satisfies
\[
 \left|\bar d-n^{-1}\one^\top M_0\one\right|
 \le\norm{E}_{\op}\le C\sqrt d.
\]
Since $d/2\le n^{-1}\one^\top M_0\one\le d$, eventually
$d/4\le\bar d\le2d$. In particular $\bar d>1$ eventually, so
\begin{equation}
 c_0p/4\le\varepsilon\le2c_0p<p/8.
 \label{sbm:floor}
\end{equation}
Together with $P_{ab}\in[p/2,p]$ and $p\le1/4$, this gives the claimed
interior distance from the clipping endpoints.

Each active raw M-step coordinate is a convex combination of a column
of $R$, even though that column may have negative entries. It therefore
lies in $[-C_Rp,C_Rp]$. Clipping makes it at most
$\max(C_Rp,\varepsilon)\le C_\mu p$. The initial mean and all clipped
candidates obey the same bound; an inactive component retains an earlier
profile. These observations prove inductively that all generated
profiles lie between $c_0p/4$ and $C_\mu p$. If
$p\le1/(2C_\mu)$, their upper bound is $1/2$. Otherwise the lower
bound in \eqref{sbm:floor} gives
$\varepsilon\ge c_0/(8C_\mu)$, providing a fixed upper gap through
clipping. Enlarge the constants to include the true profiles and the
clipped true-class means.
\end{proof}

The box in \eqref{sbm:box-def} describes profiles generated by this
algorithm and its comparison profiles. Uniform statements below concern
this common box; they do not concern arbitrary vectors in the larger
nominal interval $\mathcal I^n$.

\begin{lemma}[Exact monotonicity for signed observations]
\label{sbm:em-monotonicity}
Every exact update \eqref{sbm:em} increases or preserves $\mathcal L$.
\end{lemma}

\begin{proof}
The log-sum-exp variational identity gives
\[
 \sum_{i,a}\gamma_{ia}
 \{\log\pi_a+\ell_i(\mu_a)-\log\gamma_{ia}\}
 \le\mathcal L(\mathcal C,\pi),
\]
with equality at the E-step. For one profile coordinate and fixed
responsibilities, put $a_j=\sum_i\gamma_{ia}R_{ij}$. Its objective is
\[
 a_j\log\mu+(N_a-a_j)\log(1-\mu),
\]
whose derivative has sign equal to $a_j/N_a-\mu$ when $N_a>0$.
Thus its global maximizer on $\mathcal I$ is the clipped weighted mean,
whether or not $a_j$ lies between zero and $N_a$. The weight maximizer
is $N_a/n$. Zero mass leaves the corresponding coordinate objective
constant. Maximization of this variational lower bound proves the claim.
\end{proof}

\subsection{The finite replacement certificate and common alignment}
\label{sbm:initialization-section}

For a $K$-profile list and weights, define the empirical reference loss
\[
 F_R(\mathcal C,\pi)
 =\sum_i\ell(R_i;p_{z_i,*})-\mathcal L(\mathcal C,\pi),
\]
and its population version
\[
 F_0(\mathcal C,\pi)
 =-\sum_i\log\sum_a\pi_a
                   e^{-\KL(p_{z_i,*}\Vert\mu_a)}.
\]
Here a vector Bernoulli divergence means the coordinate sum
\[
 \KL(p_{c*}\Vert\mu)=\sum_j
 \KL(\Ber(P_{c,z_j})\Vert\Ber(\mu_j)).
\]
The first term defining $F_R$ is a comparison constant independent of
the fitted list and weights. Hence every implemented likelihood
comparison is precisely the reverse comparison of $F_R$. The population
loss $F_0$ is nonnegative. Also put
\begin{align*}
 G_R(\mathcal C)
 &=\sum_i\{\ell(R_i;p_{z_i,*})-\max_a\ell_i(\mu_a)\},\\
 G_0(\mathcal C)
 &=\sum_c n_c\min_a\KL(p_{c*}\Vert\mu_a).
\end{align*}
These empirical reference losses are defined for signed $R$.

\begin{lemma}[Uniform comparison of signed and population objectives]
\label{sbm:objective-comparison}
On the full-graph good event, uniformly over center lists in
$\mathcal P_n$ and simplex weights,
\begin{equation}
 |F_R-F_0|+|G_R-G_0|\le Cn\sqrt d.
 \label{sbm:objective-comparison-bound}
\end{equation}
Always $G_R\le F_R$. With uniform weights,
$F_R\le G_R+n\log K$. Every true class supplies a candidate $\nu_j$
with $\norm{\nu_j-p_{c*}}_1\le C\sqrt d$; all these candidates are
available in every replacement trial.
\end{lemma}

\begin{proof}
For a true class $c$ and fitted profile $\mu_a$, the difference
$\ell(v;\mu_a)-\ell(v;p_{c*})$ is affine in $v$, with coefficient
$\operatorname{logit}(\mu_a)-\operatorname{logit}(p_{c*})$ bounded
uniformly on $\mathcal P_n$. The maximum and log-sum-exp operations are
1-Lipschitz for changes bounded in maximum norm. Therefore each
comparison in \eqref{sbm:objective-comparison-bound} is at most
$C\sum_i\norm{R_i-p_{z_i,*}}_1$. By
\eqref{sbm:reconstruction-bounds}, this sum is at most
$n\norm{R-M_0}_F\le Cn\sqrt d$. The common $\log p$ terms have
canceled in the score contrasts used in this argument.

The inequalities between hard and soft reference losses follow from
\[
 K^{-1}e^{\max_a s_a}\le K^{-1}\sum_a e^{s_a}
 \le e^{\max_a s_a}
\]
for uniform weights, and
$\sum_a\pi_a e^{s_a}\le e^{\max_a s_a}$ for arbitrary weights.
Finally, each class has at least $c_*n$ rows, so one of its rows has
squared error at most $Cd/n$ by the Frobenius bound. Its $\ell_1$
error is at most $C\sqrt d$. Clipping toward an interval containing
the true profile contracts coordinatewise absolute error.
\end{proof}

\begin{lemma}[Finite replacement certificate]
\label{sbm:replacement}
At the selected endpoint of the replacement stage,
\begin{equation}
 F_R\le \xi_n nd,\qquad
 \xi_n=C\{(1-1/K)^{\lceil\log n\rceil}+d^{-1/2}\}\longrightarrow0.
 \label{sbm:certificate}
\end{equation}
The same conclusion holds if replacement stops because a complete
round does not strictly improve the incumbent.
\end{lemma}

\begin{proof}
Assign each true profile a nearest fitted center in population
divergence. A true class contributes at least $G_0/K$. If a fitted
center is unused by this assignment, replacing that center by this
true profile removes its contribution without increasing any other
class's loss. If every center is used, the assignment between the
$K$ true classes and $K$ fitted centers is a bijection; replacing the
center of the highest-loss class has the same property. Consequently
some ideal replacement has population hard loss at most
$(1-1/K)G_0$.

The derivative of $\KL(p_{c*}\Vert\mu)$ in a parameter coordinate is
$(\mu_j-P_{c,z_j})/[\mu_j(1-\mu_j)]$, uniformly bounded on the common
box. Replacing the ideal profile by the available candidate from
Lemma~\ref{sbm:objective-comparison} costs at most $Cn\sqrt d$.
Combining the objective comparisons gives a candidate list $\mathcal C'$
with
\[
 G_R(\mathcal C')\le(1-1/K)F_R(\mathcal C,\pi)+Cn\sqrt d.
\]

For a fixed deletion $a$, the empirical hard loss after inserting
$\nu_j$ is exactly
\[
 \sum_i\{\ell(R_i;p_{z_i,*})-h_i^{(-a)}
                         -(\ell_i(\nu_j)-h_i^{(-a)})_+\}.
\]
Thus the implemented maximum positive gain minimizes this loss among
candidates. All deletions are tried. Uniform trial weights add at most
$n\log K$ to the hard loss, and the subsequent exact inner fit cannot
increase $F_R$ by Lemma~\ref{sbm:em-monotonicity}. Including the
incumbent among the choices proves the recurrence
\begin{equation}
 F_{R,t+1}\le(1-1/K)F_{R,t}+Cn\sqrt d+n\log K,
 \qquad F_{R,t+1}\le F_{R,t}.
 \label{sbm:replacement-recurrence}
\end{equation}

Every generated profile lies in the common box, so each population
divergence is at most $Cd$. Hence $F_{0,0}\le Cnd$ for arbitrary
simplex weights, and the same bound holds for $F_{R,0}$ by
\eqref{sbm:objective-comparison-bound}. No preliminary accuracy guarantee
for the growing stage is needed: its generated list supplies this finite
starting bound. Iterating
\eqref{sbm:replacement-recurrence} for $\lceil\log n\rceil$ rounds
gives \eqref{sbm:certificate}, after enlarging its constant. If a
complete round gives no strict improvement, its chosen loss is the
incumbent's loss. The first inequality in
\eqref{sbm:replacement-recurrence} then immediately yields
$F_R\le K(Cn\sqrt d+n\log K)$, which has the same conclusion.
\end{proof}

\begin{lemma}[Common matching and refined responsibilities]
\label{sbm:responsibilities}
The selected fit admits a single permutation after which
\begin{equation}
 \max_a\norm{\mu_a-p_{a*}}_1=o(d),\qquad
 -\log\min_a\pi_a=o(d).
 \label{sbm:consistent-fit}
\end{equation}
Its next E-step satisfies
\begin{equation}
 \frac1n\sum_{i,a}|\gamma_{ia}-Z_{ia}|\le C/d.
 \label{sbm:responsibility-radius}
\end{equation}
\end{lemma}

\begin{proof}
The preceding lemmas give $F_0\le\zeta_nnd$ with $\zeta_n\to0$.
Since $F_0\ge G_0$, each true class has a fitted center whose
divergence is $o(d)$. On the common box, Bernoulli entropy curvature
and Cauchy--Schwarz imply
\[
 \KL(p_{a*}\Vert\mu)\ge
 \frac{\norm{p_{a*}-\mu}_2^2}{Cp}
 \ge\frac{\norm{p_{a*}-\mu}_1^2}{Cd}.
\]
Thus each true profile has a fitted center at $o(d)$ in $\ell_1$.
Uniform full rank and positive class sizes give
\[
 \min_{a\ne b}\norm{p_{a*}-p_{b*}}_2^2\ge cnp^2,
 \qquad
 \min_{a\ne b}\norm{p_{a*}-p_{b*}}_1\ge cd.
\]
The nearby centers must therefore be distinct. There are exactly $K$
of them, proving one common matching and the first assertion of
\eqref{sbm:consistent-fit}.

For a row of true class $a$, each wrong-center population divergence
is at least $c_1d$ eventually, by this matching and the common-box
Lipschitz bound. Its population soft loss is at least
$-\log(\pi_a+e^{-c_1d})$. It is the same for all $n_a$ rows of that
class, and every population soft loss is nonnegative. Consequently
\[
 \pi_a+e^{-c_1d}\ge e^{-\zeta_nnd/n_a},\qquad
 \pi_a\ge\tfrac12e^{-C\zeta_nd}
\]
for all sufficiently large $n$. This proves the weight assertion.

Fix a sufficiently small constant $t_0>0$. If
$\norm{R_i-p_{a*}}_1\le t_0d$ and $z_i=a$, the population separation,
\eqref{sbm:consistent-fit}, and bounded contrast coefficients give
a fitted correct-class margin at least $c_2d$. The wrong responsibility
mass on this row is at most $(K-1)e^{-c_2d}$. Every remaining row has
squared error at least $t_0^2d^2/n$. The Frobenius bound limits their
number to $Cn/d$. Summing their responsibility masses, which are at
most one, proves \eqref{sbm:responsibility-radius}.
\end{proof}

\subsection{Fitted profiles change the relevant score by a bounded amount}
\label{sbm:fitted-section}

Define the centered block fluctuation
\begin{equation}
 \mathcal B
 =\max_{a,c}
 \left|\frac{\one_{C_a}^\top(A-\E A)\one_{C_c}}
                   {\sqrt{n_an_c}}\right|.
 \label{sbm:block-fluctuation}
\end{equation}

\begin{lemma}[Block moments and deletion]
\label{sbm:block-moments}
For every fixed $t\ge0$, $\E e^{t\mathcal B}\le C_t$. For a tested
vertex $i$, let $\mathcal B_i$ be the analogous fluctuation calculated
from $B=A_{-i,-i}$ and the deleted class sizes. Then
$\E e^{t\mathcal B_i}\le C_t$, and on $D_{ii}\le Cd$,
\begin{equation}
 \mathcal B\le C(1+\mathcal B_i).
 \label{sbm:fluctuation-deletion}
\end{equation}
\end{lemma}

\begin{proof}
Each normalized block sum in \eqref{sbm:block-fluctuation} is a sum
of independent centered upper-triangular Bernoulli entries with
coefficients of magnitude at most $C/n$. Its sum of coefficient-squared
edge probabilities is at most $Cp$. The elementary centered Bernoulli
transform bound therefore gives, for each fixed real $t$,
\[
 \log\E\exp\left\{
 t\frac{\one_{C_a}^\top(A-\E A)\one_{C_c}}
         {\sqrt{n_an_c}}\right\}
 \le C t^2 e^{C|t|/n}p.
\]
For example, this follows by applying
$e^u-1-u\le u^2e^{|u|}/2$ to every independent edge and multiplying
their transforms. Taking both signs and summing over the fixed number
of block pairs bounds $\E e^{t\mathcal B}$.

Deleting one vertex preserves the positive class-size bounds, so the
same argument applies to $\mathcal B_i$. The unnormalized difference
between a full and deleted block sum consists of at most twice the
tested row's centered edge sum. Its absolute value is at most
$C(D_{ii}+d)$. Dividing by $\sqrt{n_an_c}\ge cn$ gives $O(p)$.
The factor converting a deleted normalized block sum to the full
denominator is bounded by a constant. This proves
\eqref{sbm:fluctuation-deletion}.
\end{proof}

For each class define the raw true-class mean and its clipped version
\[
 \bar r_a=n_a^{-1}\sum_{l\in C_a}R_l,
 \qquad \bar\mu_a=\operatorname{clip}_{\mathcal I}(\bar r_a).
\]
These are comparison quantities only. For $a\ne b$, also define
\begin{equation}
 w_j^{ab}=\log\frac{P_{a,z_j}(1-P_{b,z_j})}
                         {P_{b,z_j}(1-P_{a,z_j})},
 \qquad
 r_j^{ab}=\log\frac{1-P_{a,z_j}}{1-P_{b,z_j}}.
 \label{sbm:true-weights}
\end{equation}

\begin{lemma}[Uniform fitted-score comparison]
\label{sbm:fitted-score}
On the full-graph good event, every M-step from responsibilities
satisfying \eqref{sbm:responsibility-radius}, after their common
matching, has weights bounded below and obeys
\begin{equation}
 \left|
 \ell_i(\mu_a)-\ell_i(\mu_b)+\log\frac{\pi_a}{\pi_b}
 -\left\{R_iw^{ab}+\sum_jr_j^{ab}\right\}
 \right|
 \le C(1+\mathcal B)
 \label{sbm:fitted-score-bound}
\end{equation}
uniformly over $i,a,b$ and all such responsibility matrices.
\end{lemma}

\begin{proof}
Write $e=n^{-1}\sum_{l,a}|\gamma_{la}-Z_{la}|\le C/d$.
Then $N_a=n_a+O(ne)$, so $N_a\ge cn$ eventually and the updated
weights $N_a/n$ are bounded below. By
\eqref{sbm:reconstruction-bounds}, $\max_l\norm{R_l}_1\le Cd$.
The difference of the raw responsibility-weighted mean and $\bar r_a$
can be written as a numerator perturbation and a denominator
perturbation. Their combined $\ell_1$ norm is at most
\[
 \frac{Cd}{N_a}\sum_l|\gamma_{la}-Z_{la}|
 +\frac{Cd}{N_a}|N_a-n_a|
 \le Ced\le C.
\]
Clipping is 1-Lipschitz, whence
\begin{equation}
 \norm{\mu_a-\bar\mu_a}_1\le C.
 \label{sbm:actual-perfect-mean}
\end{equation}

Jensen's inequality and the Frobenius bound give
\[
 \norm{\bar r_a-p_{a*}}_2^2
 \le n_a^{-1}\sum_{l\in C_a}\norm{R_l-p_{a*}}_2^2\le Cp.
\]
Clipping toward an interval containing $p_{a*}$ preserves this bound.
The interior distance from Lemma~\ref{sbm:box} also gives the scalar
inequality
\[
 |\operatorname{clip}_{\mathcal I}(x)-x|
 \le\frac{(x-u)^2}{cp}
\]
for each true coordinate $u$. Thus
$\norm{\bar\mu_a-\bar r_a}_1\le C$.

For a true block $C_c$, decompose the block sum of
$\bar r_a-p_{a*}$ using $R-A$, $A-\E A$, and $\E A-M_0$.
The first contributes at most
\[
 \sqrt{n_c/n_a}\,
 \norm{\Pi_0(R-A)\Pi_0}_{\op}\le C,
\]
by \eqref{sbm:block-residual}; the second is at most $C\mathcal B$.
The diagonal correction contributes at most $P_{aa}\le p$.
Including the clipping correction proves, for
$\delta_a=\bar\mu_a-p_{a*}$,
\begin{equation}
 \norm{\delta_a}_2^2\le Cp,
 \qquad
 \left|\sum_{j\in C_c}\delta_{aj}\right|
 \le C(1+\mathcal B).
 \label{sbm:mean-errors}
\end{equation}

For any profile in the common box, the score derivatives are
\[
 \frac{\partial\ell_i}{\partial\mu_j}
 =\frac{R_{ij}-\mu_j}{\mu_j(1-\mu_j)},\qquad
 \frac{\partial^2\ell_i}{\partial\mu_j^2}
 =-\frac{R_{ij}}{\mu_j^2}
  -\frac{1-R_{ij}}{(1-\mu_j)^2}.
\]
Their absolute values are at most $C$ and $C/p$, respectively.
Therefore \eqref{sbm:actual-perfect-mean} changes the score by $O(1)$,
and the second-order remainder in a Taylor expansion at $p_{a*}$ is
at most $C\norm{\delta_a}_2^2/p\le C$.

Split the first-order gradient at $p_{a*}$ into its projection onto
$\col(Z)$ and its orthogonal complement. Its block-constant
projection has bounded coordinates; paired with $\delta_a$, it is
bounded by \eqref{sbm:mean-errors}. The diagonal multiplier
$[P_{a,z_j}(1-P_{a,z_j})]^{-1}$ is block constant and commutes with
$\Pi_0$. Hence the complementary gradient has norm at most
\[
 \frac C p\norm{R_i(I-\Pi_0)}_2\le C/\sqrt p.
\]
Its inner product with $\delta_a$, whose norm is at most $C\sqrt p$,
is $O(1)$. We have proved
$|\ell_i(\mu_a)-\ell_i(p_{a*})|\le C(1+\mathcal B)$.
Apply the same calculation to $b$ and include the bounded updated
log-weight ratio. The exact true-profile contrast is the expression
in braces in \eqref{sbm:fitted-score-bound}, proving the result.
\end{proof}

\subsection{Euclidean projected evidence at the oracle tilt}
\label{sbm:tilted-section}

Fix a tested vertex $i$ and labels $a\ne b$. Let $B=A_{-i,-i}$,
$x=A_{-i,i}$, and let $\Pi_B$ be the Euclidean projector onto the
deleted generalized signal space from Lemma~\ref{lem:loo}. Write
$\Pi_{0,B}$ for the true block projector on the remaining vertices.
For this paragraph, $w$ and $r$ denote the restrictions of
\eqref{sbm:true-weights} to these vertices. The exact oracle log
likelihood ratio is
\[
 L_{i,ab}=b_{i,ab}+x^\top w,
 \qquad b_{i,ab}=\sum_{j\ne i}r_j^{ab}.
\]
Let $I_{i,ab}$ be its exact Bernoulli Chernoff information in
\eqref{eq:sbm-information} and $t_*$
its optimizing parameter, using the transform under class $a$ with
factor $e^{-t_*L_{i,ab}}$. Under the corresponding optimizing tilt,
the row coordinates are independent with block-constant probabilities
\begin{equation}
 q_j=\frac{P_{a,z_j}^{1-t_*}P_{b,z_j}^{t_*}}
 {P_{a,z_j}^{1-t_*}P_{b,z_j}^{t_*}
 +(1-P_{a,z_j})^{1-t_*}(1-P_{b,z_j})^{t_*}}.
 \label{sbm:tilted-probability}
\end{equation}
These probabilities are comparable to $p$. The log-odds are bounded,
and uniform full rank gives $\sum_{j\ne i}w_j^2\asymp n$.
For all tilt parameters in $[0,1]$, the log transform therefore has
second derivative between $cd$ and $Cd$. Its two endpoint values
are zero. It follows that $I_{i,ab}\asymp d$, the optimizing parameter
stays inside a fixed compact subinterval of $(0,1)$, and its tilted
variance is comparable to $d$. Thus the oracle conditions of
Lemma~\ref{lem:sharp} hold uniformly.

\begin{lemma}[Euclidean tilted cancellation]
\label{sbm:tilted-cancellation}
On the deleted-graph good event, let $\beta=(\Pi_B-I)w$. Then
\begin{equation}
 \max_j|\beta_j|\le C,
 \qquad |q^\top\beta|\le C,
 \qquad \sum_jq_j\beta_j^2\le C.
 \label{sbm:tilted-errors}
\end{equation}
Consequently, conditional on $B$, for every real $u$,
\begin{equation}
 \PP_a\{b_{i,ab}+x^\top\Pi_Bw\le u\mid B\}
 \le C\frac{e^{-I_{i,ab}+t_*u}}{\sqrt d}.
 \label{sbm:conditional-tail}
\end{equation}
\end{lemma}

\begin{proof}
The deleted Euclidean basis has row norms at most $C/\sqrt n$, so
bounded entries of $w$ imply bounded entries of $\Pi_Bw$.
Let $\delta=\norm{(I-\Pi_B)\Pi_{0,B}}_{\op}\le C/\sqrt d$,
as follows from \eqref{eq:euclidean-angle} for the deleted graph.
Both $q$ and $w$ belong to $\col(\Pi_{0,B})$. Orthogonality gives
\[
 q^\top\beta
 =-\ip{(I-\Pi_B)q}{(I-\Pi_B)w}.
\]
Since $\norm q_2\le Cp\sqrt n$ and $\norm w_2\le C\sqrt n$,
its absolute value is at most $Cnp\delta^2\le C$.
Also $\norm\beta_2^2\le\delta^2\norm w_2^2\le Cn/d$, so
$\sum_jq_j\beta_j^2\le Cp\norm\beta_2^2\le C$.
This proves \eqref{sbm:tilted-errors}.

The conditional transform identity is
\[
 \E_a e^{-t_*(b_{i,ab}+x^\top\Pi_Bw)}
 =e^{-I_{i,ab}}\E_{q}e^{-t_*x^\top\beta}.
\]
Independence, $\log(1+y)\le y$, and
$e^v-1-v\le v^2e^{|v|}/2$ bound the final log transform by
$C|q^\top\beta|+C\sum_jq_j\beta_j^2\le C$.
Under tilting by this estimated score, the row is still a product of
Bernoulli variables. Its probabilities and complements are comparable
to their oracle-tilted counterparts because the coefficient errors
are bounded. Its score variance is of order $d$: in the oracle-tilted
quadratic form,
\[
 \sum_jq_j(1-q_j)(w_j+\beta_j)^2
 \ge\tfrac12\sum_jq_j(1-q_j)w_j^2
       -\sum_jq_j\beta_j^2\ge cd
\]
eventually, and the corresponding upper bound is $Cd$.
The bounded summands have total third absolute centered moment $O(d)$.
The bounds \eqref{sbm:tilted-errors} give all three premises in
\eqref{eq:tilt-error}, with zero additive intercept error and fixed
constants. Lemma~\ref{lem:sharp} now gives
\eqref{sbm:conditional-tail}.
A threshold $u$ only contributes the displayed factor $e^{t_*u}$;
it does not change this variance or the interval-concentration bound.
\end{proof}

\begin{proposition}[Sharp error bound for the finite fitted decoder]
\label{sbm:good-event-risk}
Align the selected fit by the common matching from
Lemma~\ref{sbm:responsibilities}. Let $\mathcal E_i$ be the intersection
of the full-graph and vertex-deleted good events, including the tested
degree bound, used in Lemmas~\ref{lem:normalized-rows} and~\ref{lem:loo}.
For a vertex with true class $a$ and any $b\ne a$, the final decoder
\eqref{sbm:decode} satisfies
\begin{equation}
 \PP_z\{\widehat z_i\text{ is aligned to }b,\ \mathcal E_i\}
 \le C\frac{e^{-I_{i,ab}}}{\sqrt d}.
 \label{sbm:good-event-bound}
\end{equation}
In particular, summing over competitors and averaging over vertices
bounds the good-event expected misclassification loss by
\begin{equation}
 \frac Cn\sum_i\sum_{b\ne z_i}
              \frac{e^{-I_{i,z_i b}}}{\sqrt d}.
 \label{sbm:good-average}
\end{equation}
\end{proposition}

\begin{proof}
The finite replacement certificate and its mandatory final E/M update
give the premise of Lemma~\ref{sbm:fitted-score}. A wrong assignment
from $a$ to $b$, including an unfavorable score tie, implies on the
full-graph good event that
\[
 R_iw^{ab}+\sum_jr_j^{ab}\le C(1+\mathcal B).
\]
The full-to-deleted Euclidean identity \eqref{eq:loo-euclidean} applies
to the fixed bounded vector $w^{ab}$. Together with
\eqref{sbm:fluctuation-deletion}, it gives the deterministic event
containment
\begin{equation}
 \{\widehat z_i\text{ is aligned to }b,\ \mathcal E_i\}
 \subseteq
 \{b_{i,ab}+x^\top\Pi_Bw_{-i}^{ab}
                \le C(1+\mathcal B_i),\ B\text{ good}\}.
 \label{sbm:event-containment}
\end{equation}
The nonedge intercept changes by $r_i^{ab}=O(p)$ when the tested
diagonal is removed, which is absorbed into the constant.

Use the full-graph event only for this deterministic containment.
The right-hand event does not condition the tested row on full-graph
eigenvectors, fitted parameters, or a degree event. Conditional on
$B$, the incident row has its original independent Bernoulli law.
Applying \eqref{sbm:conditional-tail} with
$u=C(1+\mathcal B_i)$, integrating over good $B$, and using
Lemma~\ref{sbm:block-moments} gives
\[
 \PP_z\{\widehat z_i\text{ is aligned to }b,\ \mathcal E_i\}
 \le C\frac{e^{-I_{i,ab}}}{\sqrt d}
       \E e^{C\mathcal B_i}
 \le C'\frac{e^{-I_{i,ab}}}{\sqrt d}.
\]
This proves \eqref{sbm:good-event-bound}. The common permutation is
the same for every vertex, so the loss minimized over permutations
is at most the loss under this matching. Summing the pairwise bounds
proves \eqref{sbm:good-average}.
\end{proof}

The sharp Bernoulli oracle comparison in Lemma~\ref{lem:sharp}
identifies the terms in \eqref{sbm:good-average} with local oracle
errors up to fixed constants, including the $d^{-1/2}$ factor
\cite{zhouli}. The probability closure pays the complement of these
good events at that same oracle scale and completes
Theorem~\ref{thm:sbm}.

\paragraph{Extension to general positive block profiles.}
The deterministic and conditional arguments also apply to a fixed
positive full-rank block class
\[
 c_-p\le P_{ab}\le c_+p,\qquad
 \max_{a,b}P_{ab}\le1-\epsilon_0,
\]
with positive comparable class sizes and $d\to\infty$; $p$ need not tend
to zero. In particular $p\le p_{\max}:=(1-\epsilon_0)/c_-$.
The average degree is eventually between $c_-d/2$ and $2c_+d$.
Choose
\[
 0<c_0<\min\{1/16,\ c_-/(4c_+),\ \epsilon_0/2\}.
\]
Then $\varepsilon\le2c_0c_+p<c_-p/2$, giving a distance at least
$c_-p/2$ between every true probability and the lower clipping endpoint.
For $n\ge2$ the elementary graph bound $\bar d\le n-1$ also gives
$\varepsilon\le c_0<\epsilon_0/2$. Hence the distance to the upper
endpoint is at least $\epsilon_0/2\ge
\{\epsilon_0/(2p_{\max})\}p$. This proves the required interior
condition for dense as well as sparse profiles.

The generated common-box proof, all likelihood derivatives, and all
block-fluctuation moments now hold with fixed constants depending on
$c_-,c_+,\epsilon_0$. The exact oracle-tilted probabilities remain
between the two candidate probabilities. Both these means and the exact
log-odds vectors are block constant, so the Euclidean tilted calculation
uses no sparse approximation. All arguments above therefore apply on
the proved good event to this general class. Proposition~\ref{prop:sbm-margin}
pays that event's complement under an explicit information--degree
margin; Corollary~\ref{cor:known} supplies the sparse equal-exposure
case where the margin is automatic.

\section{The earlier unregularized DC profile decoder}
\label{sec:dcsbm-proof}

This section proves the degree-corrected result in \cref{thm:dc}.
The reconstruction used here is the projection in the observed-degree
inner product.  It is an explicit decoder specification: the preceding
Euclidean reconstruction is not silently identified with it.
All likelihoods below are Bernoulli likelihoods.  The condition
\(np^2\le C_1\) controls their approximation at a bounded \emph{log-likelihood}
scale.

Write \(Z_{ia}=\ind{z_i=a}\), \(\Theta=\diag(\theta_1,\ldots,\theta_n)\),
\(B=pC\), and
\[
 \Omega=\Theta ZBZ^T\Theta,\qquad
 D_0=\diag(\Omega\one),\qquad
 D=\diag(D_1,\ldots,D_n),\quad D_i=\sum_jA_{ij}.
\]
The matrix \(\Omega\) includes its diagonal; the observed graph is hollow.
Thus \(\|\Omega-\E A\|_{\op}=O(p)\).  In particular, every raw noise
bound used below is unchanged up to its fixed constant by this completion.
Throughout this section a generic constant may depend on the fixed
parameter bounds in \cref{thm:dc}, but never on \(n\).  Constants depending
also on a fixed initialization-basin radius \(M\) are denoted \(C(M)\).

We use the common good event \(\mathcal G\) supplied by
\cref{lem:normalized-rows} and the probability calculation in
\cref{sec:probability}.  Its consequences needed here are
\begin{align}
 c d\le D_i\le C d,\qquad
 \|A-\Omega\|_{\op}\le C\sqrt d,\qquad
 \max_i\|U_{i\cdot}\|_2\le Cn^{-1/2},                    \label{dcsbm:good}\\
 \max_{a,b}\abs{[Z^T(A-\Omega)Z]_{ab}}\le Cn,           \label{dcsbm:block-event}
\end{align}
together with the corresponding normalized spectral and degree
consequences for the genuinely deleted graphs.  The retained eigenvectors
\(U\) are orthonormal, and their eigenvalues \(\mu_1,\ldots,\mu_K\)
are the \(K\) largest in absolute value.  The degree eigenvector is retained.
The common lemmas establish these consequences from the model and pay
their failure probabilities; they are not extra statistical assumptions.
We now prove all remaining deterministic and local probabilistic steps.

\subsection{The degree-metric reconstruction and its hollow correction}

Put \(S=D^{-1/2}AD^{-1/2}\), \(\Pi=UU^T\), and
\begin{equation}
 H=D^{-1/2}\Pi D^{1/2},\qquad
 R=AH=D^{1/2}U\diag(\mu_1,\ldots,\mu_K)U^TD^{1/2}.
 \label{dcsbm:canonical}
\end{equation}
The second formula evaluates \(R\) from retained spectral information.
Although \(H\) need not be symmetric, \(H^2=H\) and \(DH=H^TD\):
it is an orthogonal projection for
\(\langle x,y\rangle_D=x^TDy\).  Since the degree eigenvector
is proportional to \(D^{1/2}\one\), \(H\one=\one\) and
\(R\one=D\one\).

For \(h_i=R_{ii}\), define
\begin{equation}
 s_i=\frac{h_i}{1-\Pi_{ii}},\qquad
 R^\circ=R+\diag(s_1,\ldots,s_n)(H-I).
 \label{dcsbm:hollow}
\end{equation}
On \(\mathcal G\), the denominators are positive for all sufficiently
large \(n\).  Direct substitution gives
\begin{equation}
 R^\circ_{ii}=0,\qquad R^\circ\one=D\one,\qquad
 R^\circ Q=AQ\quad
 \text{if }\col(Q)\subseteq\col(D^{-1/2}U).
 \label{dcsbm:exact-identities}
\end{equation}
The last identity follows because \(HQ=Q\).  The hollow correction need
not be symmetric; a symmetric block parameter is fitted by the ordered
Bernoulli criterion, equivalently by symmetrizing its pair contributions.
The row profile scores always use the original rows of \(R^\circ\).

\begin{lemma}[Entry and absolute-row bounds]
\label{dcsbm:entry-lemma}
On \(\mathcal G\),
\begin{equation}
 \max_{i,j}|R^\circ_{ij}|\le Cp,\qquad
 \max_i\sum_j|R^\circ_{ij}|\le Cd,\qquad
 \max_i|s_i|\le Cp .
 \label{dcsbm:entry-bounds}
\end{equation}
Every action in \cref{dcsbm:canonical,dcsbm:hollow} is computable from
\((D,U,\mu)\), with no access to individual edges after encoding.
\end{lemma}

\begin{proof}
The normalized adjacency has operator norm at most one.  Consequently
\[
 |R_{ij}|
 \le \sqrt{D_iD_j}\,\|U_{i\cdot}\|_2\|U_{j\cdot}\|_2
 \le Cd/n=Cp.
\]
Moreover
\(
 |H_{ij}|
 \le\sqrt{D_j/D_i}\,\|U_{i\cdot}\|_2\|U_{j\cdot}\|_2
 \le C/n
\).
These estimates give \(|s_i|\le Cp\) and the entry bound for \(R^\circ\);
summing its entries in absolute value gives the row bound.
The displayed spectral formulas also prove the computational statement.
\end{proof}

\subsection{A finite deterministic spectral initializer}

The row-normalized population centers are particularly simple even when
the \(\theta_i\)'s vary within blocks.  To see this without replacing the
observed degrees by their expectations, put
\[
 M_D=D^{-1/2}\Omega D^{-1/2},\qquad
 G_D=Z^T\Theta D^{-1}\Theta Z.
\]
The diagonal entries of \(G_D\) are between constant multiples of \(n/d\).
The \(K\) nonzero eigenvalues of \(M_D\) have absolute values bounded
away from zero, by the lower singular-value bound on \(C\).
An orthonormal basis of its range has the form
\begin{equation}
 V=D^{-1/2}\Theta ZG_D^{-1/2}O
 \label{dcsbm:population-basis}
\end{equation}
for an orthogonal \(K\)-by-\(K\) matrix \(O\).  Each row has norm
comparable to \(n^{-1/2}\), and its normalization is exactly the
\(z_i\)-th row of \(O\).  These \(K\) unit centers have pairwise distance
\(\sqrt2\).

By \cref{dcsbm:good},
\(\|S-M_D\|_{\op}\le C/\sqrt d\).  The signal/bulk gap and
Davis--Kahan therefore give an aligned basis with
\(\|U-V\|_F^2\le C/d\).  The elementary normalization inequality
\[
 \left\|\frac{u}{\|u\|_2}-\frac{v}{\|v\|_2}\right\|_2
 \le \frac{2\|u-v\|_2}{\|v\|_2}
\]
holds when \(u\ne0\), and holds with a harmless change of constant
under any fixed unit-vector convention for \(u=0\).  For the normalized
rows \(x_i\) and true unit centers \(c_a\), it follows that
\begin{equation}
 \sum_{i=1}^n\|x_i-c_{z_i}\|_2^2\le Cn/d.
 \label{dcsbm:embedding-sse}
\end{equation}

The following initializer requires no numerical \(K\)-means certificate.
Starting with all rows remaining, repeat \(K\) times: choose a remaining
row whose closed ball of radius \(r=1/8\) contains the largest number
of remaining rows, break ties in a fixed order, record that row, and
remove all remaining rows at distance at most \(2r=1/4\) from it.
Finally assign every original row to its nearest recorded representative.
If a step is undefined, use the algorithm's fixed fallback partition.
The procedure uses \(O(Kn^2)\) distance comparisons in fixed dimension.

\begin{lemma}[Ball-count initialization]
\label{dcsbm:initializer}
Suppose every true block contains at least \(\gamma n\) rows and
the unit centers are separated by \(\sqrt2\).  Put
\(\epsilon=1/32\).  If
\(\sum_i\|x_i-c_{z_i}\|_2^2\le E\) and
\(m=E/\epsilon^2<\gamma n/3\), the procedure above selects exactly
one representative for each true center and, after one common
permutation of labels, makes at most \(m\) errors.  In particular,
\cref{dcsbm:embedding-sse} implies an initial error fraction at most
\(C/d\) on \(\mathcal G\).
\end{lemma}

\begin{proof}
Call a row good if it is within \(\epsilon\) of its true center.
There are at most \(m\) bad rows.  Inductively, all good rows of each
previously selected block have been removed and all good rows of each
unselected block remain.  A radius-\(r\) ball centered at a good row
of any unselected block contains every good row of that block, since
\(2\epsilon<r\).  Its count is at least \(\gamma n-m>m\).
Hence the maximizing ball contains a good row from an unselected block.
Its selected representative is within \(r+\epsilon=5/32\) of that
block's center.  Every good row of this block is within
\(r+2\epsilon=3/16<2r\) of the representative, while every good row
of another block is at least
\(\sqrt2-r-2\epsilon>2r\) away.  This proves the induction.
After \(K\) selections, a good row is at most \(3/16\) from its own
representative and at least \(\sqrt2-3/16\) from every other one.
Thus final nearest-representative assignment classifies all good
rows correctly.
\end{proof}

\subsection{Block-total accuracy and estimation of the nuisance parameters}

We first record the residual estimate behind both the block fit and the
tilted cancellation.

\begin{lemma}[Block residual in the degree metric]
\label{dcsbm:block-residual}
On \(\mathcal G\), with \(F=(I-\Pi)D^{1/2}Z\),
\begin{equation}
 \|F\|_{\op}\le\|F\|_F\le C\sqrt n,\qquad
 \sum_i\frac{(D_i-D_{0,i})^2}{D_i}\le Cn.
 \label{dcsbm:residual-estimates}
\end{equation}
The same conclusions hold for any genuinely deleted graph under its
corresponding common good event.
\end{lemma}

\begin{proof}
The degree-vector bound follows from
\[
 \|(D-D_0)\one\|_2
 =\|(A-\Omega)\one\|_2\le C\sqrt{nd}
\]
and \(\min_iD_i\ge cd\).
Write
\[
 D^{1/2}Z
 =D^{-1/2}D_0Z+D^{-1/2}(D-D_0)Z.
\]
Since \(D_{0,i}=p\theta_i\sum_c n_cC_{z_i c}\), the first term
belongs to \(\col(D^{-1/2}\Theta Z)=\col(M_D)\).
Its Frobenius norm is \(O(\sqrt{nd})\), while the angle between this
space and \(\col(U)\) is \(O(d^{-1/2})\).  Its residual is therefore
\(O(\sqrt n)\).  The squared Frobenius norm of the second term is
exactly the degree sum in \cref{dcsbm:residual-estimates}, because
every row of \(Z\) contains one unit entry.  This proves the assertion.
Deleting one vertex preserves all fixed group-mass bounds and the same
calculation applies to the minor.
\end{proof}

\begin{lemma}[Canonical reconstructed block totals]
\label{dcsbm:block-totals}
Fix \(M<\infty\).  Uniformly over every possibly data-dependent
partition \(\widehat Z\) with at most \(Mn/d\) mistakes after a common
permutation,
\begin{equation}
 [\widehat Z^TR^\circ\widehat Z]_{ab}
 =n_an_bpC_{ab}+O_M(n).
 \label{dcsbm:totals-accuracy}
\end{equation}
The usual diagonal exclusion is included in the \(O_M(n)\) term.
\end{lemma}

\begin{proof}
Let \(S_{\rm tail}=(I-\Pi)S(I-\Pi)\).  Spectral invariance gives the
exact quadratic identity
\begin{equation}
 Z^T(R-A)Z=-F^TS_{\rm tail}F .
 \label{dcsbm:quadratic-identity}
\end{equation}
Because \(M_D\) has rank \(K\), Weyl's inequality gives
\(\|S_{\rm tail}\|_{\op}\le C/\sqrt d\).
By \cref{dcsbm:block-residual}, the right side of
\cref{dcsbm:quadratic-identity} has norm at most \(Cn/\sqrt d\).
The hollow correction changes a true block total by at most \(Cd\):
indeed its total absolute contribution is at most
\[
 \max_i|s_i|\sum_i\sum_j|(H-I)_{ij}|
 \le Cp\,Cn=Cd.
\]
Together with \cref{dcsbm:block-event} and
\(Z^T\Omega Z=\diag(n_1,\ldots,n_K)(pC)\diag(n_1,\ldots,n_K)\),
this proves the claim for true labels.

Changing at most \(Mn/d\) labels changes at most \(C Mn^2/d\)
ordered pair memberships.  Each reconstructed entry has size at
most \(Cp\), so the additional error is at most \(C Mn\).
This deterministic argument is uniform over the entire basin and
requires no independence of \(\widehat Z\).
\end{proof}

For fitted labels \(\widehat z\), let \(\widehat n_a\) be their sizes,
\(\overline D_a^{\,\widehat z}\) their mean degrees, and set
\begin{equation}
 \widehat\theta_j=\frac{D_j}{\overline D_{\widehat z_j}^{\,\widehat z}}.
 \label{dcsbm:theta-fit}
\end{equation}
Thus \(\sum_{\widehat z_j=a}\widehat\theta_j=\widehat n_a\)
exactly.  Fit the symmetric matrix \(\widehat B\) by maximizing
\begin{equation}
 \frac12\sum_{i\ne j}\left\{
 R^\circ_{ij}\log(\widehat\theta_i\widehat\theta_j B_{\widehat z_i\widehat z_j})
 +(1-R^\circ_{ij})
       \log(1-\widehat\theta_i\widehat\theta_j B_{\widehat z_i\widehat z_j})
 \right\}
 \label{dcsbm:block-criterion}
\end{equation}
over positive symmetric \(B\), under a fixed probability cap below one.
This is a precisely defined signed Bernoulli working criterion.
The next lemma proves that its relevant roots exist and lie in the
interior on \(\mathcal G\).

\begin{lemma}[Uniform fitted-parameter bounds]
\label{dcsbm:nuisance}
For every fixed \(M\), uniformly over the basin in
\cref{dcsbm:block-totals}, the degree estimates are bounded above
and below by positive fixed constants and
\begin{align}
 \frac{\|\widehat\theta-\theta\|_2}{\sqrt n}
 &\le C(M)d^{-1/2},                                      \label{dcsbm:theta-rate}\\
 \max_{a,b}\left|\frac{\widehat B_{ab}}{pC_{ab}}-1\right|
 &\le C(M)(d^{-1}+p)\le C(M)d^{-1}.                       \label{dcsbm:B-rate}
\end{align}
The last inequality uses \(dp=np^2\le C_1\).
\end{lemma}

\begin{proof}
Every fitted group has size comparable to \(n\), and
\cref{dcsbm:good} bounds its mean degree above and below by
constant multiples of \(d\).  This proves boundedness of
\(\widehat\theta\).
With true labels, write \(D_i=dS_{z_i}\theta_i+\varepsilon_i\),
where \(S_a=\sum_c(n_c/n)C_{ac}\) and
\(\|\varepsilon\|_2\le C\sqrt{nd}\).
Division by the true-group mean degree, followed by subtraction
of \(\theta_i\), gives normalized squared error \(O(d^{-1})\).
Changing an \(\eta\)-fraction of labels changes each fitted-group
mean degree by at most \(Cd\eta\).  For correctly labeled vertices
this changes the estimated theta by \(O(\eta)\); incorrectly labeled
vertices are an \(\eta\)-fraction and have bounded error.  Hence
\[
 \frac{\|\widehat\theta-\theta\|_2}{\sqrt n}
 \le C(d^{-1/2}+\eta+\sqrt\eta),
\]
which yields \cref{dcsbm:theta-rate} for \(\eta\le M/d\).

For the block fit, group the ordered pairs in
\cref{dcsbm:block-criterion} belonging to the unordered block pair
\(\{a,b\}\).  Denote this ordered set by \(\mathcal J_{ab}\), put
\(\nu_{ab}=2-\ind{a=b}\), \(t_{ij}=\widehat\theta_i\widehat\theta_j\),
and \(E_{ab}^{\,\circ}=\sum_{\mathcal J_{ab}}R^\circ_{ij}\).
Its score equation is
\begin{equation}
 E_{ab}^{\,\circ}
 = b\sum_{(i,j)\in\mathcal J_{ab}}
      \frac{(1-R^\circ_{ij})t_{ij}}{1-bt_{ij}} .
 \label{dcsbm:block-root}
\end{equation}
All \(1-R^\circ_{ij}\) are positive for large \(n\).
The right side is strictly increasing in \(b\) below the cap.
By \cref{dcsbm:block-totals},
\[
 E_{ab}^{\,\circ}
 =\nu_{ab}n_an_bpC_{ab}+O_M(n)\asymp n^2p>0.
\]
For \(b\asymp p\), bounded theta and \(|R^\circ_{ij}|\le Cp\)
give
\[
 \sum_{\mathcal J_{ab}}
   \frac{(1-R^\circ_{ij})t_{ij}}{1-bt_{ij}}
 =\sum_{\mathcal J_{ab}}t_{ij}\{1+O(p)\}
 =\nu_{ab}\widehat n_a\widehat n_b
       \{1+O(p)+O(n^{-1})\}.
\]
The last equality uses exact fitted-block mean normalization;
for \(a=b\), removing the diagonal subtracts only
\(\sum_{\widehat z_i=a}\widehat\theta_i^2=O(n)\).
Evaluating \cref{dcsbm:block-root} at sufficiently small and large
fixed multiples of \(p\) brackets its unique root, so
\[
 \widehat B_{ab}
 =\frac{E_{ab}^{\,\circ}}{\nu_{ab}\widehat n_a\widehat n_b}
       \{1+O(p)+O(n^{-1})\}.
\]
Since \(|\widehat n_a-n_a|\le Mn/d\), this proves
\cref{dcsbm:B-rate}.  The root is \(O(p)\), hence satisfies the
fixed cap strictly for large \(n\).
\end{proof}

\subsection{Exact Bernoulli profiling at bounded score error}

Fix a tested vertex \(i\) with true block \(a\).  For candidate
block \(b\), define true neighbor masses and normalized profiles by
\begin{equation}
 M_c^{(i)}=\sum_{\substack{j\ne i\\z_j=c}}\theta_j,\qquad
 T_b^{(i)}=\sum_c M_c^{(i)}C_{bc},\qquad
 q_{bc}^{(i)}=\frac{M_c^{(i)}C_{bc}}{T_b^{(i)}} .
 \label{dcsbm:q-true}
\end{equation}
The fitted analogues are
\begin{equation}
 \widehat M_c^{(i)}
 =\sum_{\substack{j\ne i\\\widehat z_j=c}}\widehat\theta_j
 =\widehat n_c-\widehat\theta_i\ind{\widehat z_i=c},
 \qquad
 \widehat q_{bc}^{(i)}
 =\frac{\widehat M_c^{(i)}\widehat B_{bc}}
        {\sum_e\widehat M_e^{(i)}\widehat B_{be}}.
 \label{dcsbm:q-fitted}
\end{equation}
The factor \(p\) is not needed to evaluate the fitted quantity.
All these probabilities are bounded away from zero.
By \cref{dcsbm:nuisance} and fitted group-size accuracy,
\begin{equation}
 \max_{i,a,b,c}
 \left|\log\frac{\widehat q_{ac}^{(i)}}{\widehat q_{bc}^{(i)}}
          -\log\frac{q_{ac}^{(i)}}{q_{bc}^{(i)}}\right|
 \le C(M)d^{-1}.
 \label{dcsbm:q-rate}
\end{equation}
Indeed normalization and logarithms are smooth on the fixed compact
positive parameter set, and the relative mass error is
\(O_M(d^{-1}+n^{-1})\).

\begin{lemma}[Bernoulli profile reduction]
\label{dcsbm:profile-reduction}
Let a hollow row \(r=(r_j)_{j\ne i}\) have
\(\sum_jr_j=D_i\ge0\), \(\sum_j|r_j|\le Ld\), and \(D_i\le Ld\).
Suppose either \(r_j\in\{0,1\}\), or
\(\max_j|r_j|\le Lp=o(1)\).
For bounded positive neighbor parameters and block probabilities
of order \(p\), maximize the exact Bernoulli row criterion over
the target parameter \(t\ge0\), with a fixed probability cap below
one.  After dropping terms independent of the candidate label,
the profile contrast between \(a\) and \(b\) is
\begin{equation}
 L^{ab}_{r,\rm prof}
 =\sum_c N_c^r\log\frac{q_{ac}^{(i)}}{q_{bc}^{(i)}}
       +O(np^2),\qquad
 N_c^r=\sum_{\substack{j\ne i\\z_j=c}}r_j .
 \label{dcsbm:profile-linear}
\end{equation}
The same assertion holds with fitted labels, thetas, block probabilities,
and profiles.  Fixed positive candidate prior probabilities add their
bounded log-ratio to both sides.
\end{lemma}

\begin{proof}
For true parameters, omit the candidate-independent terms
\(\sum_jr_j\log\theta_j+D_i\log p\).  The remaining score is
\[
 \ell_r(b,t)
 =D_i\log t+\sum_cN_c^r\log C_{bc}
   +\sum_{j\ne i}(1-r_j)\log(1-tp\theta_jC_{bz_j}).
\]
For bounded \(t\), expansion of the logarithm gives uniformly
\begin{equation}
 \ell_r(b,t)
 =D_i\log t+\sum_cN_c^r\log C_{bc}-tpT_b^{(i)}
     +O(np^2).
 \label{dcsbm:profile-expansion}
\end{equation}
The restored edge part of the linear exposure is bounded by
\(Cp\sum_j|r_j|\le Cdp=Cnp^2\).
The quadratic remainder is bounded by
\(Cp^2(n+\sum_j|r_j|)\le Cnp^2\).

This expansion remains valid after maximizing.  In both allowed
row types \(1-r_j\ge0\), so the exact score is concave in \(t\).
For a raw row, at least \(n-O(d)\) nonedge terms remain; for a
reconstructed row, every nonedge coefficient is at least \(1/2\)
for large \(n\).  Thus the nonedge part of its derivative has
magnitude at least \(cd\), whereas the positive term is \(D_i/t\).
The derivative is negative above a sufficiently large fixed \(t_0\),
uniformly in the candidate.  The cap allows this bounded interval
for large \(n\), and the leading optimizer
\(t_b^0=D_i/(pT_b^{(i)})\) is also bounded.
Taking suprema in \cref{dcsbm:profile-expansion} is therefore valid.
The candidate-dependent leading maximum is
\(\sum_cN_c^r\log C_{bc}-D_i\log T_b^{(i)}\).
Since \(\sum_cN_c^r=D_i\), its difference is exactly the sum
in \cref{dcsbm:profile-linear}.

With fitted parameters replace \(pC\) by \(\widehat B\) and
the true masses by their fitted analogues.  Boundedness and
\(\widehat B\asymp p\) give the identical remainder calculation.
When \(D_i=0\), use the continuous extension after combining the
terms into \(D_i\log t\), with \(0\log0=0\).  Nonnegative nonedge
coefficients make \(t=0\) optimal, and the same contrast identity
holds.  In particular a raw zero row gives zero contrast.
\end{proof}

Exact preservation of \(D_i\) is essential in this argument:
it fixes the coefficient of \(\log t\) and the profiled total-degree
term.  The proof needs no \(O(d^{-1})\) estimate of an individual
target theta.

\begin{lemma}[Uniform replacement of all fitted stages]
\label{dcsbm:uniform-score}
Fix \(M\).  For every partition in the \(M/d\) basin, fitted by
\cref{dcsbm:theta-fit,dcsbm:block-criterion}, its actual row-profile
contrast satisfies, simultaneously for all \(i,a,b\),
\begin{equation}
 \widehat L_i^{ab}
 =[R^\circ w_{i,ab}]_i+O_M(1),\qquad
 (w_{i,ab})_j=\log\frac{q_{a,z_j}^{(i)}}{q_{b,z_j}^{(i)}} .
 \label{dcsbm:uniform-contrast}
\end{equation}
The formula is uniform over data-dependent labels and parameters
obtained by these fits, under the same common label permutation.
\end{lemma}

\begin{proof}
The profile remainder is \(O(np^2)=O(1)\).
Replacing \(\widehat Z\) by \(Z\) in a bounded block coefficient
vector changes row \(i\)'s score by at most
\[
 \max_j|R^\circ_{ij}|\,
     C\#\{j:\widehat z_j\ne z_j\}
 \le Cp\,CMn/d=C(M).
\]
Replacing the fitted coefficients by their true counterparts
costs at most
\(
 \sum_j|R^\circ_{ij}|\,
 C(M)d^{-1}\le C(M)
\)
by \cref{dcsbm:q-rate}.  These are deterministic bounds holding
for every partition in the basin, which proves the claimed uniformity.
\end{proof}

\subsection{The profile reference and its sharp directional tails}

For the fixed target \(i\), true block \(a\), and competitor \(b\),
write \(w_j=(w_{i,ab})_j\), \(v=\theta_i\), and
\[
 p_j=vp\theta_jC_{az_j},\qquad
 T_{i,ab}=\sum_{j\ne i}w_jA_{ij},\qquad
 K_{i,ab}(s)=\sum_{j\ne i}\log\{1+p_j(e^{-sw_j}-1)\}.
\]
Set
\begin{equation}
 J_{i,a\to b}
 =-\min_{s>0}K_{i,ab}(s).
 \label{dcsbm:profile-exponent}
\end{equation}
This is the exact Bernoulli transform of the leading profiled
contrast.  It is not asserted to be the simple-hypothesis
Chernoff information for the known-target-theta oracle.

\begin{lemma}[Uniform profile saddle and reference scale]
\label{dcsbm:profile-tail}
The minimizing \(s=s_{i,ab}\) stays in a fixed compact subinterval
of \((0,1)\).  Under the exact tilted product law,
\begin{equation}
 r_j=\frac{p_je^{-sw_j}}{1-p_j+p_je^{-sw_j}},\qquad
 \sum_j r_jw_j=0,\qquad
 \sum_j r_j(1-r_j)w_j^2\asymp d .
 \label{dcsbm:tilt}
\end{equation}
Uniformly for bounded real thresholds \(x\),
\begin{equation}
 \PP_a(T_{i,ab}<x)\asymp
 d^{-1/2}\exp\{-J_{i,a\to b}\}.
 \label{dcsbm:directional-tail}
\end{equation}
The exact full-edge Bernoulli profile-reference error at vertex \(i\)
is comparable, by fixed constants, to
\begin{equation}
 \sum_{b\ne a}d^{-1/2}\exp\{-J_{i,a\to b}\}.
 \label{dcsbm:reference-sum}
\end{equation}
\end{lemma}

\begin{proof}
Uniformly on bounded \(s\) intervals, the transform and its first
two derivatives satisfy
\[
 d^{-1}K_{i,ab}(s)
 =v\,\frac{T_a^{(i)}}n
    \left\{\sum_c(q_{ac}^{(i)})^{1-s}(q_{bc}^{(i)})^s-1\right\}
       +O(p).
\]
The derivative of the leading expression is strictly negative
at zero and strictly positive at one.  Positivity, the lower
singular-value bound on \(C\), and the group-mass bounds make
these signs and the positive second derivative uniform.
Indeed two normalized rows could coincide only if two rows of
\(C\) were proportional; compactness of the stated parameter
class excludes even an arbitrarily small normalized separation.
Strict convexity now gives the interior saddle and tilted
variance of order \(d\).

The bounded-weight, variance, and third-moment hypotheses of
\cref{lem:sharp} hold, so its bounded-shift tail conclusion proves
\cref{dcsbm:directional-tail}.  In particular, no nonlattice
assumption is necessary.  Alternatively, under the displayed
tilt, Berry--Esseen bounds every unit-interval mass by \(C/\sqrt d\);
summing against \(e^{sT}\) gives the upper bound.  A sufficiently
long fixed interval immediately below \(x\) has Gaussian mass
larger than twice the Berry--Esseen error, giving the lower bound.

The probability of \(D_i>Ld\) can be made
\(\exp(-Ld\,c_L)\) with arbitrarily large fixed exponential rate
by increasing the fixed upper cutoff.  On its complement,
\cref{dcsbm:profile-reduction} compares the exact reference
contrast to \(T_{i,ab}\) up to a fixed additive constant.
The upper-degree exception is negligible relative to the displayed
tail scales.  Each strict pairwise defeat of the true label implies
an error, while an error is contained in the union of weak pairwise
defeats.  Fixed \(K\), the two bounded-shift tail bounds, and fixed
ties therefore prove \cref{dcsbm:reference-sum}.
\end{proof}

\subsection{The complete degree-weighted tilted cancellation}

We now work on the genuinely deleted graph \(A^{(-i)}\).
To simplify notation within this subsection, \(D_-\), \(U_-\),
\(\Pi_-\), and
\[
 H_-=D_-^{-1/2}\Pi_-D_-^{1/2}
\]
refer only to this independent training graph.  Its completed mean
is \(\Omega_-\), and \(D_{0,-}=\diag(\Omega_-\one)\).
The tested row remains independent Bernoulli conditional on the
deleted graph.  Let \(w=(w_j)_{j\ne i}\) be the true block coefficient
vector above, and define \(e=(H_--I)w\).

\begin{lemma}[Bounded tilted coefficient error]
\label{dcsbm:weighted-cancellation}
On the deleted-graph good event, the exact profile tilt in
\cref{dcsbm:tilt} satisfies
\begin{equation}
 \max_j|e_j|\le C,\qquad
 \sum_jr_je_j^2\le C,\qquad
 \left|\sum_jr_je_j\right|\le C.
 \label{dcsbm:three-moments}
\end{equation}
All constants are uniform over the vertex and ordered pair of
distinct candidate blocks.
\end{lemma}

\begin{proof}
The residual lemma on the deleted graph gives
\[
 \|(I-\Pi_-)D_-^{1/2}Z_-\|_{\op}\le C\sqrt n.
\]
Consequently, for every bounded block vector \(w\),
\begin{equation}
 \|(I-H_-)w\|_{D_-}\le C\sqrt n,\qquad
 \|(I-H_-)w\|_2^2\le Cn/d.
 \label{dcsbm:weighted-residual}
\end{equation}
Incoherence bounds the coordinates of \(H_-w\):
\[
 |(H_-w)_j|
 \le D_{-,j}^{-1/2}\|U_{-,j\cdot}\|_2
          \|U_-^TD_-^{1/2}w\|_2
 \le (C/\sqrt d)(C/\sqrt n)C\sqrt{nd}\le C.
\]
This proves bounded coefficients.  Since \(r_j\le Cp\),
\cref{dcsbm:weighted-residual} gives
\(\sum_jr_je_j^2\le Cp\,Cn/d\le C\).

For the mean, exact Bernoulli tilting gives
\begin{equation}
 r_j=vp\theta_j C_{az_j}e^{-sw_j}+\varepsilon_j,
 \qquad \max_j|\varepsilon_j|\le Cp^2.
 \label{dcsbm:tilted-mean}
\end{equation}
For a remaining vertex in block \(c\), write
\[
 T_c^-=\sum_{\ell\ne i}\theta_\ell C_{cz_\ell},\qquad
 h_j=\frac{v C_{ac}e^{-sw_c}}{T_c^-}.
\]
Then \(h\) is block constant with \(\|h\|_\infty\le C/n\), and
the leading vector in \cref{dcsbm:tilted-mean} is exactly
\(D_{0,-}h\).  Because \(H_-\) is \(D_-\)-orthogonal,
\begin{align}
 r^T(H_--I)w
 ={}&-\langle(I-H_-)h,(I-H_-)w\rangle_{D_-}
       \nonumber\\
 &+\big((D_{0,-}-D_-)h\big)^T(H_--I)w
   +\varepsilon^T(H_--I)w .
 \label{dcsbm:weighted-identity}
\end{align}
This identity retains the random-degree correction explicitly.
The first factor in its first term is at most \(C/\sqrt n\):
apply the block residual estimate to the block coefficients of
\(h\), which have magnitude \(O(n^{-1})\).
The second factor is at most \(C\sqrt n\), so that term is \(O(1)\).

For the next term, the degree aggregate in
\cref{dcsbm:residual-estimates} gives
\[
 \|(D_{0,-}-D_-)h\|_{D_-^{-1}}^2
 =\sum_j\frac{(D_{0,-,j}-D_{-,j})^2h_j^2}{D_{-,j}}
 \le C/n.
\]
Cauchy--Schwarz in the degree metric and
\cref{dcsbm:weighted-residual} again give \(O(1)\).
Finally, bounded \(e_j\) and
\cref{dcsbm:tilted-mean} bound the last term by
\(Cnp^2=O(1)\).  This proves the tilted mean bound and the lemma.
\end{proof}

The middle term of \cref{dcsbm:weighted-identity} is why replacing
random degrees by population degrees without a remainder is invalid.
The product of two residual bounds is what reduces the leading
tilted bias to a constant.

\subsection{Completion of the risk bound on the common good event}

By \cref{dcsbm:initializer}, the initial partition lies in a fixed
\(M/d\) basin on \(\mathcal G\).  Apply one nuisance fit and one
profile assignment.  The uniform comparison
\cref{dcsbm:uniform-contrast} and the actual full-to-deleted-graph
scalar lemma \cref{lem:loo} give
\begin{equation}
 \widehat L_i^{ab}
 =\sum_{j\ne i}A_{ij}(H_-w_{i,ab,-i})_j+O_M(1)
 =T_{i,ab}+\sum_{j\ne i}A_{ij}e_j+O_M(1).
 \label{dcsbm:actual-to-loo}
\end{equation}
The leave-one-out graph appears only in this proof identity; the
implemented classifier uses the single full-graph encoding.
The \(O_M(1)\) error is a deterministic uniform bound, so it need
not be asserted to be independent of the tested row.

Condition on the genuinely deleted graph.  On its measurable good
event, \cref{dcsbm:weighted-cancellation} supplies all three
coefficient conditions of the sharp transfer lemma
\cref{lem:sharp}.  Therefore, for a fixed bounded threshold large
enough to absorb \cref{dcsbm:actual-to-loo},
\begin{equation}
 \PP_a\!\left\{
     \sum_{j\ne i}A_{ij}(H_-w_{i,ab,-i})_j\le C(M)
       \,\middle|\, A^{(-i)}
          \right\}
 \le C(M)d^{-1/2}e^{-J_{i,a\to b}} .
 \label{dcsbm:conditional-risk}
\end{equation}
The common full event implies the required minor event by the
common geometry lemmas.  Hence restricting an actual error to
\(\mathcal G\) is contained in the event on the left of
\cref{dcsbm:conditional-risk} together with the measurable minor
event.  This does not condition the tested row on a dependent
full-graph event.

Integrating over the deleted graph and summing over \(b\ne a\)
gives
\[
 \PP\{\widehat z_i\ne z_i,\mathcal G\}
 \le C\sum_{b\ne z_i}d^{-1/2}e^{-J_{i,z_i\to b}}.
\]
Here labels are aligned by the common initial permutation.
The permutation-invariant misclassification loss is no larger
than the loss under that alignment.  By
\cref{dcsbm:profile-tail}, the last sum is at most a constant
times the exact full-edge profile-reference error at \(i\).
Averaging proves
\begin{equation}
 \E\{\ell(\widehat z,z)\ind{\mathcal G}\}
 \le C R_{\mathrm{prof},*}.
 \label{dcsbm:on-good-risk}
\end{equation}
The calculation in \cref{sec:probability} proves
\(\PP(\mathcal G^c)\le C R_{\mathrm{prof},*}\), including the
minimum-degree exception at the correct exponential scale.
Because the fallback loss is at most one, adding that contribution
to \cref{dcsbm:on-good-risk} completes \cref{thm:dc}.

\subsection{Optional repeated hard fitting and finite adaptive stopping}

The same theorem holds if, after initialization, the algorithm
repeatedly recomputes \cref{dcsbm:theta-fit,dcsbm:block-criterion}
and makes a hard profile assignment using the same \(R^\circ\).
It may return any data-dependent finitely stopped iterate after
at least one completed profile assignment.  A finite iteration
cap is one explicit termination rule.

\begin{lemma}[A uniform hard-refit basin]
\label{dcsbm:basin}
There is a fixed \(M\) such that, on \(\mathcal G\) and for all
sufficiently large \(n\), every hard fit-and-profile update maps
the \(M/d\) error basin into itself, under the same label matching.
Every error at every updated iterate is contained in
\begin{equation}
 [R^\circ w_{i,ab}]_i\le C(M)
 \quad\text{for some }b\ne z_i .
 \label{dcsbm:common-error-event}
\end{equation}
\end{lemma}

\begin{proof}
First,
\begin{equation}
 \|R^\circ-\Omega\|_F^2\le C_{\rm spec}d.
 \label{dcsbm:frobenius}
\end{equation}
Indeed \(S_K=U\diag(\mu)U^T\) is a best rank-\(K\)
operator-norm approximation to \(S\).  Since \(M_D\) has rank
\(K\), \(\|S_K-M_D\|_{\op}\le2\|S-M_D\|_{\op}\le C/\sqrt d\).
Their difference has rank at most \(2K\), so its Frobenius norm
is also \(O(d^{-1/2})\).  Multiplication by \(D^{1/2}\) on both
sides gives \(\|R-\Omega\|_F\le C\sqrt d\).
The hollow term has Frobenius norm at most
\[
 \max_i|s_i|\,\|H-I\|_F\le Cp\,C\sqrt n=Cd/\sqrt n,
\]
whose square is at most \(Cd\), proving
\cref{dcsbm:frobenius}.

The completed population contrast has a uniform positive margin:
with \(\lambda_i=\sum_{j\ne i}p\theta_i\theta_jC_{z_i z_j}\),
\begin{equation}
 [\Omega w_{i,ab}]_i
 =\lambda_i\KL(q_a^{(i)}\|q_b^{(i)})+O(p)
 \ge c_{\rm gap}d,\qquad a=z_i,\ b\ne a.
 \label{dcsbm:population-margin}
\end{equation}
The \(O(p)\) term is the completed diagonal.  The uniform KL
gap follows from the same compactness and full-rank argument as
in \cref{dcsbm:profile-tail}.  Also
\(\|w_{i,ab}\|_2\le C\sqrt n\).

For an input partition in a fixed \(M/d\) basin,
\cref{dcsbm:uniform-score} bounds its fitted-contrast discrepancy
by \(C(M)\).  For large \(n\), this is at most
\(c_{\rm gap}d/2\).  A new erroneous assignment therefore requires
\[
 \|(R^\circ-\Omega)_{i\cdot}\|_2
 \ge \frac{c_{\rm gap}d}{2C\sqrt n}.
\]
Summing this squared requirement and applying
\cref{dcsbm:frobenius} shows that the new iterate has at most
\(C_{\rm new}n/d\) errors.  The constant \(C_{\rm new}\) depends
only on the fixed spectral and model constants, not on \(M\).
Choose \(M\) at least both \(C_{\rm new}\) and the initialization
constant.  All sufficiently large \(n\) then give basin invariance.
Group indices remain matched to the same true indices throughout:
the uniform score comparison and its population margin use that
fixed matching at every step.

Finally every erroneous assignment directly satisfies the
fixed-threshold canonical event in the statement by
\cref{dcsbm:uniform-score}.  The threshold is the same for all
partitions in the basin.
\end{proof}

By induction, \cref{dcsbm:basin} applies to every iterate on the
single graph event \(\mathcal G\).  The final error of an adaptively
chosen finite iterate is thus contained in the same canonical
event used to prove \cref{dcsbm:conditional-risk}.
Applying that probability bound once proves the same sharp risk
bound; there is no union over iterations and no stopping-time
independence assertion.  This corollary concerns these specified
hard updates.  It does not assert the same result for an unrelated
soft update or a reconstruction changed during iteration.

\subsection{Equal exposures and the original known-target-theta oracle}

We finish by proving \cref{cor:known}.  Suppose
\begin{equation}
 \sum_c\frac{n_c}{n}C_{ac}=S
 \quad\text{for every }a.
 \label{dcsbm:equal-exposure}
\end{equation}
The block means of theta equal one, so
\(
 T_a^{(i)}=nS-\theta_iC_{a,z_i}
\).
In particular \(T_a^{(i)}-T_b^{(i)}=O(1)\) and
\(\log(T_a^{(i)}/T_b^{(i)})=O(n^{-1})\).
Let \(v=\theta_i\) be supplied to the full-edge
known-target-theta oracle.  Its exact Bernoulli contrast is
\[
 L^{ab}_{A,\rm known}
 =\sum_cN_c\log\frac{C_{ac}}{C_{bc}}
  +\sum_{j\ne i}(1-A_{ij})
       \log\frac{1-vp\theta_jC_{az_j}}
                 {1-vp\theta_jC_{bz_j}} .
\]
On \(D_i\le Ld\), the same direct Bernoulli expansion used above
gives
\[
 L^{ab}_{A,\rm known}
 =\sum_cN_c\log\frac{C_{ac}}{C_{bc}}
   -vp(T_a^{(i)}-T_b^{(i)})+O(np^2).
\]
By \cref{dcsbm:profile-reduction}, the exact profile contrast is
\[
 L^{ab}_{A,\rm prof}
 =\sum_cN_c\log\frac{C_{ac}}{C_{bc}}
    -D_i\log\frac{T_a^{(i)}}{T_b^{(i)}}+O(np^2).
\]
Thus, using \cref{dcsbm:equal-exposure},
\begin{equation}
 |L^{ab}_{A,\rm prof}-L^{ab}_{A,\rm known}|
 \le C(np^2+p)=O(1).
 \label{dcsbm:known-comparison}
\end{equation}
Fixed prior log-ratios, if present, are identical on both sides.
The bounded-shift upper and lower tails in
\cref{dcsbm:directional-tail} show that the two directional
error probabilities are comparable by constants.
The upper-degree exception has arbitrarily large fixed
exponential rate after enlarging \(L\), and is negligible for
both references.  Finite unions over competitors and averaging
over vertices yield
\[
 R_{\mathrm{prof},*}\asymp R_{\theta,*}.
\]
Combining this with \cref{thm:dc} proves \cref{cor:known}.
The condition permits heterogeneous theta, unequal block sizes,
and general positive full-rank symmetric \(C\) obeying its exposure
constraints.  Its role is precisely to remove the label information
carried by the leading total-degree exposure, which profiling
otherwise changes.

\subsection{The classical balanced DC-SBM oracle and its sharp risk}
\label{dcsbm:classical-oracle}

The fundamental local testing problem used by
Gao, Ma, Zhang and Zhou~\cite{gmzz2018} fixes the target degree parameter
under both candidate labels. Its exact simple-hypothesis likelihood test
therefore corresponds to \(R_{\theta,*}\) in \cref{sec:oracle}.
The next proposition explicitly connects that reference, the profile
reference and true-group neighbor voting in the balanced planted-partition
model. This is an oracle comparison with its own assumptions:
it uses neither the spectral good event \(\mathcal G\) nor a logarithmic
lower bound on \(np\).

\begin{proposition}[Sharp alignment of the classical DC-SBM references]
\label{dcsbm:classical-bridge}
\label{prop:classical-dc-oracle-bridge}
Fix an arbitrary \(K\ge2\), positive bounds
\(\theta_-,\theta_+\), and \(0<r_-<r_+<1\).
Let \(n=Km\), let every group have \(m\) vertices, and suppose
\[
 \theta_-\le\theta_i\le\theta_+,\qquad
 \sum_{i:z_i=a}\theta_i=m\quad\text{for all }a.
\]
Specialize the block matrix \(B=pC\) to
\begin{equation}
 B_{aa}=p,\qquad B_{ab}=q\ (a\ne b),\qquad
 r_-\le q/p\le r_+ .
 \label{dcsbm:classical-model}
\end{equation}
Assume only, for the present oracle comparison,
\begin{equation}
 d=np\longrightarrow\infty,\qquad np^2\le C_1<\infty.
 \label{dcsbm:classical-oracle-range}
\end{equation}
Define \(R_{\mathrm{vote},*}\) as the average error of the local raw-row
rule which knows the neighbor labels and selects the group with the
largest incident block count. Use any fixed tie ordering.
Then, uniformly over this parameter class,
\begin{equation}
 \begin{aligned}
 R_{\theta,*}
 &\asymp R_{\mathrm{prof},*}
 \asymp R_{\mathrm{vote},*}\\
 &\asymp
 \frac1n\sum_{i=1}^n
 \frac1{\sqrt{np}}
 \exp\left\{-\theta_i\frac nK(\sqrt p-\sqrt q)^2\right\}.
 \end{aligned}
 \label{dcsbm:classical-risk}
\end{equation}
The profile reference is exactly the Bernoulli profile in
\eqref{eq:profile-oracle}, with candidate probabilities capped at \(1/2\).
The same comparison holds if the voting counts are normalized by the
self-excluded neighbor-group sizes. Constants depend only on the displayed
fixed bounds, not on \(n\). In particular,
\eqref{dcsbm:classical-oracle-range} permits \(np=o(\log n)\).
\end{proposition}

\begin{proof}
Fix a target \(i\), its true group \(a\), and a competitor \(b\ne a\).
For this proof write
\[
 N_c=\sum_{\substack{j\ne i\\z_j=c}}A_{ij},\qquad
 M_c=\sum_{\substack{j\ne i\\z_j=c}}\theta_j
     =m-\theta_i\ind{c=a},\qquad
 L=\log(p/q).
\]
Thus \(L\) is bounded above and below by positive constants.
Put
\[
 V_{i,ab}=L(N_a-N_b),\qquad
 J_i=\theta_i m(\sqrt p-\sqrt q)^2 .
\]
The assumptions imply \(J_i\asymp d\), uniformly in \(i\).
They also imply \(p\to0\), so every true candidate edge probability
is eventually less than \(1/2\).

\medskip
\noindent\textbf{The exact Bernoulli count transform.}
Under the actual true-row law, independence gives
\begin{align}
 \Psi_{i,ab}(s)
 &:=\log\E_a e^{-sV_{i,ab}}\nonumber\\
 &=\sum_{\substack{j\ne i\\z_j=a}}
   \log\{1+\theta_i\theta_jp(e^{-sL}-1)\}
   +\sum_{j:z_j=b}
   \log\{1+\theta_i\theta_jq(e^{sL}-1)\}.
 \label{dcsbm:classical-cumulant}
\end{align}
Taylor expansion, uniformly for \(s\in[0,1]\), gives
\begin{equation}
 \Psi_{i,ab}(s)
 =m\theta_i\{p(e^{-sL}-1)+q(e^{sL}-1)\}
       +O(np^2+p).
 \label{dcsbm:classical-cumulant-expansion}
\end{equation}
The same error bound holds for the first two derivatives. Bounded theta
and bounded \(L\) make the sum of quadratic Bernoulli remainders
\(O(np^2)\). The additional \(O(p)\) term accounts for the omitted
\(\theta_i\) in \(M_a=m-\theta_i\).

The leading transform in
\eqref{dcsbm:classical-cumulant-expansion} is minimized at \(s=1/2\),
with minimum \(-J_i\). Its endpoint derivatives have opposite signs,
each of magnitude at least \(cd\), and its second derivative is
between positive constant multiples of \(d\).
Since \(np^2+p=o(d)\), the exact minimizing saddle stays in a fixed
compact subinterval of \((0,1)\). Convexity makes it the global
minimum over \(s\ge0\). In particular,
\begin{equation}
 I^{\mathrm{cnt}}_{i,ab}
 :=-\min_{s\ge0}\Psi_{i,ab}(s)
 =J_i+O(np^2+p)=J_i+O(1).
 \label{dcsbm:classical-information}
\end{equation}
One can also obtain this bound directly by taking minima in the
uniform transform approximation.

Under the exact optimizing tilt, the Bernoulli probabilities on the
two relevant groups are comparable to \(p\). Their sum is \(O(d)\),
and the variance of \(V_{i,ab}\) is comparable to \(d\).
Consequently \cref{lem:sharp} applies to
\eqref{dcsbm:classical-cumulant} without any spectral event. For every
fixed \(x_0<\infty\), uniformly over \(|x|\le x_0\),
\begin{equation}
 c_{x_0}d^{-1/2}e^{-J_i}
 \le \PP_a(V_{i,ab}<x)
 \le \PP_a(V_{i,ab}\le x)
 \le C_{x_0}d^{-1/2}e^{-J_i}.
 \label{dcsbm:classical-count-tail}
\end{equation}
The strict lower bound from that lemma uses an open interval of fixed
length under the tilt. Thus lattice atoms and an arbitrary fixed
tie ordering cause no loss of the lower bound.

\medskip
\noindent\textbf{Comparison with the known-target likelihood.}
Let
\[
 \mathcal T_c=\sum_{j\ne i}\theta_j B_{c,z_j}
              =pT_c^{(i)}.
\]
By the exact group normalization,
\begin{equation}
 \begin{aligned}
 \mathcal T_a&=m[p+(K-1)q]-\theta_ip,\\
 \mathcal T_b&=m[p+(K-1)q]-\theta_iq.
 \end{aligned}
 \label{dcsbm:classical-exposure}
\end{equation}
In particular \(\mathcal T_a,\mathcal T_b\asymp d\),
\(\mathcal T_a-\mathcal T_b=O(p)\), and
\(\log(\mathcal T_a/\mathcal T_b)=O(n^{-1})\).

Fix a sufficiently large constant \(L_0\) to be chosen below, and consider
the single-target event
\[
 \mathcal E_i=\{D_i\le L_0d\}.
\]
The exact known-target Bernoulli contrast is
\[
 L^{ab}_{i,\mathrm{known}}
 =V_{i,ab}+
 \sum_{j\ne i}(1-A_{ij})
 \log\frac{1-\theta_i\theta_jB_{a,z_j}}
          {1-\theta_i\theta_jB_{b,z_j}} .
\]
Expansion of these nonedge logarithms on \(\mathcal E_i\) gives
\begin{align}
 L^{ab}_{i,\mathrm{known}}
 &=V_{i,ab}-\theta_i(\mathcal T_a-\mathcal T_b)
       +O(D_ip+np^2)\nonumber\\
 &=V_{i,ab}+O(np^2+p).
 \label{dcsbm:classical-known-contrast}
\end{align}
The \(O(D_ip)\) term restores the incident-edge terms in the linear
exposure; the quadratic remainder is \(O(np^2)\).

\medskip
\noindent\textbf{Comparison with the capped Bernoulli profile.}
For a candidate \(c\), put \(b_{cj}=\theta_jB_{c,z_j}\).
The admissible interval in \eqref{eq:profile-oracle} is
\[
 0\le t\le t_{\max,c}:=\frac1{2\max_{j\ne i}b_{cj}}.
\]
Because \(b_{cj}\asymp p\), its upper endpoint tends to infinity.
For \(t>0\) the derivative of the exact raw-row likelihood is
\[
 \frac{D_i}{t}
 -
 \sum_{j\ne i}(1-A_{ij})
       \frac{b_{cj}}{1-tb_{cj}} .
\]
On \(\mathcal E_i\), there are at least \(n-1-L_0d\ge n/2\)
nonedges for all sufficiently large \(n\). Their total intensity is
at least \(cd\), uniformly over candidates. Hence the derivative is
negative for every admissible \(t\) above a fixed constant depending
only on \(L_0\) and the structural bounds. Every exact maximizing
\(t\) is therefore bounded. The cap at \(1/2\) is inactive on this
bounded interval for all large \(n\).

Uniformly over the resulting bounded \(t\) interval, the likelihood is
\begin{equation}
 D_i\log t+\sum_{j\ne i}A_{ij}\log b_{cj}
      -t\mathcal T_c+O(np^2).
 \label{dcsbm:classical-profile-expansion}
\end{equation}
Indeed, restoring the edge terms in the linear exposure costs
\(O(D_ip)\), and the sum of quadratic nonedge remainders costs
\(O(np^2)\). Both are \(O(np^2)\) on \(\mathcal E_i\).
The leading optimizer \(D_i/\mathcal T_c\) is bounded as well, so
the same error bound holds after maximization.
Terms \(D_i\log D_i-D_i\) and
\(\sum_{j\ne i}A_{ij}\log\theta_j\) cancel in a contrast.
Thus the exact capped profile satisfies
\begin{align}
 L^{ab}_{i,\mathrm{prof}}
 &=V_{i,ab}
   -D_i\log(\mathcal T_a/\mathcal T_b)+O(np^2)\nonumber\\
 &=V_{i,ab}+O(np^2+p)
 \quad\text{on }\mathcal E_i.
 \label{dcsbm:classical-profile-contrast}
\end{align}
For \(D_i=0\), the raw row is identically zero; \(t=0\) maximizes
every candidate likelihood with value zero. The displayed formulas
then hold with the continuous convention \(0\log0=0\).

\medskip
\noindent\textbf{Single-row exceptions, directional errors, and averaging.}
Every mean degree is at most \(\lambda_+d\) for a fixed \(\lambda_+\).
For any prescribed fixed \(H\), a Bernoulli Chernoff bound permits a
fixed choice of \(L_0\) such that
\begin{equation}
 \sup_i\PP(\mathcal E_i^c)\le e^{-Hd}.
 \label{dcsbm:classical-single-exception}
\end{equation}
Also \(J_i\le C_Jd\) uniformly. Choose \(H>C_J+1\).
It follows from \(d\to\infty\) that
\[
 \sup_i
 \frac{\PP(\mathcal E_i^c)}
      {d^{-1/2}e^{-J_i}}
 \le \sqrt d\,e^{-(H-C_J)d}\longrightarrow0 .
\]
This calculation concerns one target at a time; it introduces no
factor \(n\) and requires no logarithmic lower bound on \(d\).

By \eqref{dcsbm:classical-known-contrast} and
\eqref{dcsbm:classical-profile-contrast}, a fixed constant \(B_0\)
bounds the difference between \(V_{i,ab}\) and either exact contrast
on \(\mathcal E_i\). Thus for either
\(L^{ab}_i=L^{ab}_{i,\mathrm{known}}\) or
\(L^{ab}_i=L^{ab}_{i,\mathrm{prof}}\),
\begin{align*}
 \PP_a(V_{i,ab}<-B_0)-\PP_a(\mathcal E_i^c)
 &\le \PP_a(L^{ab}_i<0)\\
 &\le \PP_a(L^{ab}_i\le0)\\
 &\le \PP_a(V_{i,ab}\le B_0)+\PP_a(\mathcal E_i^c).
\end{align*}
Applying \eqref{dcsbm:classical-count-tail} and then
\eqref{dcsbm:classical-single-exception} proves that both the strict
and nonstrict directional errors of either oracle are comparable
to \(d^{-1/2}e^{-J_i}\).
These are unconditional event inclusions; no independence from
\(\mathcal E_i\) is asserted or needed.

For the raw block-count vote, the pairwise contrast is \(V_{i,ab}\)
itself. If counts are normalized by the self-excluded group sizes,
multiplying all candidate scores by \(mL\) changes the pairwise
contrast by at most \(CLD_i/n=O(p)\) on \(\mathcal E_i\), since
those sizes are \(m\) or \(m-1\). The same tail comparison therefore
applies to either voting convention.

For each of these coherent families of candidate scores, a strict
pairwise defeat of the true group implies an error, and an error
implies a nonstrict defeat by some competitor. Taking the maximum
of the strict pair probabilities for a lower bound and their
nonstrict sum for an upper bound proves that each per-target error
is comparable to \(d^{-1/2}e^{-J_i}\), since \(K\) is fixed.
Averaging these bounds over \(i\) proves
\eqref{dcsbm:classical-risk}.
\end{proof}

\begin{remark}[Oracle alignment and the algorithmic range]
\label{dcsbm:classical-scope}
The degree-dependent exponent in \eqref{dcsbm:classical-risk} is the
balanced DC-SBM exponent appearing in the minimax analysis of
Gao, Ma, Zhang and Zhou~\cite{gmzz2018}. Their exponent-level
minimax expression does not, by itself, give the sharp prefactor
or a constant-ratio theorem for every parameter point.
The proposition supplies the exact local-oracle comparison on the
fixed-separation, bounded-theta class stated here.

For \eqref{dcsbm:classical-model}, the finite-\(n\) exposures in
\eqref{eq:equal-exposure} are exactly equal. Hence, on the overlap
with the algorithmic range of \cref{thm:dc},
\cref{cor:known,dcsbm:classical-bridge} give
\[
 \E\ell(\widehat z_D,z)\le C R_{\theta,*},
 \qquad R_{\theta,*}\asymp R_{\mathrm{prof},*}.
\]
The oracle proposition alone does not remove the logarithmic
restriction from the earlier unregularized \cref{thm:dc}.
The required sparse spectral and likelihood-transfer estimates are
now supplied in \cref{sec:sparse-restored} for the explicit regularized
decoder. Consequently \cref{thm:regularized,restore:dc-theorem} give
the displayed known-target oracle bound at every $np\to\infty$
with bounded $np^2$ on this classical class.
\end{remark}



\section{Probability closure for the historical decoders}
\label{sec:probability}

The preceding sections prove fitted-score comparisons on events whose
deterministic premises are consequences of raw noise, degree and row
concentration bounds. This section completes the unconditional theorems.
It is not enough that those events have probability tending to one.

\subsection{Events with an arbitrarily large fixed exponential rate}

For any fixed $H>0$, choose the raw norm threshold $L_H\sqrt d$ in
\eqref{eq:raw-norm-tail}. Since $\lambda_i\le\lambda_+d$ for a fixed
$\lambda_+$, the elementary Bernoulli Chernoff inequality also gives,
for a sufficiently large fixed $b_H$,
\[
 \PP\{\max_iD_i>b_Hd\}\le ne^{-Hd}.
\]
For example the exponent in the Poisson upper bound is
$b_H\log(b_H/\lambda_+)-b_H+\lambda_+$, which grows without bound as
$b_H$ grows. Here this is a bound on a Bernoulli transform, not a change
of sampling model.
The conditional events in \cref{lem:normalized-rows} contribute at most
$2Kn e^{-Hd}$. The constants in the deterministic lemmas may increase
with $H$, but remain independent of $n$.

For DC-SBM we additionally impose
\begin{equation}\label{eq:raw-block-event}
 \max_{a,b}|[Z^\top(A-\Omega)Z]_{ab}|\le L_Bn.
\end{equation}
Each centered block total is a sum of independent bounded upper-triangular
edge variables, with coefficients at most two and variance at most $Cnd$.
The completed diagonal contributes only $O(d)$.
Scalar Bernstein therefore implies, for sufficiently large fixed $L_B$
and all large $n$,
\[
 \PP\{\eqref{eq:raw-block-event}\text{ fails}\}
 \le 2K^2\exp\left\{-\frac{cL_B^2 n}{d+L_B}\right\}.
\]
Under $d^2\le C_1n$, the exponent is at least
$c'L_B^2d/C_1$ once $d\ge L_B$. Thus $L_B$ can be chosen to make this
probability at most $2K^2e^{-Hd}$ for any prescribed fixed $H$.
Ordinary free-profile fitting does not need this event: it integrates the
block-fluctuation exponential moment instead, as proved in
\cref{sec:sbm-proof}.

Let $\mathcal G$ be the intersection of these required events and the
minimum-degree event $D_i\ge a_0d$. All spectral, initialization and fitting
premises in the earlier sections follow on $\mathcal G$.
In particular all tested deleted graphs satisfy their own measurable
training premises there. Apart from the minimum-degree event, the failure
probability is at most
\begin{equation}\label{eq:other-bad-events}
 C_Kne^{-Hd}.
\end{equation}

\subsection{Ordinary SBM on the explicit parameter class}

For $P_{ab}\in[p/2,p]$ and $p\le1/4$, the log-odds coefficient of any
pairwise oracle contrast satisfies
\[
 |w_c^{ab}|\le W:=\log(7/3).
\]
Every tilted Bernoulli probability between the two candidates is at most
$p$, so the second derivative of the log affinity in
\eqref{eq:sbm-information} is at most $dW^2$.
A twice differentiable convex function vanishing at zero and one with
second derivative at most $M$ is at least $-Mt(1-t)/2$.
Consequently
\begin{equation}\label{eq:sbm-information-upper}
 I_{i,ab}\le d\,\frac{\log^2(7/3)}8<0.09d.
\end{equation}

For all sufficiently large $n$, every mean degree is at least $0.49d$.
Choose $a_0=1/16$. The Bernoulli lower-tail transform yields
\begin{align*}
 \PP(D_i<d/16)
 &\le \exp\left\{-d\left[
 0.49-\frac1{16}-\frac1{16}\log(16\cdot0.49)\right]\right\}\\
 &\le e^{-0.298d}.
\end{align*}
\begin{equation}\label{eq:sbm-low-degree}
 \PP(\min_iD_i<d/16)\le ne^{-0.298d}.
\end{equation}
Take $H=1$, increasing the other fixed thresholds if necessary.
Equations \eqref{eq:other-bad-events} and
\eqref{eq:sbm-low-degree} give
$\PP(\mathcal G^c)\le Cne^{-0.298d}$.
The information bound \eqref{eq:sbm-information-upper} applies to every
vertex and every competitor. Thus the sharp pairwise lower bound
\eqref{eq:oracle-pairs}, followed by averaging over all vertices, gives
\[
 R_{\mathrm{SBM},*}\ge
 \frac1n\sum_{i=1}^n\frac c{\sqrt d}e^{-0.09d}
 =\frac c{\sqrt d}e^{-0.09d}.
\]
No additional balancing assumption is needed: the uniform coefficient
bound, group-proportion bound and full-rank separation already hold at
every vertex. Consequently
\[
 \frac{\PP(\mathcal G^c)}{R_{\mathrm{SBM},*}}
 \le Cn\sqrt d\,e^{-0.208d}.
\]
The function $\sqrt d\,e^{-0.208d}$ is decreasing once
$d>1/(2\cdot0.208)$. Therefore, for $d\ge5\log n$ and all sufficiently
large $n$,
\[
 n\sqrt d\,e^{-0.208d}
 \le \sqrt{5\log n}\,n^{1-5\cdot0.208}
 =\sqrt{5\log n}\,n^{-0.04}=o(1).
\]
The earlier coefficient $20$ was a sufficient numerical choice obtained
by retaining only one vertex in the oracle lower bound and then bounding
$\sqrt d$ by $\sqrt n$. Both losses are unnecessary. The improved
coefficient $5$ changes neither the encoder nor the likelihood decoder;
it does not remove the global minimum-degree event from this proof.
The good-event comparison in \cref{sec:sbm-proof} and
\eqref{eq:bad-payment} complete \cref{thm:sbm}.

\subsection{The automatic margin for a general DC profile oracle}

Let $J_{i,a\to b}$ be the exact Bernoulli directional profile information
defined in \cref{sec:dcsbm-proof}. For the self-excluded neighbor masses
$M_c=\sum_{j\ne i:z_j=c}\theta_j$, set
\[
 T_a=\sum_cM_cC_{ac},\qquad
 q_{ac}=\frac{M_cC_{ac}}{T_a},\qquad
 h_{ab}^*=\min_{0\le s\le1}\sum_cq_{ac}^{1-s}q_{bc}^s.
\]
Uniform positive group proportions and entries of $C$ give
$q_{ac}\ge q_->0$ for all large $n$. Hence
\begin{equation}\label{eq:profile-overlap}
 h_{ab}^*\ge\sum_c\min(q_{ac},q_{bc})\ge Kq_->0.
\end{equation}
The exact log-transform expansion, uniformly in $s\in[0,1]$, is
\begin{align*}
 \sum_{j\ne i}\log\{1+\theta_i\theta_jpC_{az_j}
                    (e^{-sw_{z_j}}-1)\}
 &=\theta_ipT_a
       \left[\sum_cq_{ac}^{1-s}q_{bc}^s-1\right]+O(np^2).
\end{align*}
Minimizing therefore gives
\begin{equation}\label{eq:profile-leading-information}
 J_{i,a\to b}
 =\lambda_i(1-h_{ab}^*)+O(np^2),
 \qquad \lambda_i=\theta_ipT_a.
\end{equation}
This calculation uses the exact Bernoulli transform.
Under $np^2\le C_1$, \eqref{eq:profile-overlap} and
$\lambda_i\ge\lambda_-d$ imply, for a fixed $\eta>0$ and all large $n$,
\begin{equation}\label{eq:profile-degree-margin}
 J_{\min}:=\min_i\min_{b\ne z_i}J_{i,z_i\to b}
 \le\lambda_{\min}-\eta d.
\end{equation}
To see it, choose a vertex with minimum mean degree and any competitor.
Equal candidate exposures are not required for this profile margin.

Choose $a_0>0$ sufficiently small that
\[
 a_0\{1+\log(\lambda_+/a_0)\}\le\eta/2,\qquad a_0<\lambda_-.
\]
For each vertex, the Bernoulli lower-tail bound is
\[
 \PP(D_i<a_0d)
 \le\exp\{-\lambda_i+a_0d+
                   a_0d\log(\lambda_i/(a_0d))\}.
\]
The exponent is monotone in $\lambda_i$ above $a_0d$.
Together with \eqref{eq:profile-degree-margin}, it follows that
\begin{equation}\label{eq:dc-low-degree}
 \PP(\min_iD_i<a_0d)
 \le n\exp\{-J_{\min}-\eta d/2\}.
\end{equation}
The full profile reference has the lower bound
\[
 R_{\mathrm{prof},*}\ge
 \frac c{n\sqrt d}e^{-J_{\min}}.
\]
Thus the ratio of \eqref{eq:dc-low-degree} to that reference is at most
\[
 Cn^2\sqrt d\,e^{-\eta d/2}=o(1)
\]
when $d\ge C_0\log n$ with, for example, $C_0>6/\eta$ and sufficiently
large $n$. We used only $d\le n$ in bounding the polynomial factor.
The constants can be enlarged without changing their independence of $n$.

Also $J_{\min}\le C_Jd$ for a fixed $C_J$ by its bounded-weight transform.
Choosing $H>C_J+1$ in \eqref{eq:other-bad-events}, and enlarging $C_0$ if
needed, pays those remaining failures at the same reference scale.
Hence $\PP(\mathcal G^c)=O(R_{\mathrm{prof},*})$.
The good-event theorem and invariant-basin result in
\cref{sec:dcsbm-proof} now prove \cref{thm:dc}, including its finite
hard-refit version.

Under the equal-exposure condition, the profile and known-target-degree
oracle risks are comparable by the bounded-contrast proof in
\cref{sec:dcsbm-proof}. This proves the DC assertion of \cref{cor:known}.
For the native Euclidean ordinary-SBM decoder, the same sparse transform
expansion gives
$I_{i,ab}=\lambda_i(1-h_{ab}^*)+O(np^2+p)$ under equal exposures.
Thus its information--degree margin is also automatic.
The general-positive-profile version of \cref{sec:sbm-proof}, with its
block fluctuation integrated instead of imposing \eqref{eq:raw-block-event},
therefore proves the remaining assertion of \cref{cor:known}.

\subsection{An explicit extension for ordinary SBM with unequal exposures}

The entry-ratio-two restriction in \cref{thm:sbm} is a simple sufficient
way to ensure the failure-event comparison. The deterministic and
conditional likelihood proof covers a larger parameter range.

\begin{proposition}[A deterministic information--degree margin suffices]
\label{prop:sbm-margin}
Suppose ordinary SBM has fixed $K$, positive comparable group sizes,
$c_-p\le P_{ab}\le c_+p$, $\max P_{ab}\le1-\epsilon_0$,
and $\sigma_{\min}(P)\ge\kappa p$, for fixed positive constants.
If for some fixed $\eta>0$ and all sufficiently large $n$,
\[
 \lambda_{\min}-I_{\min}\ge\eta d,\qquad d=np\ge C_0\log n,
\]
then the native Euclidean signed likelihood decoder, with sufficiently
small structural floor coefficient, satisfies
$\E\ell(\widehat z_E,z)\le C R_{\mathrm{SBM},*}$
when $C_0$ is sufficiently large. No condition $np^2=O(1)$ is needed
for this ordinary-SBM proposition.
\end{proposition}

\begin{proof}
All true Bernoulli likelihood coefficients and tilted means are exactly
block constant, so the Euclidean projection, fitted-score and sharp-tail
arguments of \cref{sec:sbm-proof} apply without a sparse Bernoulli
approximation. Generated profiles remain comparable to $p$ and uniformly
away from one when the floor is chosen for the structural class.
The displayed margin, the same small degree cutoff used above and an
arbitrarily strong choice of the other exponential bounds imply
$\PP(\mathcal G^c)=O(R_{\mathrm{SBM},*})$.
Here the polynomial comparison uses $d\le Cn$, which follows from
$c_-p\le P_{ab}\le1-\epsilon_0$; it does not require $p\le1$.
\end{proof}

This proposition is a condition on the deterministic model parameters,
not an assumed success event or a hypothesis that the desired risk ratio
already holds. It is nevertheless an additional restriction. For model
families violating it, a refined analysis of exceptional low-degree
graphs or a different normalized encoder is needed; the crude loss bound
one on all such graphs does not prove a constant ratio.

\section{Why the known-target oracle cannot be the general DC-SBM benchmark}
\label{sec:impossibility}

The profile reference and an oracle supplied with the target's true degree
parameter are different experiments. The distinction cannot be removed by
a more accurate spectral reconstruction: the following obstruction holds
even when the classifier observes the entire adjacency matrix. The example
has bounded positive degree parameters, exact block normalization, balanced
groups, a positive definite block matrix, and logarithmic expected degree.
Every probability calculation below is for the Bernoulli graph itself.

\paragraph{The local known-target reference.}
For a fixed DC-SBM parameter point \(\tau=(z^\tau,\theta^\tau)\), write
\(\PP_\tau\) for the graph law. At a target vertex \(v\), give a local oracle
the values of \(p,C,\theta_v^\tau\) and all neighbor parameters
\((z_j^\tau,\theta_j^\tau)\), \(j\ne v\). Its candidate row laws are
\begin{equation}
 q_{vj}^{(a,\tau)}
   =p\theta_v^\tau\theta_j^\tau C_{a z_j^\tau},\qquad
 L_v^{(a,\tau)}(A_{v,-v})
   =\prod_{j\ne v}
     \bigl(q_{vj}^{(a,\tau)}\bigr)^{A_{vj}}
     \bigl(1-q_{vj}^{(a,\tau)}\bigr)^{1-A_{vj}}.
 \label{eq:imposs-local-laws}
\end{equation}
Let \(\delta_v^{\mathrm{known},\tau}\) maximize these two likelihoods with
equal candidate priors and a fixed tie rule, and define
\begin{equation}
 R_{\mathrm{known}}^\tau(n)
   =\frac1n\sum_{v=1}^n
     \PP_\tau\{\delta_v^{\mathrm{known},\tau}\ne z_v^\tau\}.
 \label{eq:imposs-known-risk}
\end{equation}
The target parameter is held at its supplied true value under both candidate
labels. The local comparison does not impose a global exact-group-size or
block-normalization constraint after revealing the other labels; an exact
group-size constraint would itself determine a missing label. Thus
\eqref{eq:imposs-known-risk} is the nondegenerate local likelihood benchmark,
not a globally constrained oracle experiment.

\begin{theorem}[An exact Bernoulli obstruction for the known-target oracle]
\label{thm:imposs}
Let \(n=2m\) tend to infinity through even integers, and set
\begin{equation}
 C=\begin{pmatrix}4&2\\2&11/10\end{pmatrix},\qquad
 p=\frac{M\log n}{n},\qquad M=100.
 \label{eq:imposs-model}
\end{equation}
There are two explicit parameter sequences, denoted \(+\) and \(-\), for
the independent-edge Bernoulli DC-SBM with probabilities
\(p\theta_u\theta_vC_{z_uz_v}\), such that both sequences satisfy
\begin{equation}
 n_1=n_2=m,\qquad
 \sum_{v:z_v=a}\theta_v=m\quad(a=1,2),\qquad
 \frac12\le\theta_v\le\frac32.
 \label{eq:imposs-normalization}
\end{equation}
For every possibly randomized clustering algorithm based on the full graph
and information common to the two parameter points, put
\begin{equation}
 R_{\mathrm{alg}}^\tau(n)
   =\E_\tau\left[
       \frac1n\min_{\sigma\in\mathfrak S_2}
       \sum_{v=1}^n\ind{\sigma(\widehat z_v)\ne z_v^\tau}
     \right],\qquad \tau\in\{+,-\}.
 \label{eq:imposs-alg-risk}
\end{equation}
Define
\begin{equation}
 c=\frac{(2-\sqrt2)^2}{8},\qquad
 \gamma=\frac{(\sqrt{11/10}-1)^2}{2}.
 \label{eq:imposs-constants}
\end{equation}
Then
\begin{equation}
 \max_{\tau\in\{+,-\}}
 \frac{R_{\mathrm{alg}}^\tau(n)}{R_{\mathrm{known}}^\tau(n)}
 \ge \frac{1+o(1)}4\,
       n^{M(c-2\gamma)-1}.
 \label{eq:imposs-ratio}
\end{equation}
The exponent is strictly positive; at \(M=100\),
\begin{equation}
 M(c-2\gamma)-1
   =3.05109151537556\ldots>0.
 \label{eq:imposs-positive-gap}
\end{equation}
In particular, no algorithm with unknown vertex degree parameters can have
a bounded risk ratio to \eqref{eq:imposs-known-risk} along every admissible
parameter sequence in this class. This includes every algorithm based only
on retained spectral information and observed degrees.
\end{theorem}

\begin{proof}
\textbf{Construction and exact normalization.}
Distinguish vertices \(i=1\) and \(k=2\). Partition the remaining vertices
into two fixed anchor sets \(V_1,V_2\), each of size \(m-1\). Every anchor
in \(V_a\) has label \(a\) under both parameter points. Give anchors in
\(V_1\) the common degree parameter
\[
 t_m=1+\frac1{2(m-1)},
\]
and anchors in \(V_2\) degree parameter one. At the two parameter points set
\begin{equation}
 \begin{aligned}
 +:\quad &(z_i,\theta_i)=(1,1/2),
          & (z_k,\theta_k)&=(2,1),\\
 -:\quad &(z_i,\theta_i)=(2,1),
          & (z_k,\theta_k)&=(1,1/2).
 \end{aligned}
 \label{eq:imposs-two-points}
\end{equation}
Each group has exactly \(m\) vertices. Its total degree parameter is exactly
\(m\), since
\[
 (m-1)t_m+\frac12=m,
 \qquad (m-1)+1=m.
\]
The stated degree bounds hold for \(m\ge2\). The matrix \(C\) has positive
entries, positive trace, and determinant \(2/5>0\), so it is positive
definite with a fixed positive smallest eigenvalue. All edge probabilities
belong to \((0,1)\) for sufficiently large \(n\).

\medskip
\textbf{Only \(n-2\) Bernoulli coordinates distinguish the graph laws.}
All anchor--anchor probabilities agree between the two parameter points.
The edge \(i\)--\(k\) has probability \(p\) under either point. For an edge
from either distinguished vertex to an anchor of \(V_1\), the two candidate
intensities agree exactly:
\[
 \frac12\cdot4\cdot t_m=1\cdot2\cdot t_m=2t_m.
\]
Thus the only distinct coordinates are the \(2(m-1)=n-2\) edges from
\(\{i,k\}\) to \(V_2\). Their two probabilities are \(p\) and
\((11/10)p\), exchanged between \(+\) and \(-\).

For probability mass functions \(P,Q\), let
\(\mathcal A(P,Q)=\sum_x\sqrt{P(x)Q(x)}\) be their affinity. Independence
gives the exact graph affinity
\begin{equation}
 \mathcal A_n
 =\mathcal A(\PP_+,\PP_-)
 =\left[
     \sqrt{11/10}\,p+
       \sqrt{(1-p)(1-(11/10)p)}
   \right]^{n-2}.
 \label{eq:imposs-affinity}
\end{equation}
The factor in brackets has the elementary expansion
\[
 \sqrt{11/10}\,p+
       \sqrt{(1-p)(1-(11/10)p)}
 =1-\gamma p+O(p^2).
\]
Since \(np^2=M^2(\log n)^2/n=o(1)\),
\begin{equation}
 \mathcal A_n^2
   =\exp\{-2M\gamma\log n+o(1)\}
   =(1+o(1))n^{-2M\gamma}.
 \label{eq:imposs-affinity-rate}
\end{equation}
The elementary Cauchy--Schwarz bound
\(\TV(P,Q)\le\sqrt{1-\mathcal A(P,Q)^2}\) implies that every test \(\phi\)
of the two graph laws satisfies
\begin{align}
 \frac{\PP_+(\phi=-)+\PP_-(\phi=+)}2
 &\ge \frac{1-\TV(\PP_+,\PP_-)}2 \\
 &\ge \frac{\mathcal A_n^2}{4}
  =\frac{1+o(1)}4 n^{-2M\gamma},
 \label{eq:imposs-test-lower}
\end{align}
where the second inequality uses
\(1-\sqrt{1-y}\ge y/2\). Independent randomization by the test or the
clustering algorithm does not alter this bound: adjoining a common random
seed leaves the affinity unchanged.

\medskip
\textbf{From testing to graph-average permutation loss.}
Orient the two output labels of an arbitrary estimated partition by the
permutation minimizing disagreement with the fixed anchor labels, using
a fixed tie rule. This is one measurable rule under both parameter points,
since the anchors and their labels are the same. Define a test \(\phi\) to
declare \(+\) if the oriented label of vertex \(i\) is 1, and \(-\)
otherwise.

If the estimated partition has zero permutation-invariant error, the
anchor orientation recovers its correct permutation uniquely: both anchor
sets are nonempty. Therefore an erroneous test outcome implies at least
one error in the permutation-invariant count. Pathwise, under either point,
\begin{equation}
 \frac1n\min_{\sigma\in\mathfrak S_2}
       \sum_v\ind{\sigma(\widehat z_v)\ne z_v^\tau}
 \ge\frac1n\ind{\phi\ne\tau}.
 \label{eq:imposs-orientation}
\end{equation}
Taking expectations and using \eqref{eq:imposs-test-lower} gives
\begin{equation}
 \frac{R_{\mathrm{alg}}^+(n)+R_{\mathrm{alg}}^-(n)}2
 \ge\frac{1+o(1)}4 n^{-2M\gamma-1}.
 \label{eq:imposs-graph-risk-lower}
\end{equation}
The anchors are needed only to turn a partition into a test in this proof.
Indeed, their labels and degree parameters, the block matrix, the exact
group sizes, and the block-normalization constraints could all be revealed
to the clustering algorithm: they are common to the two parameter points,
and \eqref{eq:imposs-test-lower} would still apply.

\medskip
\textbf{An upper bound on every known-target local error.}
For any target \(v\), at least \(m-2\) anchors in \(V_1\) are different
from \(v\). For each such anchor \(j\), its two local candidate edge
probabilities in \eqref{eq:imposs-local-laws} are
\[
 u=4p\theta_v\theta_j,
 \qquad v'=2p\theta_v\theta_j.
\]
They satisfy
\[
 (\sqrt u-\sqrt{v'})^2
 =p\theta_v\theta_j(2-\sqrt2)^2
 \ge\frac p2(2-\sqrt2)^2.
\]
For any Bernoulli probabilities \(u,v'\), the arithmetic--geometric mean
inequality for their nonedge probabilities yields
\begin{equation}
 \sqrt{uv'}+\sqrt{(1-u)(1-v')}
 \le 1-\frac{(\sqrt u-\sqrt{v'})^2}{2}
 \le \exp\left\{-\frac{(\sqrt u-\sqrt{v'})^2}{2}\right\}.
 \label{eq:imposs-one-edge-upper}
\end{equation}
Under either local true label, Markov's inequality applied to the square
root of the competing likelihood ratio bounds the LRT's directional error
by the product affinity of all its row coordinates. Discarding coordinates
outside these \(m-2\) anchors can only increase that upper bound. Hence,
uniformly over targets and both parameter points,
\begin{align}
 \PP_\tau\{\delta_v^{\mathrm{known},\tau}\ne z_v^\tau\}
 &\le\exp\left\{-\frac{(m-2)p}{4}(2-\sqrt2)^2\right\} \\
 &= (1+o(1))n^{-Mc}.
 \label{eq:imposs-oracle-upper}
\end{align}
In particular,
\begin{equation}
 R_{\mathrm{known}}^+(n)+R_{\mathrm{known}}^-(n)
 \le 2(1+o(1))n^{-Mc}.
 \label{eq:imposs-oracle-sum-upper}
\end{equation}

Finally, the maximum of two nonnegative risk ratios is at least their
denominator-weighted average. Combining
\eqref{eq:imposs-graph-risk-lower} and
\eqref{eq:imposs-oracle-sum-upper} therefore proves
\eqref{eq:imposs-ratio}. The constants in
\eqref{eq:imposs-constants} are
\[
 c=0.04289321881345246\ldots,
 \qquad \gamma=0.00119115182984846\ldots,
\]
which give the exponent in \eqref{eq:imposs-positive-gap}. Its positivity
does not depend on numerical evaluation: \(c>1/25\) and
\(2\gamma<1/400\), so \(100(c-2\gamma)-1>11/4>0\).
\end{proof}

\begin{remark}[Scope of the obstruction]
\label{rem:imposs-scope}
The same proof works for every fixed
\(M>(c-2\gamma)^{-1}\), so its logarithmic-degree multiplier can exceed
any prescribed fixed lower threshold. The example also satisfies
\(np^2\to0\). Thus neither a larger logarithmic expected degree nor the
sparse Bernoulli approximation removes the nuisance-information gap.

The result concerns a uniform comparison to the known-target reference,
not the full-edge profile benchmark used by the positive DC-SBM theorem.
It also does not assert that every fixed DC-SBM has an unbounded ratio:
it proves that at least one of the two explicit parameter sequences has
an unbounded limsup ratio for any algorithm. Additional restrictions can
make the stronger benchmark attainable. For example, the balanced block
strengths in this counterexample are
\(C(1/2,1/2)^T=(3,31/20)^T\), which are unequal; the example lies outside
the equal-exposure corollary.

The counterexample is also outside the assortative parameter class of
Gao, Ma, Zhang and Zhou~\cite{gmzz2018}: its smallest diagonal entry is
$1.1$, whereas its off-diagonal entry is $2$. Their class requires every
within-group strength to exceed every between-group strength. Therefore
this obstruction is not a negative result for their classical balanced
planted-partition benchmark. That benchmark is explicitly covered by
\cref{prop:classical-dc-oracle-bridge} on the stated bounded-parameter
class. Their classwise minimax theorem and a pointwise comparison to a
known-target oracle on every arbitrary block matrix are different claims.

Finally, a lower bound for an experiment retaining only block counts would
not by itself establish a lower bound for arbitrary data-dependent real
spectral coordinates. The present theorem needs no such compression claim:
it is a lower bound for the full adjacency experiment with unknown target
degree parameters.
\end{remark}

\section{Additional arguments, implementation scope, and retained records}
\label{sec:scope}
\subsection{Why a whole-graph zero-degree fallback fails}
\begin{proposition}[A whole-graph zero-degree fallback is invalid in the
sparse regime]\label{restore:fallback}
Consider a balanced two-group SBM with within probability $p$, between
probability $p/2$, and $d=np=\log\log n$. The probability of at least
one isolated vertex tends to one. For every deterministic fallback
partition fixed in advance, there is a balanced true partition for which
an algorithm returning that fallback whenever an isolated vertex exists
has expected permutation-invariant misclassification bounded below by
$1/4+o(1)$. Its local oracle risk tends to zero.
\end{proposition}

\begin{proof}
For every vertex the isolation probability is
$(1-p)^{n/2-1}(1-p/2)^{n/2}
=\exp\{-3d/4+o(1)\}$. If $N_0$ counts isolated vertices, then
$\E N_0=n\exp\{-3d/4+o(1)\}\to\infty$. For distinct vertices,
their two isolation events share only the edge between them, giving
\[
 \frac{\operatorname{Var}(N_0)}{(\E N_0)^2}
 \le \frac1{\E N_0}+\frac p{1-p}\longrightarrow0.
\]
Choose a balanced true partition whose overlap with the fixed fallback
has error at least $1/4$ after either label permutation. Such a partition
exists by splitting each fallback class approximately evenly. The loss
lower bound follows. Standard incident-row Chernoff comparison gives
oracle risk at most $Ce^{-cd}\to0$ on this separated model.
\end{proof}

Thus the logarithmic restriction cannot be removed merely by deleting
text from the unregularized theorem. The valid repair is to restore a
positive floor and local risk accounting. With the regularized encoder,
a zero-degree vertex has a zero reconstructed row and causes no
whole-graph fallback.

\subsection{The existing degree-plus-spectrum action remains valid}

Write $\mathbf d=A\one$ for the raw degree vector. For any strictly
positive diagonal matrix $W$, retain
\[
 WAWU=U\Lambda,\qquad U^\top U=I,\qquad
 X=[\one,WU],\qquad Y=[\mathbf d,W^{-1}U\Lambda].
\]
Then the exact identities
\begin{equation}\label{eq:sparse-old-action}
 AX=Y,\qquad YX^\dagger=A\Pi_X,\qquad A\Pi_X\one=\mathbf d
\end{equation}
hold for every graph, including graphs with zero degrees. Here $\Pi_X$
is the Euclidean orthogonal projector onto the column space of $X$;
no full-column-rank assumption is necessary. The earlier choice
$W_{ii}=\min(1,2\bar D/D_i)$ for $D_i>0$, with $W_{ii}=1$ for
$D_i=0$ and $\bar D=n^{-1}\sum_iD_i$, supplies a positive reweighting
when $\bar D>0$. The all-zero graph can be assigned a fixed output.
The block-count error is exactly $-A(I-\Pi_X)Z$.

These facts need no logarithmic expected-degree assumption. They are
computational identities, however: their validity does not give the
sharp tail of the remaining random block-count error. In DC-SBM a
normalization must also preserve the effective likelihood coefficients,
as opposed to preserving only a population mean.


\subsection{Algorithm distinctions and scope}
The current ordinary article uses the regularized block Bernoulli LRT,
with fixed degree parameters equal to one and no target profiling.
The current DC result in this companion uses clamped, group-normalized
nuisance estimates and profiles the target separately for each candidate.
Both retain raw degrees and the informative normalized eigensystem,
apply exact degree correction and simple hollowing, and control finite
hard refits through a uniform invariant basin.

The older ordinary free-profile growing/replacement/soft-EM schedule
in \cref{sec:sbm-proof} is a different finite algorithm. Its separate
proof also covers probabilities outside bounded $np^2$. The old common-box
proof uses a simultaneous entry bound and must not be attached without
change to the present sparse encoder. In particular, exact degree
preservation together with $|R_{ij}|\le Cp$ would force all raw degrees
to be at most $Cd$, which is not a valid sparse uniform premise.
The logarithmic restrictions written in the historical part describe
those earlier unregularized proofs, not the current theorem.

The earlier DC reconstruction $AH_D$ is different from the Euclidean
$A\Pi_X$, even though both can be evaluated from the same retained
input. Earlier experimental curves do not automatically evaluate either
the new empirical floor or its clamped nuisance fit. The present claims
are mathematical risk bounds, not conclusions inferred from simulations.
Numerical eigensolver tolerances and optimization tolerances require their
own control of induced score errors; exact arithmetic defines these
algorithms. The largest absolute eigenvalues are retained so that a
full-rank indefinite signal is included.

The constants may depend on the fixed signal gap, group proportions,
bounded degree-parameter range, and $C_1$. A bounded risk ratio is not
a claim that the ratio tends to one. General DC-SBM with unbounded
$np^2$, a vanishing signal gap, or degree parameters approaching zero
is outside the present theorem. The general known-target-degree
benchmark is obstructed by \cref{thm:imposs}; equal exposures retain
the positive result.

\subsection{Sparse proof history and review records}
The 9 October degree-plus-spectrum draft contained the sparse truncated
row-energy mechanism at $d\to\infty$, $d=o(\log n)$. Its score errors
were $O(\sqrt d)$ or $o(d)$ and therefore established a leading exponent,
not the present constant risk ratio. The current argument keeps its
same-threshold mechanism, upgrades the local full-to-deleted score
comparison and fitted contrast errors to $O(1)$, and preserves the sharp
$d^{-1/2}$ prefactor. The earlier source remains in the repository at
\href{https://github.com/zhixin0825/spectral-lrt/blob/de60b64c7d7f127ae6dc1d010c840af707a080da/paper/previous_draft/sbm_spectral_lrt_20261009.tex}{commit \texttt{de60b64}}.

The source package includes the \texttt{audits/} directory and its index.
It retains the sparse-history audit, the independent local-floor LOO
review, the same-threshold tail and conditioning review, the high-scale
probability review, the ordinary algorithm-scope review, the final
combined review, and earlier DC and regularization audit snapshots.
These are internal working checks, not external peer review. Statements
of then-open gaps in older snapshots are superseded by the current
proofs. The full pre-split revision is preserved separately in the
repository; it is an archive rather than the reading order of these notes.

An inventory in \texttt{README.md} records where each old source module
now belongs. The original implementation/audit narrative is retained as
\texttt{archive/scope\_audit.tex}; the obsolete end-of-revision status
narrative is retained as \texttt{archive/sparse\_interfaces.tex}. These
source records are not assertions of a remaining logarithmic lower bound.

\bibliographystyle{plain}
\bibliography{references}
\end{document}
